\documentclass[11pt,a4paper]{article}
% preamble.tex - shared house preamble for /explain-projects documents.
% Every package here is VERIFIED present in this machine's TeX Live install.
% Do not add packages without running: kpsewhich <pkg>.sty
% Notably ABSENT (do not \usepackage them, the build will fail):
%   algorithm2e, algorithm, algorithmic, algpseudocode  (texlive-science is not installed)
%   -> use the \begin{pseudocode} tcolorbox environment defined below instead.
%   lmodern, charter, mathptmx, times, mathpazo, courier, newtx*, tgtermes
%   -> their .sty files may resolve but the FONT FILES are absent, so the build
%      dies with "I can't find file `bchr8t'" or similar. Verified 2026-07-16 by
%      test-compiling each one. Only Computer Modern (the default) and helvet
%      actually render on this box. Do not "upgrade" the font stack without
%      test-compiling first: kpsewhich lying about a font is the exact trap here.
%
% To get a nicer font stack + algorithm2e, install:
%   sudo apt install texlive-fonts-recommended texlive-science
% Until then, this preamble deliberately uses the default font.
%
% Usage from a document:
%   \documentclass[11pt,a4paper]{article}
%   % preamble.tex - shared house preamble for /explain-projects documents.
% Every package here is VERIFIED present in this machine's TeX Live install.
% Do not add packages without running: kpsewhich <pkg>.sty
% Notably ABSENT (do not \usepackage them, the build will fail):
%   algorithm2e, algorithm, algorithmic, algpseudocode  (texlive-science is not installed)
%   -> use the \begin{pseudocode} tcolorbox environment defined below instead.
%   lmodern, charter, mathptmx, times, mathpazo, courier, newtx*, tgtermes
%   -> their .sty files may resolve but the FONT FILES are absent, so the build
%      dies with "I can't find file `bchr8t'" or similar. Verified 2026-07-16 by
%      test-compiling each one. Only Computer Modern (the default) and helvet
%      actually render on this box. Do not "upgrade" the font stack without
%      test-compiling first: kpsewhich lying about a font is the exact trap here.
%
% To get a nicer font stack + algorithm2e, install:
%   sudo apt install texlive-fonts-recommended texlive-science
% Until then, this preamble deliberately uses the default font.
%
% Usage from a document:
%   \documentclass[11pt,a4paper]{article}
%   % preamble.tex - shared house preamble for /explain-projects documents.
% Every package here is VERIFIED present in this machine's TeX Live install.
% Do not add packages without running: kpsewhich <pkg>.sty
% Notably ABSENT (do not \usepackage them, the build will fail):
%   algorithm2e, algorithm, algorithmic, algpseudocode  (texlive-science is not installed)
%   -> use the \begin{pseudocode} tcolorbox environment defined below instead.
%   lmodern, charter, mathptmx, times, mathpazo, courier, newtx*, tgtermes
%   -> their .sty files may resolve but the FONT FILES are absent, so the build
%      dies with "I can't find file `bchr8t'" or similar. Verified 2026-07-16 by
%      test-compiling each one. Only Computer Modern (the default) and helvet
%      actually render on this box. Do not "upgrade" the font stack without
%      test-compiling first: kpsewhich lying about a font is the exact trap here.
%
% To get a nicer font stack + algorithm2e, install:
%   sudo apt install texlive-fonts-recommended texlive-science
% Until then, this preamble deliberately uses the default font.
%
% Usage from a document:
%   \documentclass[11pt,a4paper]{article}
%   \input{preamble}          % this file is copied next to the .tex

\usepackage[T1]{fontenc}
\usepackage[utf8]{inputenc}
% Font: Computer Modern (default). See the note above before changing this.
\usepackage[a4paper,margin=2.4cm,headheight=14pt]{geometry}
% microtype WITHOUT font expansion. This is load-bearing, do not "restore" it:
% cm-super is not installed, so T1 Computer Modern resolves to BITMAP fonts, and
% pdfTeX's auto-expansion needs scalable ones. With expansion on, any document
% long enough to reach a second page dies with
%   "pdfTeX error (font expansion): auto expansion is only possible with scalable fonts"
% Protrusion (the bigger visual win) still works and stays on.
% `sudo apt install cm-super` would lift this; until then, leave expansion off.
\usepackage[protrusion=true,expansion=false]{microtype}
\usepackage{graphicx}
\usepackage{xcolor}
\usepackage{booktabs}
\usepackage{enumitem}
\usepackage{fancyhdr}
\usepackage{titlesec}
\usepackage{tcolorbox}
\usepackage{listings}
\usepackage{float}
\usepackage{caption}
\usepackage{tikz}
\usepackage{forest}
\usepackage{pgfplots}
\pgfplotsset{compat=1.18}
\usetikzlibrary{arrows.meta,positioning,shapes.geometric,fit,backgrounds,calc,decorations.pathreplacing}
\tcbuselibrary{skins,breakable}
\usepackage[hidelinks,bookmarks=true,bookmarksnumbered=true]{hyperref}

% ===== palette ===========================================================
\definecolor{inkblue}{HTML}{1F3A5F}
\definecolor{accent}{HTML}{2E7DAF}
\definecolor{softgrey}{HTML}{F4F6F8}
\definecolor{linegrey}{HTML}{C8D0D8}
\definecolor{codebg}{HTML}{F7F7F4}
\definecolor{warnred}{HTML}{A33B3B}
\definecolor{okgreen}{HTML}{2E7D53}

% ===== headings ==========================================================
\titleformat{\section}{\Large\bfseries\color{inkblue}}{\thesection}{0.7em}{}
\titleformat{\subsection}{\large\bfseries\color{inkblue}}{\thesubsection}{0.6em}{}
\titleformat{\subsubsection}{\normalsize\bfseries\color{accent}}{\thesubsubsection}{0.5em}{}

\pagestyle{fancy}
\fancyhf{}
\fancyhead[L]{\small\color{inkblue}\nouppercase{\leftmark}}
\fancyhead[R]{\small\color{inkblue}\thepage}
\renewcommand{\headrulewidth}{0.4pt}
\renewcommand{\headrule}{\hbox to\headwidth{\color{linegrey}\leaders\hrule height \headrulewidth\hfill}}

% ===== code ==============================================================
\lstdefinestyle{house}{
  backgroundcolor=\color{codebg},
  basicstyle=\ttfamily\footnotesize,
  keywordstyle=\color{inkblue}\bfseries,
  commentstyle=\color{gray}\itshape,
  stringstyle=\color{okgreen},
  numbers=left, numberstyle=\tiny\color{gray}, numbersep=8pt,
  frame=single, rulecolor=\color{linegrey},
  breaklines=true, breakatwhitespace=false,
  showstringspaces=false, tabsize=2,
  captionpos=b, xleftmargin=14pt, framexleftmargin=14pt,
  columns=fullflexible, keepspaces=true,
  upquote=true  % render ' and " as straight quotes, not curly
}
\lstset{style=house}

% ===== callout boxes =====================================================
% \begin{keyidea}{Title} ... \end{keyidea}   - a concept a beginner must grasp
\newtcolorbox{keyidea}[1]{
  enhanced, breakable, colback=softgrey, colframe=accent,
  boxrule=0.4pt, arc=2pt, left=8pt, right=8pt, top=6pt, bottom=6pt,
  title={\bfseries #1}, coltitle=white, colbacktitle=accent, fonttitle=\small
}

% \begin{jargon}{Term} definition \end{jargon}  - define a term at first use
\newtcolorbox{jargon}[1]{
  enhanced, breakable, colback=white, colframe=linegrey,
  boxrule=0.4pt, arc=2pt, left=8pt, right=8pt, top=6pt, bottom=6pt,
  title={\bfseries Jargon: #1}, coltitle=inkblue, colbacktitle=softgrey, fonttitle=\small
}

% \begin{honest}{Title} ... \end{honest}  - brutally honest finding (rule 18)
\newtcolorbox{honest}[1]{
  enhanced, breakable, colback=white, colframe=warnred,
  boxrule=0.6pt, arc=2pt, left=8pt, right=8pt, top=6pt, bottom=6pt,
  title={\bfseries #1}, coltitle=white, colbacktitle=warnred, fonttitle=\small
}

% \begin{provenance}{path/to/file} ... \end{provenance}
% Where a file came from: hand-written, generated by a tool, scaffolded by a
% framework, copied or adapted from an upstream project, ported from an older
% module. State the EVIDENCE (a header comment, a git log entry, an upstream
% URL, a generator command) and label the confidence. If no evidence exists,
% say so in the box. Never guess a lineage into existence.
\newtcolorbox{provenance}[1]{
  enhanced, breakable, colback=white, colframe=okgreen,
  boxrule=0.4pt, arc=2pt, left=8pt, right=8pt, top=6pt, bottom=6pt,
  title={\bfseries Where this came from: #1}, coltitle=white,
  colbacktitle=okgreen, fonttitle=\small
}

% \begin{pseudocode}{Caption} ... \end{pseudocode}
% Stands in for algorithm2e, which is NOT installed. Use \Step / \Indent inside.
\newtcolorbox{pseudocode}[1]{
  enhanced, breakable, colback=codebg, colframe=inkblue,
  boxrule=0.4pt, arc=2pt, left=10pt, right=8pt, top=6pt, bottom=6pt,
  fontupper=\ttfamily\small,
  title={\bfseries\rmfamily #1}, coltitle=white, colbacktitle=inkblue, fonttitle=\small
}
\newcommand{\Step}[1]{\par\noindent #1}
\newcommand{\Indent}[1]{\par\noindent\hspace*{2em}#1}

% ===== tikz house styles =================================================
\tikzset{
  box/.style   = {rectangle, rounded corners=2pt, draw=inkblue, fill=softgrey,
                  text=inkblue, align=center, inner sep=6pt, minimum height=9mm,
                  font=\small},
  extbox/.style= {rectangle, rounded corners=2pt, draw=linegrey, fill=white,
                  text=black!70, align=center, inner sep=6pt, minimum height=9mm,
                  font=\small, dashed},
  store/.style = {cylinder, shape border rotate=90, aspect=0.25, draw=accent,
                  fill=white, text=inkblue, align=center, inner sep=5pt,
                  font=\small},
  dec/.style   = {diamond, aspect=1.6, draw=accent, fill=white, text=inkblue,
                  align=center, inner sep=2pt, font=\scriptsize},
  flow/.style  = {-{Stealth[length=2.4mm]}, draw=inkblue, thick},
  dflow/.style = {-{Stealth[length=2.4mm]}, draw=accent, thick, dashed},
  % align=center is load-bearing: without it, a \\ inside an edge label throws
  % the misleading "perhaps a missing \item" error.
  lbl/.style   = {font=\scriptsize\itshape, text=black!65, fill=white,
                  inner sep=2pt, align=center}
}

% forest default for directory trees
\forestset{
  dirtree/.style={
    for tree={
      font=\ttfamily\small, grow'=0, child anchor=west, parent anchor=south,
      anchor=west, calign=first, inner sep=2pt,
      edge path={\noexpand\path [draw, \forestoption{edge}]
        (!u.south west) +(9pt,0) |- (.child anchor) \forestoption{edge label};},
      before typesetting nodes={if n=1{insert before={[,phantom]}}{}},
      fit=band, before computing xy={l=15pt}
    }
  }
}

% ===== misc ==============================================================
\captionsetup{font=small, labelfont={bf,color=inkblue}}
\setlist[itemize]{leftmargin=1.2em, itemsep=2pt, topsep=3pt}
\setlist[enumerate]{leftmargin=1.4em, itemsep=2pt, topsep=3pt}
\newcommand{\file}[1]{\texttt{\small #1}}
\newcommand{\code}[1]{\texttt{\small #1}}
\setcounter{tocdepth}{2}
\setcounter{secnumdepth}{3}
          % this file is copied next to the .tex

\usepackage[T1]{fontenc}
\usepackage[utf8]{inputenc}
% Font: Computer Modern (default). See the note above before changing this.
\usepackage[a4paper,margin=2.4cm,headheight=14pt]{geometry}
% microtype WITHOUT font expansion. This is load-bearing, do not "restore" it:
% cm-super is not installed, so T1 Computer Modern resolves to BITMAP fonts, and
% pdfTeX's auto-expansion needs scalable ones. With expansion on, any document
% long enough to reach a second page dies with
%   "pdfTeX error (font expansion): auto expansion is only possible with scalable fonts"
% Protrusion (the bigger visual win) still works and stays on.
% `sudo apt install cm-super` would lift this; until then, leave expansion off.
\usepackage[protrusion=true,expansion=false]{microtype}
\usepackage{graphicx}
\usepackage{xcolor}
\usepackage{booktabs}
\usepackage{enumitem}
\usepackage{fancyhdr}
\usepackage{titlesec}
\usepackage{tcolorbox}
\usepackage{listings}
\usepackage{float}
\usepackage{caption}
\usepackage{tikz}
\usepackage{forest}
\usepackage{pgfplots}
\pgfplotsset{compat=1.18}
\usetikzlibrary{arrows.meta,positioning,shapes.geometric,fit,backgrounds,calc,decorations.pathreplacing}
\tcbuselibrary{skins,breakable}
\usepackage[hidelinks,bookmarks=true,bookmarksnumbered=true]{hyperref}

% ===== palette ===========================================================
\definecolor{inkblue}{HTML}{1F3A5F}
\definecolor{accent}{HTML}{2E7DAF}
\definecolor{softgrey}{HTML}{F4F6F8}
\definecolor{linegrey}{HTML}{C8D0D8}
\definecolor{codebg}{HTML}{F7F7F4}
\definecolor{warnred}{HTML}{A33B3B}
\definecolor{okgreen}{HTML}{2E7D53}

% ===== headings ==========================================================
\titleformat{\section}{\Large\bfseries\color{inkblue}}{\thesection}{0.7em}{}
\titleformat{\subsection}{\large\bfseries\color{inkblue}}{\thesubsection}{0.6em}{}
\titleformat{\subsubsection}{\normalsize\bfseries\color{accent}}{\thesubsubsection}{0.5em}{}

\pagestyle{fancy}
\fancyhf{}
\fancyhead[L]{\small\color{inkblue}\nouppercase{\leftmark}}
\fancyhead[R]{\small\color{inkblue}\thepage}
\renewcommand{\headrulewidth}{0.4pt}
\renewcommand{\headrule}{\hbox to\headwidth{\color{linegrey}\leaders\hrule height \headrulewidth\hfill}}

% ===== code ==============================================================
\lstdefinestyle{house}{
  backgroundcolor=\color{codebg},
  basicstyle=\ttfamily\footnotesize,
  keywordstyle=\color{inkblue}\bfseries,
  commentstyle=\color{gray}\itshape,
  stringstyle=\color{okgreen},
  numbers=left, numberstyle=\tiny\color{gray}, numbersep=8pt,
  frame=single, rulecolor=\color{linegrey},
  breaklines=true, breakatwhitespace=false,
  showstringspaces=false, tabsize=2,
  captionpos=b, xleftmargin=14pt, framexleftmargin=14pt,
  columns=fullflexible, keepspaces=true,
  upquote=true  % render ' and " as straight quotes, not curly
}
\lstset{style=house}

% ===== callout boxes =====================================================
% \begin{keyidea}{Title} ... \end{keyidea}   - a concept a beginner must grasp
\newtcolorbox{keyidea}[1]{
  enhanced, breakable, colback=softgrey, colframe=accent,
  boxrule=0.4pt, arc=2pt, left=8pt, right=8pt, top=6pt, bottom=6pt,
  title={\bfseries #1}, coltitle=white, colbacktitle=accent, fonttitle=\small
}

% \begin{jargon}{Term} definition \end{jargon}  - define a term at first use
\newtcolorbox{jargon}[1]{
  enhanced, breakable, colback=white, colframe=linegrey,
  boxrule=0.4pt, arc=2pt, left=8pt, right=8pt, top=6pt, bottom=6pt,
  title={\bfseries Jargon: #1}, coltitle=inkblue, colbacktitle=softgrey, fonttitle=\small
}

% \begin{honest}{Title} ... \end{honest}  - brutally honest finding (rule 18)
\newtcolorbox{honest}[1]{
  enhanced, breakable, colback=white, colframe=warnred,
  boxrule=0.6pt, arc=2pt, left=8pt, right=8pt, top=6pt, bottom=6pt,
  title={\bfseries #1}, coltitle=white, colbacktitle=warnred, fonttitle=\small
}

% \begin{provenance}{path/to/file} ... \end{provenance}
% Where a file came from: hand-written, generated by a tool, scaffolded by a
% framework, copied or adapted from an upstream project, ported from an older
% module. State the EVIDENCE (a header comment, a git log entry, an upstream
% URL, a generator command) and label the confidence. If no evidence exists,
% say so in the box. Never guess a lineage into existence.
\newtcolorbox{provenance}[1]{
  enhanced, breakable, colback=white, colframe=okgreen,
  boxrule=0.4pt, arc=2pt, left=8pt, right=8pt, top=6pt, bottom=6pt,
  title={\bfseries Where this came from: #1}, coltitle=white,
  colbacktitle=okgreen, fonttitle=\small
}

% \begin{pseudocode}{Caption} ... \end{pseudocode}
% Stands in for algorithm2e, which is NOT installed. Use \Step / \Indent inside.
\newtcolorbox{pseudocode}[1]{
  enhanced, breakable, colback=codebg, colframe=inkblue,
  boxrule=0.4pt, arc=2pt, left=10pt, right=8pt, top=6pt, bottom=6pt,
  fontupper=\ttfamily\small,
  title={\bfseries\rmfamily #1}, coltitle=white, colbacktitle=inkblue, fonttitle=\small
}
\newcommand{\Step}[1]{\par\noindent #1}
\newcommand{\Indent}[1]{\par\noindent\hspace*{2em}#1}

% ===== tikz house styles =================================================
\tikzset{
  box/.style   = {rectangle, rounded corners=2pt, draw=inkblue, fill=softgrey,
                  text=inkblue, align=center, inner sep=6pt, minimum height=9mm,
                  font=\small},
  extbox/.style= {rectangle, rounded corners=2pt, draw=linegrey, fill=white,
                  text=black!70, align=center, inner sep=6pt, minimum height=9mm,
                  font=\small, dashed},
  store/.style = {cylinder, shape border rotate=90, aspect=0.25, draw=accent,
                  fill=white, text=inkblue, align=center, inner sep=5pt,
                  font=\small},
  dec/.style   = {diamond, aspect=1.6, draw=accent, fill=white, text=inkblue,
                  align=center, inner sep=2pt, font=\scriptsize},
  flow/.style  = {-{Stealth[length=2.4mm]}, draw=inkblue, thick},
  dflow/.style = {-{Stealth[length=2.4mm]}, draw=accent, thick, dashed},
  % align=center is load-bearing: without it, a \\ inside an edge label throws
  % the misleading "perhaps a missing \item" error.
  lbl/.style   = {font=\scriptsize\itshape, text=black!65, fill=white,
                  inner sep=2pt, align=center}
}

% forest default for directory trees
\forestset{
  dirtree/.style={
    for tree={
      font=\ttfamily\small, grow'=0, child anchor=west, parent anchor=south,
      anchor=west, calign=first, inner sep=2pt,
      edge path={\noexpand\path [draw, \forestoption{edge}]
        (!u.south west) +(9pt,0) |- (.child anchor) \forestoption{edge label};},
      before typesetting nodes={if n=1{insert before={[,phantom]}}{}},
      fit=band, before computing xy={l=15pt}
    }
  }
}

% ===== misc ==============================================================
\captionsetup{font=small, labelfont={bf,color=inkblue}}
\setlist[itemize]{leftmargin=1.2em, itemsep=2pt, topsep=3pt}
\setlist[enumerate]{leftmargin=1.4em, itemsep=2pt, topsep=3pt}
\newcommand{\file}[1]{\texttt{\small #1}}
\newcommand{\code}[1]{\texttt{\small #1}}
\setcounter{tocdepth}{2}
\setcounter{secnumdepth}{3}
          % this file is copied next to the .tex

\usepackage[T1]{fontenc}
\usepackage[utf8]{inputenc}
% Font: Computer Modern (default). See the note above before changing this.
\usepackage[a4paper,margin=2.4cm,headheight=14pt]{geometry}
% microtype WITHOUT font expansion. This is load-bearing, do not "restore" it:
% cm-super is not installed, so T1 Computer Modern resolves to BITMAP fonts, and
% pdfTeX's auto-expansion needs scalable ones. With expansion on, any document
% long enough to reach a second page dies with
%   "pdfTeX error (font expansion): auto expansion is only possible with scalable fonts"
% Protrusion (the bigger visual win) still works and stays on.
% `sudo apt install cm-super` would lift this; until then, leave expansion off.
\usepackage[protrusion=true,expansion=false]{microtype}
\usepackage{graphicx}
\usepackage{xcolor}
\usepackage{booktabs}
\usepackage{enumitem}
\usepackage{fancyhdr}
\usepackage{titlesec}
\usepackage{tcolorbox}
\usepackage{listings}
\usepackage{float}
\usepackage{caption}
\usepackage{tikz}
\usepackage{forest}
\usepackage{pgfplots}
\pgfplotsset{compat=1.18}
\usetikzlibrary{arrows.meta,positioning,shapes.geometric,fit,backgrounds,calc,decorations.pathreplacing}
\tcbuselibrary{skins,breakable}
\usepackage[hidelinks,bookmarks=true,bookmarksnumbered=true]{hyperref}

% ===== palette ===========================================================
\definecolor{inkblue}{HTML}{1F3A5F}
\definecolor{accent}{HTML}{2E7DAF}
\definecolor{softgrey}{HTML}{F4F6F8}
\definecolor{linegrey}{HTML}{C8D0D8}
\definecolor{codebg}{HTML}{F7F7F4}
\definecolor{warnred}{HTML}{A33B3B}
\definecolor{okgreen}{HTML}{2E7D53}

% ===== headings ==========================================================
\titleformat{\section}{\Large\bfseries\color{inkblue}}{\thesection}{0.7em}{}
\titleformat{\subsection}{\large\bfseries\color{inkblue}}{\thesubsection}{0.6em}{}
\titleformat{\subsubsection}{\normalsize\bfseries\color{accent}}{\thesubsubsection}{0.5em}{}

\pagestyle{fancy}
\fancyhf{}
\fancyhead[L]{\small\color{inkblue}\nouppercase{\leftmark}}
\fancyhead[R]{\small\color{inkblue}\thepage}
\renewcommand{\headrulewidth}{0.4pt}
\renewcommand{\headrule}{\hbox to\headwidth{\color{linegrey}\leaders\hrule height \headrulewidth\hfill}}

% ===== code ==============================================================
\lstdefinestyle{house}{
  backgroundcolor=\color{codebg},
  basicstyle=\ttfamily\footnotesize,
  keywordstyle=\color{inkblue}\bfseries,
  commentstyle=\color{gray}\itshape,
  stringstyle=\color{okgreen},
  numbers=left, numberstyle=\tiny\color{gray}, numbersep=8pt,
  frame=single, rulecolor=\color{linegrey},
  breaklines=true, breakatwhitespace=false,
  showstringspaces=false, tabsize=2,
  captionpos=b, xleftmargin=14pt, framexleftmargin=14pt,
  columns=fullflexible, keepspaces=true,
  upquote=true  % render ' and " as straight quotes, not curly
}
\lstset{style=house}

% ===== callout boxes =====================================================
% \begin{keyidea}{Title} ... \end{keyidea}   - a concept a beginner must grasp
\newtcolorbox{keyidea}[1]{
  enhanced, breakable, colback=softgrey, colframe=accent,
  boxrule=0.4pt, arc=2pt, left=8pt, right=8pt, top=6pt, bottom=6pt,
  title={\bfseries #1}, coltitle=white, colbacktitle=accent, fonttitle=\small
}

% \begin{jargon}{Term} definition \end{jargon}  - define a term at first use
\newtcolorbox{jargon}[1]{
  enhanced, breakable, colback=white, colframe=linegrey,
  boxrule=0.4pt, arc=2pt, left=8pt, right=8pt, top=6pt, bottom=6pt,
  title={\bfseries Jargon: #1}, coltitle=inkblue, colbacktitle=softgrey, fonttitle=\small
}

% \begin{honest}{Title} ... \end{honest}  - brutally honest finding (rule 18)
\newtcolorbox{honest}[1]{
  enhanced, breakable, colback=white, colframe=warnred,
  boxrule=0.6pt, arc=2pt, left=8pt, right=8pt, top=6pt, bottom=6pt,
  title={\bfseries #1}, coltitle=white, colbacktitle=warnred, fonttitle=\small
}

% \begin{provenance}{path/to/file} ... \end{provenance}
% Where a file came from: hand-written, generated by a tool, scaffolded by a
% framework, copied or adapted from an upstream project, ported from an older
% module. State the EVIDENCE (a header comment, a git log entry, an upstream
% URL, a generator command) and label the confidence. If no evidence exists,
% say so in the box. Never guess a lineage into existence.
\newtcolorbox{provenance}[1]{
  enhanced, breakable, colback=white, colframe=okgreen,
  boxrule=0.4pt, arc=2pt, left=8pt, right=8pt, top=6pt, bottom=6pt,
  title={\bfseries Where this came from: #1}, coltitle=white,
  colbacktitle=okgreen, fonttitle=\small
}

% \begin{pseudocode}{Caption} ... \end{pseudocode}
% Stands in for algorithm2e, which is NOT installed. Use \Step / \Indent inside.
\newtcolorbox{pseudocode}[1]{
  enhanced, breakable, colback=codebg, colframe=inkblue,
  boxrule=0.4pt, arc=2pt, left=10pt, right=8pt, top=6pt, bottom=6pt,
  fontupper=\ttfamily\small,
  title={\bfseries\rmfamily #1}, coltitle=white, colbacktitle=inkblue, fonttitle=\small
}
\newcommand{\Step}[1]{\par\noindent #1}
\newcommand{\Indent}[1]{\par\noindent\hspace*{2em}#1}

% ===== tikz house styles =================================================
\tikzset{
  box/.style   = {rectangle, rounded corners=2pt, draw=inkblue, fill=softgrey,
                  text=inkblue, align=center, inner sep=6pt, minimum height=9mm,
                  font=\small},
  extbox/.style= {rectangle, rounded corners=2pt, draw=linegrey, fill=white,
                  text=black!70, align=center, inner sep=6pt, minimum height=9mm,
                  font=\small, dashed},
  store/.style = {cylinder, shape border rotate=90, aspect=0.25, draw=accent,
                  fill=white, text=inkblue, align=center, inner sep=5pt,
                  font=\small},
  dec/.style   = {diamond, aspect=1.6, draw=accent, fill=white, text=inkblue,
                  align=center, inner sep=2pt, font=\scriptsize},
  flow/.style  = {-{Stealth[length=2.4mm]}, draw=inkblue, thick},
  dflow/.style = {-{Stealth[length=2.4mm]}, draw=accent, thick, dashed},
  % align=center is load-bearing: without it, a \\ inside an edge label throws
  % the misleading "perhaps a missing \item" error.
  lbl/.style   = {font=\scriptsize\itshape, text=black!65, fill=white,
                  inner sep=2pt, align=center}
}

% forest default for directory trees
\forestset{
  dirtree/.style={
    for tree={
      font=\ttfamily\small, grow'=0, child anchor=west, parent anchor=south,
      anchor=west, calign=first, inner sep=2pt,
      edge path={\noexpand\path [draw, \forestoption{edge}]
        (!u.south west) +(9pt,0) |- (.child anchor) \forestoption{edge label};},
      before typesetting nodes={if n=1{insert before={[,phantom]}}{}},
      fit=band, before computing xy={l=15pt}
    }
  }
}

% ===== misc ==============================================================
\captionsetup{font=small, labelfont={bf,color=inkblue}}
\setlist[itemize]{leftmargin=1.2em, itemsep=2pt, topsep=3pt}
\setlist[enumerate]{leftmargin=1.4em, itemsep=2pt, topsep=3pt}
\newcommand{\file}[1]{\texttt{\small #1}}
\newcommand{\code}[1]{\texttt{\small #1}}
\setcounter{tocdepth}{2}
\setcounter{secnumdepth}{3}


% ===========================================================================
% Three characters in this repository's source files cannot be typeset by the
% font this document uses, so they are mapped to the closest faithful stand-in.
% The mapping is stated openly in section 1.4 so no reader is misled by it.
%   U+2500  horizontal box-drawing rule   ->  hyphen
%   U+2550  double box-drawing rule       ->  equals sign
%   U+2764  heavy black heart             ->  heart symbol
% U+00B7 (middle dot) renders exactly and needs no substitution.
% ===========================================================================
\DeclareUnicodeCharacter{2500}{-}
\DeclareUnicodeCharacter{2550}{=}
\DeclareUnicodeCharacter{2764}{\ensuremath{\heartsuit}}
\lstset{extendedchars=true,
  literate={·}{{\textperiodcentered}}1
           {─}{{-}}1
           {═}{{=}}1
           {❤}{{\ensuremath{\heartsuit}}}1
}

\title{AI Digitalise\\[2mm]\large Deep Dive: every folder, every file, every line}
\author{Digitalise Agency}
\date{2026-08-11}

\begin{document}
\maketitle
\tableofcontents
\newpage

% ===========================================================================
\section{What this project is}
% ===========================================================================

This repository holds a marketing web page for \textbf{ai.digitalise.agency}: a
single-page site whose job is to persuade a visitor that Digitalise Agency can
build AI systems for them, and then get that visitor to book a call.

It is six files and 1{,}591 lines. There is no framework, no build step, no
package manager, no database and no server-side code. A web browser can open
the page straight off a disk. That makes it one of the simplest projects in the
workspace, and a good one to read first if you have never read code before.

\begin{keyidea}{The one thing to understand before anything else}
This repository does \emph{not} contain the web page that is currently being
served at ai.digitalise.agency. The design here was built, then deliberately
parked, and the machinery that used to publish it was switched off. Something
else is serving that address today. Section 1.2 shows the evidence, and the
whole document is written with that split in mind.
\end{keyidea}

\subsection{The problem it was built to solve}

An agency that sells AI services needs a page that a prospective client can be
sent to. The page has to say what is on offer, show that other people have been
happy, and make booking a conversation a single click. That is the entire
purpose. Everything in these six files serves one of those three jobs.

\begin{figure}[H]\centering
\begin{tikzpicture}[node distance=10mm and 18mm]
  \node[box]                        (v)  {A visitor\\\scriptsize arrives from a link or an ad};
  \node[box, right=of v]            (p)  {The page\\\scriptsize \file{index.html}};
  \node[extbox, right=of p]         (cal){cal.com\\\scriptsize booking calendar};
  \node[box, below=of p]            (out){A booked call\\\scriptsize the only outcome that counts};
  \draw[flow]  (v) -- node[lbl]{opens} (p);
  \draw[dflow] (p) -- node[lbl]{Book a call} (cal);
  \draw[flow]  (cal) |- node[lbl, pos=0.25]{picks a time} (out);
\end{tikzpicture}
\caption{The whole project in one picture. The dashed box is a third-party
service the agency does not run. Everything in this repository exists to move a
visitor from the left box to the bottom box.}
\end{figure}

\subsection{The split between this repository and the live site}

The last commit made to this repository, \code{da508a7} on 2026-07-29, did two
things in one move: it pushed every file down into a folder named
\file{new-design-reference/}, and it renamed the automatic publishing workflow so
that it could no longer run. The person who made that commit explained why in the
commit message itself, and it is worth quoting because it is the single most
important fact about this repository:

\begin{quote}
\emph{``The live site is the generated agency variant built from the portfolio
repo, not this one. Parking this design for reference and turning the workflow
off so it cannot overwrite the real site.''}
\end{quote}

That claim was checked against the live address on 2026-08-11 while this document
was being written. The address answers, and it does not serve this page:

\begin{table}[H]\centering\small
\begin{tabular}{@{}ll@{}}
\toprule
What was checked & What came back \\
\midrule
HTTP status of \code{https://ai.digitalise.agency/} & \code{200}, the site is up \\
Page title served & \emph{Digitalise Agency} followed by \emph{AI \& Full-Stack Engineering} \\
Page title in this repository & \emph{Digitalise AI: Done-for-you AI systems for entrepreneurs} \\
Canonical URL declared by the live page & \code{https://salmanadnan.com/} \\
\bottomrule
\end{tabular}
\caption{Verified by fetching the live address on 2026-08-11. The titles do not
match, so the live page is not the page in this repository, exactly as commit
\code{da508a7} said.}
\end{table}

\begin{honest}{The live site tells search engines to index a different domain}
The page currently served at ai.digitalise.agency declares
\code{<link rel="canonical" href="https://salmanadnan.com/">}. A canonical link
is an instruction to a search engine saying ``the real home of this content is
over there, index that instead of this''. So the live ai.digitalise.agency is
telling Google to credit salmanadnan.com for its content.

That is almost certainly a copy-and-paste accident in whatever repository
produces the live page, not a decision. It is out of scope for \emph{this}
repository, which cannot fix it, but anybody reading this document is likely to
be the person who can. Verified by fetching the page on 2026-08-11.
\end{honest}

\begin{jargon}{Repository, often shortened to ``repo''}
A folder of files whose complete history is recorded, so that every change ever
made can be listed, inspected, and undone. The recording is done by a program
called \emph{git}. When this document says ``this repository'', it means the
folder \file{AI Digitalise} and the history stored alongside it.
\end{jargon}

\subsection{Reading order, and why it is not alphabetical}

Files are explained in the order that makes each one understandable when you
reach it, not in the order a file listing would print them:

\begin{enumerate}
\item \file{new-design-reference/index.html} first, because it is the page
  itself. Every other file exists to serve it.
\item \file{new-design-reference/css/site.css} second, because the page names it
  and is unreadable without it.
\item \file{new-design-reference/js/main.js} third, because it animates what the
  first two files have already put on screen.
\item \file{new-design-reference/netlify.toml} fourth: a hosting setting for a
  host that was never used.
\item \file{new-design-reference/assets/ASSETS.md} fifth: a checklist of work
  that was never finished.
\item \file{.github/workflows-disabled/deploy.yml.disabled} last, because it is
  about publishing, which only makes sense once you know what is being
  published.
\end{enumerate}

\subsection{A note on three characters}

Three characters used in this repository's source files cannot be printed by the
typeface this document uses. Rather than silently drop them, they are printed as
the closest stand-in, listed here so that no quoted line misleads you:

\begin{table}[H]\centering\small
\begin{tabular}{@{}lll@{}}
\toprule
In the real file & Printed here as & Where it appears \\
\midrule
A long horizontal rule (U+2500) & a hyphen & comment banners in the CSS and JavaScript \\
A double horizontal rule (U+2550) & an equals sign & comment banners in the HTML \\
A heavy black heart (U+2764) & a heart symbol & one decorative tile in the HTML \\
\bottomrule
\end{tabular}
\caption{Character substitutions used in every quoted listing in this document.}
\end{table}

The middle dot (U+00B7), which appears in the page footer and in several file
header comments, prints exactly as it does in the file and is not substituted.

% ===========================================================================
\section{The ten-minute mental model}
% ===========================================================================

If you read only one page of this document, read this one.

A web page is made of three kinds of file, and this project has exactly one of
each. That is not a coincidence: it is the oldest and simplest way to build a
web page, and the project chose it deliberately.

\begin{figure}[H]\centering
\begin{tikzpicture}[node distance=11mm and 14mm]
  \node[box]                 (html) {\file{index.html}\\\scriptsize the words and the structure};
  \node[box, right=of html]  (css)  {\file{css/site.css}\\\scriptsize what it all looks like};
  \node[box, right=of css]   (js)   {\file{js/main.js}\\\scriptsize what moves when you scroll};
  \node[extbox, below=of css](br)   {The visitor's web browser\\\scriptsize assembles all three into a page};
  \draw[flow] (html) -- node[lbl]{names} (css);
  \draw[flow] (css)  -- node[lbl]{then} (js);
  \draw[flow] (html) |- (br);
  \draw[flow] (css)  -- (br);
  \draw[flow] (js)   |- (br);
\end{tikzpicture}
\caption{The three-file model. The browser downloads all three and combines them.
Nothing runs on a server: everything happens on the visitor's own machine.}
\end{figure}

\begin{jargon}{HTML}
Short for HyperText Markup Language. It is the file that holds the actual
content of a page (the headings, the paragraphs, the links) wrapped in labels
called \emph{tags} that say what each piece of content \emph{is}. A tag looks
like \code{<h1>} to open and \code{</h1>} to close. HTML says what things are; it
says almost nothing about what they look like.
\end{jargon}

\begin{jargon}{CSS}
Short for Cascading Style Sheets. It is the file that decides appearance: colour,
size, spacing, position, and movement. It works by naming a kind of element (``every
heading'', ``anything labelled \code{btn}'') and listing the appearance rules that
apply to it. The word ``cascading'' means that when two rules disagree, there is a
defined pecking order that decides which one wins.
\end{jargon}

\begin{jargon}{JavaScript}
A programming language that runs inside the browser, after the page has loaded.
It is what makes a page \emph{react}: to a scroll, a click, a keypress. Unlike
HTML and CSS, which describe things, JavaScript gives instructions that happen in
sequence.
\end{jargon}

\subsection{The one paragraph that matters}

\file{index.html} lays out seven stacked bands down the page: a navigation bar, a
full-screen hero with floating tiles behind it, a row of four big numbers, a
five-step pipeline diagram, a grid of six past projects, three client result
cards, a grid of six services, a booking panel, and a footer. \file{site.css}
gives all of that its dark navy look, its rounded corners and its layout, using a
small set of named colours defined once at the top. \file{main.js} adds three
behaviours and nothing else: the navigation bar grows a dark glassy background
once you scroll past 40 pixels, the four big numbers count up from zero when they
scroll into view, and most blocks fade upward into place as they appear. The
``Book a call'' buttons are ordinary links: two of them jump down the page to the
booking panel, and the button inside that panel opens a cal.com calendar in a new
tab.

\begin{keyidea}{Why ``no build step'' is a real design decision}
Most modern web projects are written in one form and then \emph{compiled} into
another before they can be published, by a tool the developer has to install and
maintain. This project skips that entirely: the files you read are byte for byte
the files a browser receives. The benefit is that anybody can open, understand
and change it with nothing but a text editor. The cost is that there is nothing
to catch a mistake before a visitor sees it, and no way to share code with
another page except by copying it.
\end{keyidea}

% ===========================================================================
\section{Background a beginner needs}
% ===========================================================================

This section defines everything the project assumes you already know. If you
know it, skip to section 4.

\subsection{What happens when someone opens the page}

\begin{figure}[H]\centering
\begin{tikzpicture}[node distance=26mm]
  \node[box] (b) {Browser};
  \node[box, right=of b] (s) {Web server\\\scriptsize (Caddy on the VPS)};
  \foreach \n in {b,s} \draw[dashed, linegrey] (\n.south) -- ++(0,-4.0);
  \draw[flow]  ($(b.south)+(0,-0.6)$) -- node[lbl,above]{give me \file{/}} ($(s.south)+(0,-0.6)$);
  \draw[flow]  ($(s.south)+(0,-1.3)$) -- node[lbl,above]{\file{index.html}} ($(b.south)+(0,-1.3)$);
  \draw[flow]  ($(b.south)+(0,-2.0)$) -- node[lbl,above]{give me \file{/css/site.css}} ($(s.south)+(0,-2.0)$);
  \draw[flow]  ($(s.south)+(0,-2.7)$) -- node[lbl,above]{the stylesheet} ($(b.south)+(0,-2.7)$);
  \draw[flow]  ($(b.south)+(0,-3.4)$) -- node[lbl,above]{give me \file{/js/main.js}} ($(s.south)+(0,-3.4)$);
\end{tikzpicture}
\caption{The request sequence. The browser reads \file{index.html} first,
discovers the other two files named inside it, and asks for those next. This is
why the \code{<link>} and \code{<script>} tags matter so much: they are the only
reason the browser knows the other files exist.}
\end{figure}

\begin{jargon}{Web server}
A program that sits on a computer that is always on, waits for requests for
files, and sends those files back. This project is served by a web server called
\emph{Caddy} running on the agency's rented computer. Caddy does no thinking
here: it hands over files exactly as they sit on disk.
\end{jargon}

\begin{jargon}{VPS}
Virtual Private Server: a computer the agency rents from a hosting company and
controls remotely, rather than one sitting in an office. The deploy workflow in
section 12 names its address as \code{46.202.194.85}.
\end{jargon}

\subsection{Tags, attributes, classes and selectors}

Four words do most of the work in this project.

A \textbf{tag} marks a piece of content: \code{<p>} is a paragraph,
\code{<section>} is a major band of a page, \code{<a>} is a link. Most tags come
in pairs that open and close, with content in between.

An \textbf{attribute} is extra information written inside the opening tag, as a
name and a value: in \code{<a href="\#book">}, the attribute \code{href} has the
value \code{\#book}, meaning ``this link points at the thing named book on this
same page''.

A \textbf{class} is a label you invent and attach to any number of elements, with
the \code{class} attribute. It carries no meaning to the browser; it exists so
the CSS can find those elements again.

A \textbf{selector} is how CSS names what it wants to style. \code{.btn} means
``every element labelled \code{btn}''. \code{.btn-primary:hover} means ``every
element labelled \code{btn-primary}, but only while the mouse is over it''.

\begin{jargon}{Element}
One thing on the page, created by a tag: a heading, a button, a box. When this
document says ``element'', it means one of these, not a chemical one.
\end{jargon}

\subsection{Custom properties: colours defined once}

At the top of the stylesheet is a block that gives names to values:
\code{-{}-accent: \#1967ff;} means ``from now on, the name \code{-{}-accent} stands
for this particular blue''. Everywhere else, the stylesheet writes
\code{var(-{}-accent)} instead of the raw colour. Change the definition once and
every use changes with it.

\begin{jargon}{Hex colour}
A colour written as a hash sign and six characters, like \code{\#1967ff}. The six
characters are three pairs giving the amount of red, green and blue, each from
\code{00} (none) to \code{ff} (maximum). \code{\#1967ff} is a strong blue:
little red, a moderate amount of green, almost full blue.
\end{jargon}

\begin{jargon}{Custom property, also called a CSS variable}
A name defined in CSS whose value can be reused anywhere. It always starts with
two hyphens when defined, and is read back with \code{var(...)}. This project
defines twenty of them and reads most of them dozens of times.
\end{jargon}

% ===========================================================================
\section{Repository layout}
% ===========================================================================

\begin{figure}[H]\centering
\begin{forest}dirtree
[AI Digitalise
  [.github
    [workflows-disabled
      [deploy.yml.disabled]
    ]
  ]
  [new-design-reference
    [assets
      [ASSETS.md]
    ]
    [css
      [site.css]
    ]
    [js
      [main.js]
    ]
    [index.html]
    [netlify.toml]
  ]
]
\end{forest}
\caption{The complete repository. Nothing is folded away: these six files are
everything git tracks. The complete list with line counts is in Appendix A.}
\end{figure}

Two things about this shape are unusual and both are deliberate.

First, \textbf{every single file lives under \file{new-design-reference/}}. On a
normal web project, \file{index.html} sits at the top so that a web server
pointed at the folder finds it immediately. Here it is one level down, which
means a server pointed at this folder would find no home page at all. That is the
parking described in section 1.2: the folder name is a message to the next
person, saying ``this is a reference, not a live site''.

Second, \textbf{the automation folder is misspelled on purpose}. GitHub only runs
files it finds in a folder called exactly \file{.github/workflows}. This one is
called \file{.github/workflows-disabled}, and the file inside it ends in
\file{.disabled} rather than \file{.yml}. Either change alone would be enough to
stop it running. Both were made, which tells you the intent was not to pause it
briefly.

\begin{jargon}{Hidden folder}
A folder whose name begins with a dot, like \file{.github}. On Mac and Linux
these are hidden from ordinary file listings by default. Nothing else is special
about them; the dot is just a convention meaning ``tooling lives here, not
content''.
\end{jargon}

% ===========================================================================
\section{Architecture}
% ===========================================================================

There is no architecture in the usual sense, because there is no running program
to have one. What exists instead is a set of relationships between six files, and
those are worth drawing once before walking them.

\begin{figure}[H]\centering
\begin{tikzpicture}[node distance=12mm and 20mm]
  \node[box]                          (html) {\file{index.html}\\\scriptsize 597 lines};
  \node[box, above right=of html]     (css)  {\file{css/site.css}\\\scriptsize 768 lines};
  \node[box, below right=of html]     (js)   {\file{js/main.js}\\\scriptsize 79 lines};
  \node[extbox, left=of html]         (fonts){Google Fonts\\\scriptsize two typefaces};
  \node[extbox, below=of html]        (cal)  {cal.com\\\scriptsize \code{davonex/30min}};
  \node[box, above=of css]            (assets){\file{assets/ASSETS.md}\\\scriptsize a checklist, not code};
  \draw[flow]  (html) -- node[lbl]{\code{<link>}} (css);
  \draw[flow]  (html) -- node[lbl]{\code{<script>}} (js);
  \draw[dflow] (html) -- node[lbl]{\code{<link>}} (fonts);
  \draw[dflow] (html) -- node[lbl]{\code{href}} (cal);
  \draw[flow, dashed] (assets) -- node[lbl]{describes} (css);
\end{tikzpicture}
\caption{What names what. Solid arrows are files this project owns; dashed boxes
are services on somebody else's computers. \file{netlify.toml} and the deploy
workflow are not shown because nothing on the page refers to them: they are
instructions for machines that publish the page, not part of it.}
\end{figure}

The important structural point is that \textbf{the arrows only go one way}.
\file{index.html} knows the names of the other two files. Neither of them knows
anything about it: \file{site.css} styles class names that it hopes exist, and
\file{main.js} looks for elements that may or may not be there. Every one of its
lookups is written to do nothing quietly if the element is missing. That is why
the page does not break when a section is deleted, and also why a typo in a class
name fails silently rather than loudly.

\begin{honest}{Silent failure is the price of this simplicity}
Nothing anywhere checks that the class names in \file{index.html} and the
selectors in \file{site.css} agree. Rename \code{.hero-badge} in one file and not
the other and the page still loads, still scores as ``working'', and simply looks
wrong. Section 14 lists three places where this has already happened.
\end{honest}

% ===========================================================================
\section{Data flow, end to end}
% ===========================================================================

\begin{figure}[H]\centering
\begin{tikzpicture}[node distance=8mm and 12mm]
  \node[box]                    (a) {Browser downloads\\\file{index.html}};
  \node[box, right=of a]        (b) {Reads \code{<head>},\\requests fonts and CSS};
  \node[box, right=of b]        (c) {Draws the page\\top to bottom};
  \node[box, below=of c]        (d) {Reaches \code{<script>}\\at the very end};
  \node[box, left=of d]         (e) {\file{main.js} runs,\\attaches watchers};
  \node[box, left=of e]         (f) {Visitor scrolls,\\watchers fire};
  \draw[flow] (a) -- (b);
  \draw[flow] (b) -- (c);
  \draw[flow] (c) -- (d);
  \draw[flow] (d) -- (e);
  \draw[flow] (e) -- (f);
\end{tikzpicture}
\caption{The full life of one page view. Note that the JavaScript is the last
thing to run, because its \code{<script>} tag is the last thing in the document.
That ordering is deliberate and is explained in section 7.3.}
\end{figure}

There is no data flow beyond this. No form is submitted, nothing is saved,
nothing is sent anywhere. The only outbound step a visitor can trigger is
following the cal.com link, at which point they have left this page entirely.

\begin{honest}{Nothing measures whether the page works}
There is no analytics of any kind: no counter, no tag manager, no event when
somebody presses ``Book a call''. So the question this page exists to answer,
``does it get bookings'', cannot be answered from the page itself. Any evidence
would have to come from the cal.com side. Whether that is a privacy choice or an
oversight is not recorded anywhere in the repository, so this document does not
guess.
\end{honest}

% ===========================================================================
\section{The file-by-file walkthrough}
% ===========================================================================

Every file git tracks gets its own subsection below, quoted from its first line
to its last. Nothing is summarized away and nothing is skipped.

% ===========================================================================
\subsection{\file{new-design-reference/index.html}}
% ===========================================================================

\textbf{HTML, 597 lines.} The page itself: every word a visitor reads, and the
structure that holds those words. If you delete every other file in this
repository, this one still opens in a browser and still says everything the
project wants to say. It would simply look plain.

\begin{provenance}{new-design-reference/index.html}
Hand-written. It arrived complete in the very first commit, \code{0d77d51} on
2026-07-28 (``Launch ai.digitalise.agency''), at the path \file{index.html}, and
was moved to its present path in \code{da508a7} on 2026-07-29 with no change to
its contents. There is no generator banner, no framework scaffold, and no
\file{package.json} that could have produced it.

Its icons are a different matter and are worth naming precisely. Every icon in
the file is an inline SVG drawing whose path data matches the Heroicons outline
set, a widely used free icon library: the shapes, the \code{24 24} view box, the
\code{stroke-linecap} and \code{stroke-linejoin} attributes and the
\code{1.5}-based coordinates are that library's house conventions. \textbf{The
file carries no attribution comment saying so}, so this is stated as an inference
from the path data, not as a documented fact. \textbf{Confidence: high for
hand-written, inferred for the icon origin.}
\end{provenance}

\subsubsection*{Why it exists}

Without this file there is no page. The stylesheet styles class names that would
belong to nothing, and the JavaScript would search for elements that do not
exist and quietly do nothing, which is exactly how it is written to behave.

\subsubsection*{Line by line}

\paragraph{Lines 1 to 17: the head, which the visitor never sees.}

\begin{lstlisting}[language=HTML,firstnumber=1,caption={\file{index.html}, lines 1--17}]
<!DOCTYPE html>
<html lang="en">
<head>
<meta charset="UTF-8" />
<meta name="viewport" content="width=device-width, initial-scale=1.0" />
<title>Digitalise AI: Done-for-you AI systems for entrepreneurs</title>
<meta name="description" content="Done-for-you AI systems that answer, follow up, and publish, so your business keeps growing without adding headcount." />
<meta property="og:title" content="Digitalise AI" />
<meta property="og:description" content="Done-for-you AI systems that answer, follow up, and publish." />
<meta property="og:type" content="website" />
<meta name="theme-color" content="#05070d" />

<link rel="preconnect" href="https://fonts.googleapis.com" />
<link rel="preconnect" href="https://fonts.gstatic.com" crossorigin />
<link href="https://fonts.googleapis.com/css2?family=Space+Grotesk:wght@500;600;700&family=Inter:wght@400;500;600;700&display=swap" rel="stylesheet" />
<link rel="stylesheet" href="/css/site.css" />
</head>
\end{lstlisting}

Line 1 is not a tag at all. \code{<!DOCTYPE html>} is a declaration that tells the
browser ``interpret this as modern HTML''. It is a historical artefact: without
it browsers fall back to a compatibility mode that imitates the bugs of software
from the 1990s. It must be the first thing in the file, and it takes no closing
partner.

Line 2 opens the document and declares its language as English with
\code{lang="en"}. That attribute is small and does real work: screen readers use
it to choose a pronunciation, and search engines use it to decide who to show the
page to.

Line 3 opens the \code{<head>}. Everything between here and line 17 is
information \emph{about} the page rather than content \emph{of} it. Nothing in
this block appears on screen.

Line 4 declares the character encoding as UTF-8, which is the scheme that lets a
file contain any character in any language. This matters concretely here: the
page contains a middle dot and a heart symbol, and without this line those would
render as mojibake. It should appear as early as possible, because the browser
has to guess the encoding until it is told.

\begin{jargon}{Meta tag}
A tag inside \code{<head>} carrying a fact about the page rather than content for
the page. It has no closing partner and displays nothing. Different meta tags are
read by completely different audiences: browsers, search engines, and social
networks each look for their own.
\end{jargon}

Line 5 is the viewport declaration and it is the single line that makes the page
work on a phone. \code{width=device-width} tells the browser to treat the page
width as the real width of the device, and \code{initial-scale=1.0} says do not
zoom out to start with. Without it, a phone pretends to be about 980 pixels wide
and shrinks everything, which is why old sites look like tiny unreadable
newspapers on a phone. Every media query in the stylesheet depends on this line
being present.

Line 6 is the \code{<title>}: the text in the browser tab and the blue headline
in a search result. Note it for later, because it is how section 1.2 proved this
page is not the live one.

Line 7 is the description, which search engines often show as the grey text under
the title. It is 118 characters, comfortably inside the roughly 155 that get
displayed.

Lines 8 to 10 are Open Graph tags, a convention invented by Facebook and now read
by nearly every chat app and social network to build the preview card when
somebody pastes a link. Note that \code{og:title} here says only ``Digitalise
AI'', shorter than the real title on line 6, which is a sensible choice: preview
cards truncate.

\begin{honest}{The Open Graph set is incomplete in the way that matters most}
The three tags present are \code{og:title}, \code{og:description} and
\code{og:type}. The two missing ones are \code{og:image} and \code{og:url}.
Without \code{og:image}, a pasted link produces a preview card with no picture,
which is the difference between a link that gets clicked and one that does not.
\file{ASSETS.md} lists \code{og-image.png} as required before launch, and it was
never created. See section 7.5.
\end{honest}

Line 11 sets \code{theme-color} to \code{\#05070d}, the same near-black used as
the page background. Mobile browsers paint their address bar this colour, so the
browser chrome blends into the page instead of framing it in white.

Line 12 is blank. Blank lines have no effect on how HTML renders; they are purely
for the human reading the file, and here the author uses them to separate the
page's identity from its external resources.

Lines 13 and 14 are \code{preconnect} hints. They tell the browser: you will
shortly need to talk to these two Google servers, so start the handshake now
rather than waiting. Establishing a secure connection takes several network round
trips, and doing that early rather than at the moment of need typically saves a
few hundred milliseconds. The \code{crossorigin} attribute on line 14 is required
on the second one, because font files are fetched under stricter rules than
stylesheets, and a preconnect made without it would open a connection the font
request cannot reuse. That subtlety catches out experienced developers
constantly.

Line 15 actually fetches the fonts: two typefaces, Space Grotesk in three weights
and Inter in four. \code{display=swap} at the end is the important part. It tells
the browser to render text immediately in a fallback typeface and swap in the
real one when it arrives, rather than hiding the text for up to three seconds
while it waits. The visible cost is that text reflows slightly; the benefit is
that a visitor on a slow connection can read the page at all.

\begin{jargon}{Font weight}
How heavy a typeface looks, on a numeric scale where 400 is normal and 700 is
bold. Requesting \code{wght@500;600;700} downloads three separate files. Each
weight requested is another download, which is why the list is short.
\end{jargon}

Line 16 links the stylesheet. This is the path discussed in section 11: the
leading slash in \code{/css/site.css} means ``from the top of the site'', which
is correct when served by a web server and wrong when the file is opened directly
off a disk.

Line 17 closes the head.

\paragraph{Lines 18 to 35: the navigation bar.}

\begin{lstlisting}[language=HTML,firstnumber=18,caption={\file{index.html}, lines 18--35}]
<body>

<!-- ══════════════════════════════════════════════
     NAV
══════════════════════════════════════════════ -->
<nav class="nav" aria-label="Site navigation">
  <div class="nav-inner">
    <a href="/" class="nav-logo" aria-label="Digitalise AI home">
      <svg viewBox="0 0 28 28" fill="none" aria-hidden="true">
        <rect width="28" height="28" rx="8" fill="#1967ff"/>
        <path d="M7 14c0-3.866 3.134-7 7-7s7 3.134 7 7-3.134 7-7 7" stroke="#fff" stroke-width="2.2" stroke-linecap="round"/>
        <circle cx="14" cy="14" r="2.5" fill="#fff"/>
      </svg>
      digitalise <span class="logo-ai">AI</span>
    </a>
    <a href="#book" class="btn btn-primary">Book a call</a>
  </div>
</nav>
\end{lstlisting}

Line 18 opens \code{<body>}: everything from here to line 596 is what a visitor
actually sees.

Line 19 is blank, and lines 20 to 22 are a comment banner. Anything between
\code{<!-{}-} and \code{-{}->} is ignored entirely by the browser. This file uses
seven of these banners as signposts, one per major band of the page, each drawn
with a rule of repeated double-line characters. They are the author's table of
contents, and they are the reason a 597-line file stays navigable.

Line 23 opens the navigation bar. Two things are happening. \code{<nav>} is a
\emph{semantic} tag: it does not look different from a plain \code{<div>}, but it
announces to screen readers and search engines that this is the site's
navigation. \code{aria-label="Site navigation"} gives that region a spoken name,
so a screen-reader user hearing a list of page regions hears something meaningful.

\begin{jargon}{ARIA}
Accessible Rich Internet Applications: a set of attributes, all starting with
\code{aria-}, that describe a page to assistive software such as screen readers.
They change nothing visually. \code{aria-label} supplies a name;
\code{aria-hidden="true"} tells assistive software to skip something entirely
because it is purely decorative.
\end{jargon}

Line 24 opens a \code{<div>} labelled \code{nav-inner}. A \code{<div>} is a
generic container with no meaning at all; it exists purely to be styled. The
reason for a second container inside \code{<nav>} is a common layout pattern: the
outer element spans the full window width so its background can stretch edge to
edge, while the inner one is capped at a maximum width and centred, so the logo
and button line up with the content further down the page.

Line 25 opens a link to \code{/}, the site's home page. Since this \emph{is} the
home page it reloads the page, which is the conventional behaviour of a logo.
Its \code{aria-label} matters here: the link's visible content is a drawing plus
the words ``digitalise AI'', and without the label a screen reader would announce
the drawing awkwardly.

Lines 26 to 30 are the logo, drawn in SVG.

\begin{jargon}{SVG}
Scalable Vector Graphics: a picture described as shapes and coordinates rather
than as a grid of coloured dots. Because it is instructions rather than pixels,
it stays perfectly sharp at any size, and because it is text it can sit directly
inside the HTML file with no extra download.
\end{jargon}

\code{viewBox="0 0 28 28"} sets up a private coordinate system 28 units square;
every coordinate inside is relative to that, so the whole drawing can be scaled to
any size by CSS without touching the numbers. \code{fill="none"} sets the default
to no fill, which each shape then overrides. \code{aria-hidden="true"} hides it
from screen readers, correctly: the link already has a text label, and announcing
the drawing too would be duplication.

Line 27 draws a 28 by 28 rounded square (\code{rx="8"} is the corner radius) in
the brand blue. Line 28 draws the swoosh: \code{d} is a path description in a
compact language where \code{M7 14} means ``move to the point (7,14)'' and the
\code{c} and \code{s} commands that follow are curve instructions. It is written
to be generated, not read, and nobody edits it by hand. Line 29 puts a white dot
in the middle. Line 30 closes the drawing.

Line 31 is the wordmark. ``digitalise'' is plain text and ``AI'' is wrapped in a
\code{<span>} labelled \code{logo-ai} so the stylesheet can colour those two
letters blue. A \code{<span>} is the inline equivalent of a \code{<div>}: no
meaning, no visual effect, purely a hook for styling.

Line 32 closes the logo link. Line 33 is the navigation's only button, a link to
\code{\#book}. Note it is a link and not a \code{<button>}, which is correct: it
navigates somewhere rather than performing an action. It carries two classes,
\code{btn} for the shared button shape and \code{btn-primary} for the filled blue
variant. That two-class pattern repeats throughout the file and is explained in
section 7.2.

Lines 34 and 35 close the containers.

\paragraph{Lines 36 to 57: the hero opens, and the first floating tile.}

\begin{lstlisting}[language=HTML,firstnumber=36,caption={\file{index.html}, lines 36--57}]


<!-- ══════════════════════════════════════════════
     HERO
══════════════════════════════════════════════ -->
<section class="hero" aria-label="Hero">

  <!-- animated background tiles -->
  <div class="hero-bg" aria-hidden="true">

    <!-- Tile 1: LinkedIn Video -->
    <div class="tile tile-blue">
      <div class="tile-thumb">
        <svg viewBox="0 0 24 24" fill="none" stroke="currentColor" stroke-width="1.8" aria-hidden="true">
          <path stroke-linecap="round" stroke-linejoin="round" d="M15.75 10.5l4.72-4.72a.75.75 0 011.28.53v11.38a.75.75 0 01-1.28.53l-4.72-4.72M4.5 18.75h9a2.25 2.25 0 002.25-2.25v-9a2.25 2.25 0 00-2.25-2.25h-9A2.25 2.25 0 002.25 7.5v9a2.25 2.25 0 002.25 2.25z"/>
        </svg>
      </div>
      <div class="tile-meta">
        <span class="tile-label">LinkedIn Video</span>
        <span class="tile-metric">847K views</span>
      </div>
    </div>
\end{lstlisting}

Lines 36 to 39 are a blank line and the second comment banner.

Line 40 opens the hero, the full-screen band at the top of the page.
\code{<section>} is another semantic tag, meaning ``a distinct part of the
document''.

Line 42 is a comment, and line 43 opens the container for the eight drifting
cards. \code{aria-hidden="true"} on this container is the single most important
accessibility decision in the file: it hides all eight tiles, and every number
inside them, from screen readers. That is correct, because they are decoration.
It is also the reason the invented numbers inside them are a presentation problem
rather than an accessibility one.

Lines 46 to 57 are the first tile. \code{class="tile tile-blue"} again uses the
two-class pattern: \code{tile} carries the shared card shape and the floating
animation, \code{tile-blue} adds only the blue gradient behind the icon.

Inside, line 47 opens the thumbnail area, and lines 48 to 51 draw a video-camera
icon. \code{stroke="currentColor"} on line 48 is worth understanding: it means
``draw the outline in whatever text colour applies here'', so the icon inherits
its colour from CSS instead of hard-coding one. That single attribute is why one
icon definition can appear in white on one background and blue on another.

Lines 53 to 56 are the tile's caption row: a label on the left and a metric on
the right.

\begin{honest}{The tile metrics are invented}
\code{847K views} is not a measurement. Neither are the other five numbers in
these tiles: 234K impressions, 1.2M views, 512K views, 1.8M views and 19K
reactions. Nothing in this repository or its history connects them to a real
account or a real campaign, and they appeared fully formed in the first commit
alongside the rest of the design.

They are decorative and hidden from assistive software, and a designer would call
them realistic placeholder content, which is a normal thing to put in a mockup.
But a sighted visitor cannot tell decoration from a claim, and reads them as
results this agency has produced. That is worth raising specifically because the
agency's main website was built to a stated standard of using no fabricated
numbers.
\end{honest}

\paragraph{Lines 58 to 70: the second tile, built from bars instead of an icon.}

\begin{lstlisting}[language=HTML,firstnumber=58,caption={\file{index.html}, lines 58--70}]

    <!-- Tile 2: Post tile -->
    <div class="tile tile-orange">
      <div class="tile-lines" style="padding:10px 0;">
        <div class="tile-line"></div>
        <div class="tile-line"></div>
        <div class="tile-line"></div>
      </div>
      <div class="tile-meta">
        <span class="tile-label">LinkedIn Post</span>
        <span class="tile-metric">234K impr.</span>
      </div>
    </div>
\end{lstlisting}

This tile represents a text post rather than a video, so instead of an icon it
shows three grey bars standing in for lines of text. Each \code{<div
class="tile-line">} is completely empty: it has no content at all, and exists
only so that CSS can give it a height, a width and a rounded grey background.
This is a standard technique for drawing an abstract shape without an image, and
the widths (100\%, 85\%, 60\%) are set in the stylesheet at lines 216 to 218 to
suggest a paragraph ending mid-line.

\begin{honest}{An inline style that fights the stylesheet}
\code{style="padding:10px 0;"} on line 61 sets the spacing directly on the
element. The stylesheet also has a rule for this class, \code{.tile-lines
\{ padding: 0; \}} at line 209, and an inline style always wins over a stylesheet
rule. So that rule never has any effect, in either of the two places the class is
used. One of the two should go; the stylesheet rule is the dead one.
\end{honest}

\begin{jargon}{Inline style}
Appearance written directly on an element with a \code{style} attribute, rather
than in the stylesheet. It overrides almost anything the stylesheet says, which
makes it useful for one-offs and dangerous as a habit: rules written in one file
stop being able to control what appears in the other.
\end{jargon}

\paragraph{Lines 71 to 80: the third tile, a statistic.}

\begin{lstlisting}[language=HTML,firstnumber=71,caption={\file{index.html}, lines 71--80}]

    <!-- Tile 3: Stats tile -->
    <div class="tile" style="background:var(--s1);">
      <div class="tile-stat-num">1.2M</div>
      <div class="tile-stat-lbl">Views this month</div>
      <div class="tile-meta" style="margin-top:12px;">
        <span class="tile-label">AI Campaign</span>
        <span class="tile-metric">+340%</span>
      </div>
    </div>
\end{lstlisting}

This tile has no icon and no bars, just a large number, a caption and the usual
meta row. Note \code{style="background:var(-{}-s1);"} on line 73: an inline style
that reads a custom property defined in the stylesheet. It works, and it is
redundant, because the \code{.tile} rule at line 155 already sets exactly that
background. Two more tiles (lines 96 and 121) repeat the same redundant inline
background.

\paragraph{Lines 81 to 93: the fourth tile.}

\begin{lstlisting}[language=HTML,firstnumber=81,caption={\file{index.html}, lines 81--93}]

    <!-- Tile 4: Video -->
    <div class="tile tile-teal">
      <div class="tile-thumb">
        <svg viewBox="0 0 24 24" fill="none" stroke="currentColor" stroke-width="1.8" aria-hidden="true">
          <path stroke-linecap="round" stroke-linejoin="round" d="M5.25 5.653c0-.856.917-1.398 1.667-.986l11.54 6.348a1.125 1.125 0 010 1.971l-11.54 6.347a1.125 1.125 0 01-1.667-.985V5.653z"/>
        </svg>
      </div>
      <div class="tile-meta">
        <span class="tile-label">AI Video</span>
        <span class="tile-metric">512K views</span>
      </div>
    </div>
\end{lstlisting}

Structurally identical to tile 1, with two differences: the colour class is
\code{tile-teal}, and the icon is a play triangle rather than a video camera.
That same play-triangle path appears six more times in this file, once on each
portfolio card. Nothing shares it: the same 180 characters of path data are
repeated seven times in total, which is the cost of having no templating system.
It is a real duplication, and it is also the correct trade for a project that
deliberately has no build step.

\paragraph{Lines 94 to 105: the fifth tile, a status badge.}

\begin{lstlisting}[language=HTML,firstnumber=94,caption={\file{index.html}, lines 94--105}]

    <!-- Tile 5: Published badge -->
    <div class="tile" style="background:var(--s1);">
      <div style="display:flex;align-items:center;gap:8px;margin-bottom:10px;">
        <svg viewBox="0 0 24 24" fill="none" stroke="#34d399" stroke-width="2" width="20" height="20" aria-hidden="true">
          <path stroke-linecap="round" stroke-linejoin="round" d="M9 12.75L11.25 15 15 9.75M21 12a9 9 0 11-18 0 9 9 0 0118 0z"/>
        </svg>
        <span style="font-family:var(--font-d);font-size:13px;font-weight:700;color:#34d399;">Published</span>
      </div>
      <div class="tile-stat-num" style="font-size:22px;">47 posts</div>
      <div class="tile-stat-lbl">this week, zero manual work</div>
    </div>
\end{lstlisting}

This is the most heavily inline-styled element in the file: line 97 lays out a
row, line 101 styles the word ``Published'', and line 103 overrides a font size.
None of it has a class, so none of it can be restyled from the stylesheet.

Note the green on lines 98 and 101 is written as the literal \code{\#34d399},
even though the stylesheet defines exactly that colour as \code{-{}-green} at line
15. Change the token and this tile would not follow. The same literal appears
again at line 276 in the pipeline. This is the clearest illustration in the file
of what inline styles cost: a colour system that only mostly applies.

\paragraph{Lines 106 to 118: the sixth tile.}

\begin{lstlisting}[language=HTML,firstnumber=106,caption={\file{index.html}, lines 106--118}]

    <!-- Tile 6: Violet video -->
    <div class="tile tile-violet">
      <div class="tile-thumb">
        <svg viewBox="0 0 24 24" fill="none" stroke="currentColor" stroke-width="1.8" aria-hidden="true">
          <path stroke-linecap="round" stroke-linejoin="round" d="M15.75 10.5l4.72-4.72a.75.75 0 011.28.53v11.38a.75.75 0 01-1.28.53l-4.72-4.72M4.5 18.75h9a2.25 2.25 0 002.25-2.25v-9a2.25 2.25 0 00-2.25-2.25h-9A2.25 2.25 0 002.25 7.5v9a2.25 2.25 0 002.25 2.25z"/>
        </svg>
      </div>
      <div class="tile-meta">
        <span class="tile-label">AI Video</span>
        <span class="tile-metric">1.8M views</span>
      </div>
    </div>
\end{lstlisting}

Identical in structure to tile 1, and using the identical video-camera icon path,
with a violet gradient and a different number. This is the third of four tiles
built from the same template.

\paragraph{Lines 119 to 130: the seventh tile, a schedule.}

\begin{lstlisting}[language=HTML,firstnumber=119,caption={\file{index.html}, lines 119--130}]

    <!-- Tile 7: Schedule -->
    <div class="tile" style="background:var(--s1);">
      <div style="display:flex;align-items:center;gap:8px;margin-bottom:10px;">
        <svg viewBox="0 0 24 24" fill="none" stroke="var(--accentH)" stroke-width="1.8" width="18" height="18" aria-hidden="true">
          <path stroke-linecap="round" stroke-linejoin="round" d="M6.75 3v2.25M17.25 3v2.25M3 18.75V7.5a2.25 2.25 0 012.25-2.25h13.5A2.25 2.25 0 0121 7.5v11.25m-18 0A2.25 2.25 0 005.25 21h13.5A2.25 2.25 0 0021 18.75m-18 0v-7.5A2.25 2.25 0 015.25 9h13.5A2.25 2.25 0 0121 11.25v7.5"/>
        </svg>
        <span style="font-family:var(--font-d);font-size:12px;font-weight:600;color:var(--accentH);">Scheduled</span>
      </div>
      <div class="tile-stat-num" style="font-size:22px;">34 posts</div>
      <div class="tile-stat-lbl">queued this week</div>
    </div>
\end{lstlisting}

The mirror image of tile 5: same layout, calendar icon instead of a tick, and
``queued'' instead of ``published''. Together the pair tells a small story in the
background of the hero, that work is both finished and lined up. Note that here
the colour \emph{is} read from a token, \code{var(-{}-accentH)} on lines 123 and
126, unlike tile 5 which used a literal. The two tiles were written minutes apart
and are inconsistent with each other.

\paragraph{Lines 131 to 145: the eighth tile, and the end of the backdrop.}

\begin{lstlisting}[language=HTML,firstnumber=131,caption={\file{index.html}, lines 131--145}]

    <!-- Tile 8: Reactions -->
    <div class="tile tile-blue">
      <div class="tile-lines" style="padding:10px 0;">
        <div class="tile-line"></div>
        <div class="tile-line"></div>
        <div class="tile-line" style="width:50%;"></div>
      </div>
      <div class="tile-meta">
        <span class="tile-label">Viral Post</span>
        <span class="tile-metric">19K ❤</span>
      </div>
    </div>

  </div><!-- /hero-bg -->
\end{lstlisting}

The last tile repeats the bar pattern from tile 2, overriding the third bar's
width to 50\% inline so the two text tiles do not look identical. Line 141 is
where the heart character appears; as section 1.4 explained, this document prints
it as a heart symbol.

Line 145 closes the tile container, with a trailing comment naming what was
closed. That habit is worth copying: in a file with this much nesting, a bare
\code{</div>} tells the reader nothing.

\begin{keyidea}{Why the tile order in the file is not the order on screen}
The eight tiles are written top to bottom in the file, but they appear scattered
around the hero. Nothing here places them. Each one is positioned by the
stylesheet, using rules that count children: \code{.tile:nth-child(1)},
\code{(2)} and so on, at lines 162 to 169, each setting a corner and a delay
before that tile starts drifting. \textbf{This means the order of these blocks in
the HTML is load-bearing.} Move tile 5 above tile 3 in this file, without touching
the CSS, and both tiles jump to new positions on screen and start animating at
different times.
\end{keyidea}

\paragraph{Lines 146 to 182: the overlay and the hero's actual content.}

\begin{lstlisting}[language=HTML,firstnumber=146,caption={\file{index.html}, lines 146--182}]

  <div class="hero-overlay" aria-hidden="true"></div>

  <!-- hero content -->
  <div class="hero-content">
    <div class="hero-badge">
      <span class="hero-badge-dot"></span>
      Digitalise Agency
    </div>

    <h1 class="hero-h1">
      AI systems built<br>
      to grow while you <em>sleep.</em>
    </h1>

    <div class="hero-actions">
      <a href="#book" class="btn btn-primary">
        Book a call
        <svg width="16" height="16" viewBox="0 0 24 24" fill="none" stroke="currentColor" stroke-width="2.5" aria-hidden="true">
          <path stroke-linecap="round" stroke-linejoin="round" d="M13.5 4.5L21 12m0 0l-7.5 7.5M21 12H3"/>
        </svg>
      </a>
      <a href="#portfolio" class="btn btn-ghost">See the work</a>
    </div>

    <div class="hero-social-proof">
      <div class="avatars" aria-hidden="true">
        <div class="avatar av1">K</div>
        <div class="avatar av2">H</div>
        <div class="avatar av3">M</div>
        <div class="avatar av4">A</div>
      </div>
      <p class="hero-proof-text"><strong>17+ businesses</strong> growing with Digitalise AI</p>
    </div>
  </div>

</section>
\end{lstlisting}

Line 147 is an empty \code{<div>} that exists purely to darken the middle of the
hero so the headline stays readable over the drifting tiles. The stylesheet fills
it with two stacked gradients (lines 232 to 243). An empty element used as a
visual layer is a common and legitimate technique.

Lines 150 to 153 are the small pill above the headline. \code{hero-badge-dot} is
another empty element, styled into a small circle that pulses.

Lines 156 to 159 are the headline, in an \code{<h1>}. A page should have exactly
one \code{<h1>}, and this page does: it is the document's title as far as a screen
reader or a search engine is concerned. \code{<br>} on line 157 forces a line
break so the headline splits where the designer wants rather than where the window
happens to be narrow.

\begin{honest}{\code{<em>} is used for colour, not for emphasis}
Line 158 wraps ``sleep.'' in \code{<em>}, and the stylesheet then removes the
italics that tag normally implies and colours the word blue instead (lines 285 to
288). \code{<em>} means emphasis, and a screen reader may change its intonation
for it. Using it as a colour hook works, and it slightly misrepresents the
content to assistive software; a \code{<span>} with a class would have been the
honest choice. This is a small point and the intent is plainly visual.
\end{honest}

Lines 161 to 169 hold the two hero buttons. The first is the primary blue one,
containing both text and an arrow icon; it points at \code{\#book}. The second is
the ghost variant, pointing at \code{\#portfolio}. Both are ordinary links, and
both are intercepted by the JavaScript at line 70 to scroll smoothly rather than
jump.

Lines 171 to 179 are the social proof row: four overlapping circles bearing the
letters K, H, M and A, and a line of text. The circles are \code{aria-hidden}
because they carry no information.

\begin{honest}{The four avatar initials do not correspond to anything}
K, H, M and A do not match the three clients named further down the page (Hi-Tec
Bearings, DigiMarvixx, Synergy Event, whose own avatars use K, D and S). They are
placeholder initials that were never replaced. The ``17+ businesses'' claim beside
them matches the \code{data-target="17"} in the metrics band, so that number at
least is internally consistent; whether it is externally true cannot be
established from this repository.
\end{honest}

Line 182 closes the hero section.

\paragraph{Lines 183 to 210: the metrics band.}

\begin{lstlisting}[language=HTML,firstnumber=183,caption={\file{index.html}, lines 183--210}]


<!-- ══════════════════════════════════════════════
     METRICS BAND
     UPDATE: replace data-target values with real numbers
══════════════════════════════════════════════ -->
<section class="metrics" aria-label="Results">
  <div class="container">
    <div class="metrics-grid stagger">
      <div class="metric-item" style="--i:0">
        <div class="metric-num" data-target="17" data-suffix="+">0</div>
        <div class="metric-label">Projects Delivered</div>
      </div>
      <div class="metric-item" style="--i:1">
        <div class="metric-num" data-target="1.6" data-suffix="M+">0</div>
        <div class="metric-label">Organic Views for Clients</div>
      </div>
      <div class="metric-item" style="--i:2">
        <div class="metric-num" data-target="10" data-suffix="K+">0</div>
        <div class="metric-label">Users Served</div>
      </div>
      <div class="metric-item" style="--i:3">
        <div class="metric-num" data-target="300" data-suffix="%">0</div>
        <div class="metric-label">Avg. Traffic Growth</div>
      </div>
    </div>
  </div>
</section>
\end{lstlisting}

This is the first band that the JavaScript genuinely drives, and it introduces
the two mechanisms the rest of the page reuses.

The first is the \textbf{data attribute}. Look at line 193: the element's visible
content is the single character \code{0}, and the real number lives in
\code{data-target="17"}, with \code{data-suffix="+"} beside it.

\begin{jargon}{Data attribute}
Any attribute whose name starts with \code{data-}. HTML ignores it completely and
displays nothing; it exists so that JavaScript can attach its own information to
an element. \file{main.js} reads these two at lines 20 and 21 to know what number
to count towards and what to print after it.
\end{jargon}

Why is the visible content \code{0} rather than \code{17}? Because that is what a
visitor sees before the counter runs. It is also what a visitor sees
\emph{forever} if JavaScript is switched off or fails to load, which is the
weakness of this pattern: the page would show four zeros and claim nothing. The
robust alternative is to write the real number in the HTML and have the script
count towards it from zero, which fails gracefully. This file chose the fragile
form.

The second mechanism is the \textbf{stagger}. Line 191 gives the grid the class
\code{stagger}, and each child carries \code{style="-{}-i:0"}, \code{-{}-i:1} and so
on. That custom property is multiplied by 90 milliseconds in the stylesheet (line
90) to delay each child's fade-in, so the four items appear in sequence rather
than together. Setting a custom property inline is the cleanest way to pass a
per-element number to CSS, and it is used four times in this file.

Note that this band still carries an \code{UPDATE} comment on line 187 telling the
reader to replace the values with real numbers, even though commit \code{2c66a7a}
claimed to have done exactly that. See section 14.

\paragraph{Lines 211 to 235: the pipeline section and its first node.}

\begin{lstlisting}[language=HTML,firstnumber=211,caption={\file{index.html}, lines 211--235}]


<!-- ══════════════════════════════════════════════
     HOW IT WORKS (PIPELINE)
══════════════════════════════════════════════ -->
<section class="pipeline-section" aria-label="How it works">
  <div class="container">
    <div class="reveal">
      <p class="section-kicker">How it works</p>
      <h2 class="section-title">From idea to live<br>in one click.</h2>
    </div>

    <div class="pipeline" aria-label="Content production pipeline">

      <!-- Node 1 -->
      <div class="pipeline-node">
        <div class="pipeline-icon">
          <svg viewBox="0 0 24 24" fill="none" stroke="currentColor" stroke-width="1.8" aria-hidden="true">
            <path stroke-linecap="round" stroke-linejoin="round" d="M12 18v-5.25m0 0a6.01 6.01 0 001.5-.189m-1.5.189a6.01 6.01 0 01-1.5-.189m3.75 7.478a12.06 12.06 0 01-4.5 0m3.75 2.355a14.998 14.998 0 01-3 0M14.25 18v-.192c0-.983.658-1.823 1.508-2.316a7.5 7.5 0 10-7.517 0c.85.493 1.509 1.333 1.509 2.316V18"/>
          </svg>
        </div>
        <span class="pipeline-lbl">Your Idea</span>
      </div>

      <div class="pipeline-connector" aria-hidden="true"></div>
\end{lstlisting}

Lines 218 to 221 establish a pattern used by all four remaining sections: a
wrapper with the class \code{reveal}, containing a small uppercase ``kicker'' line
and a large heading. The \code{reveal} class is what the JavaScript watches for at
line 56, fading the pair upward into place when they scroll into view.

Line 220 is an \code{<h2>}. Heading levels form the document's outline: the single
\code{<h1>} in the hero, then an \code{<h2>} per major section. Getting that order
right is how a screen-reader user navigates a long page, and this file gets it
right throughout.

Line 223 opens the pipeline, and line 235 is the first connector: an entirely
empty element that the stylesheet turns into a short line between two nodes, and
then fills with a blue gradient when the section becomes visible. The nodes and
connectors alternate, which is why the stylesheet targets odd children for nodes
and even children for connectors (lines 400 to 426).

\paragraph{Lines 236 to 285: the remaining four pipeline nodes.}

\begin{lstlisting}[language=HTML,firstnumber=236,caption={\file{index.html}, lines 236--285}]

      <!-- Node 2 -->
      <div class="pipeline-node">
        <div class="pipeline-icon">
          <svg viewBox="0 0 24 24" fill="none" stroke="currentColor" stroke-width="1.8" aria-hidden="true">
            <path stroke-linecap="round" stroke-linejoin="round" d="M16.862 4.487l1.687-1.688a1.875 1.875 0 112.652 2.652L10.582 16.07a4.5 4.5 0 01-1.897 1.13L6 18l.8-2.685a4.5 4.5 0 011.13-1.897l8.932-8.931zm0 0L19.5 7.125M18 14v4.75A2.25 2.25 0 0115.75 21H5.25A2.25 2.25 0 013 18.75V8.25A2.25 2.25 0 015.25 6H10"/>
          </svg>
        </div>
        <span class="pipeline-lbl">AI Script</span>
      </div>

      <div class="pipeline-connector" aria-hidden="true"></div>

      <!-- Node 3 -->
      <div class="pipeline-node">
        <div class="pipeline-icon">
          <svg viewBox="0 0 24 24" fill="none" stroke="currentColor" stroke-width="1.8" aria-hidden="true">
            <path stroke-linecap="round" stroke-linejoin="round" d="M15.75 10.5l4.72-4.72a.75.75 0 011.28.53v11.38a.75.75 0 01-1.28.53l-4.72-4.72M4.5 18.75h9a2.25 2.25 0 002.25-2.25v-9a2.25 2.25 0 00-2.25-2.25h-9A2.25 2.25 0 002.25 7.5v9a2.25 2.25 0 002.25 2.25z"/>
          </svg>
        </div>
        <span class="pipeline-lbl">AI Video</span>
      </div>

      <div class="pipeline-connector" aria-hidden="true"></div>

      <!-- Node 4 -->
      <div class="pipeline-node">
        <div class="pipeline-icon">
          <svg viewBox="0 0 24 24" fill="none" stroke="currentColor" stroke-width="1.8" aria-hidden="true">
            <path stroke-linecap="round" stroke-linejoin="round" d="M6 12L3.269 3.126A59.768 59.768 0 0121.485 12 59.77 59.77 0 013.27 20.876L5.999 12zm0 0h7.5"/>
          </svg>
        </div>
        <span class="pipeline-lbl">Auto-Post</span>
      </div>

      <div class="pipeline-connector" aria-hidden="true"></div>

      <!-- Node 5 -->
      <div class="pipeline-node">
        <div class="pipeline-icon" style="background:rgba(52,211,153,.08);border-color:rgba(52,211,153,.25);">
          <svg viewBox="0 0 24 24" fill="none" stroke="#34d399" stroke-width="1.8" aria-hidden="true">
            <path stroke-linecap="round" stroke-linejoin="round" d="M9 12.75L11.25 15 15 9.75M21 12a9 9 0 11-18 0 9 9 0 0118 0z"/>
          </svg>
        </div>
        <span class="pipeline-lbl" style="color:#34d399;">Live</span>
      </div>

    </div><!-- /pipeline -->
  </div>
</section>
\end{lstlisting}

Four more nodes, each built to the identical shape: a container, an icon box, an
SVG, and a label. The icons in order are a lightbulb (an idea), a pencil on a
document (a script), a video camera (a video, reusing the same path as tile 1 and
tile 6), and a paper plane (posting).

Node 5 is the only one that differs, and it differs three times over: it overrides
its background and border colour inline on line 275, hard-codes the same green in
the SVG stroke on line 276, and colours its label inline on line 280. Everything
is written to make the final step green rather than blue, and none of it uses the
\code{-{}-green} token that exists for exactly this purpose. The same criticism as
tile 5 applies, and this is the fourth and last place the literal \code{\#34d399}
appears.

The five steps read as a claim about the product: an idea becomes a script,
becomes a video, gets posted automatically, and is live. That is the same pipeline
the Gantry project in this workspace actually implements, so the page is
advertising a real capability rather than an invented one, though nothing in this
repository connects the two.

\paragraph{Lines 286 to 318: the portfolio section and its first card.}

\begin{lstlisting}[language=HTML,firstnumber=286,caption={\file{index.html}, lines 286--318}]


<!-- ══════════════════════════════════════════════
     PORTFOLIO
     UPDATE: replace gradient .pt* classes with real video poster images.
     Add poster="path/to/thumb.jpg" to each <video> or swap divs for <img>.
══════════════════════════════════════════════ -->
<section class="portfolio-section" id="portfolio" aria-label="Work portfolio">
  <div class="container">
    <div class="reveal" style="margin-bottom:48px;">
      <p class="section-kicker">The work</p>
      <h2 class="section-title">Real content.<br>Real results.</h2>
    </div>

    <div class="portfolio-grid stagger">

      <div class="port-card" style="--i:0">
        <div class="port-thumb pt1">
          <div class="port-badge">847K views</div>
          <div class="port-thumb-overlay">
            <div class="port-play">
              <svg viewBox="0 0 24 24" fill="white" aria-hidden="true"><path d="M5.25 5.653c0-.856.917-1.398 1.667-.986l11.54 6.348a1.125 1.125 0 010 1.971l-11.54 6.347a1.125 1.125 0 01-1.667-.985V5.653z"/></svg>
            </div>
          </div>
        </div>
        <div class="port-info">
          <div class="port-platform">
            <svg viewBox="0 0 24 24" fill="none" stroke="currentColor" stroke-width="1.8" aria-hidden="true"><path stroke-linecap="round" stroke-linejoin="round" d="M13.5 4.5L21 12m0 0l-7.5 7.5M21 12H3"/></svg>
            E-commerce
          </div>
          <div class="port-title">Hi-Tec Bearings: 5,000+ SKUs, 300% organic traffic growth</div>
        </div>
      </div>
\end{lstlisting}

Line 293 carries \code{id="portfolio"}, which is what makes the hero's ``See the
work'' button work. An \code{id} is a unique name for one element, and a link to
\code{\#portfolio} scrolls to whatever carries that name.

\begin{jargon}{id, and how it differs from class}
An \code{id} names exactly one element on the page and must be unique. A
\code{class} labels any number of them. Use an \code{id} when something needs to
be linked to or found individually; use a \code{class} when you want to style a
group. This page has exactly two ids, \code{portfolio} and \code{book}, and both
exist to be jumped to.
\end{jargon}

Each card is built from two halves. The top half, \code{port-thumb}, is where a
real video thumbnail would go; today it is a gradient supplied by the class
\code{pt1}, defined at line 504 of the stylesheet. Over it sit a badge in the
corner and an overlay containing a play button, which the stylesheet keeps
invisible until the mouse is over the card. The bottom half, \code{port-info},
holds a small platform label with an arrow icon and the project title.

The six projects named across these cards are Hi-Tec Bearings, DigiMarvixx,
Synergy Event, IBA GPT, Prestige Sleep and IBA CED. Unlike the hero tiles, these
are specific, named and checkable, which is a meaningfully different kind of
claim.

\begin{honest}{The play buttons promise videos that do not exist}
Every card has a play button that appears on hover, and none of them plays
anything: the card is not a link and has no click handler. The stylesheet even
sets \code{cursor: pointer} on it (line 465), so the pointer changes to a hand,
promising an action that never happens. \file{ASSETS.md} confirms the intent was
real video, and the video was never supplied. As it stands the hover behaviour is
a small deception, and the cheapest honest fix is to remove the play overlay
until the videos exist.
\end{honest}

\paragraph{Lines 319 to 354: portfolio cards two and three.}

\begin{lstlisting}[language=HTML,firstnumber=319,caption={\file{index.html}, lines 319--354}]

      <div class="port-card" style="--i:1">
        <div class="port-thumb pt2">
          <div class="port-badge">3x leads</div>
          <div class="port-thumb-overlay">
            <div class="port-play">
              <svg viewBox="0 0 24 24" fill="white" aria-hidden="true"><path d="M5.25 5.653c0-.856.917-1.398 1.667-.986l11.54 6.348a1.125 1.125 0 010 1.971l-11.54 6.347a1.125 1.125 0 01-1.667-.985V5.653z"/></svg>
            </div>
          </div>
        </div>
        <div class="port-info">
          <div class="port-platform">
            <svg viewBox="0 0 24 24" fill="none" stroke="currentColor" stroke-width="1.8" aria-hidden="true"><path stroke-linecap="round" stroke-linejoin="round" d="M13.5 4.5L21 12m0 0l-7.5 7.5M21 12H3"/></svg>
            Web + WhatsApp
          </div>
          <div class="port-title">DigiMarvixx: 3x monthly WhatsApp leads, 60% mobile traffic</div>
        </div>
      </div>

      <div class="port-card" style="--i:2">
        <div class="port-thumb pt3">
          <div class="port-badge">1.2s load</div>
          <div class="port-thumb-overlay">
            <div class="port-play">
              <svg viewBox="0 0 24 24" fill="white" aria-hidden="true"><path d="M5.25 5.653c0-.856.917-1.398 1.667-.986l11.54 6.348a1.125 1.125 0 010 1.971l-11.54 6.347a1.125 1.125 0 01-1.667-.985V5.653z"/></svg>
            </div>
          </div>
        </div>
        <div class="port-info">
          <div class="port-platform">
            <svg viewBox="0 0 24 24" fill="none" stroke="currentColor" stroke-width="1.8" aria-hidden="true"><path stroke-linecap="round" stroke-linejoin="round" d="M13.5 4.5L21 12m0 0l-7.5 7.5M21 12H3"/></svg>
            Events
          </div>
          <div class="port-title">Synergy Event: 800 participants, 98/100 Lighthouse score</div>
        </div>
      </div>
\end{lstlisting}

Structurally identical to card one, differing only in the stagger index
(\code{-{}-i:1} and \code{-{}-i:2}), the gradient class (\code{pt2}, \code{pt3}), the
badge text, the platform label and the title. Note that the same arrow icon is
used as the ``platform'' icon on every card regardless of what the platform
actually is, so it is decorative rather than informative.

``98/100 Lighthouse score'' on line 352 is a specific, verifiable claim.
Lighthouse is Google's page-quality auditing tool, and 98 out of 100 is a strong
result.

\paragraph{Lines 355 to 412: portfolio cards four, five and six.}

\begin{lstlisting}[language=HTML,firstnumber=355,caption={\file{index.html}, lines 355--412}]

      <div class="port-card" style="--i:3">
        <div class="port-thumb pt4">
          <div class="port-badge">AI chat</div>
          <div class="port-thumb-overlay">
            <div class="port-play">
              <svg viewBox="0 0 24 24" fill="white" aria-hidden="true"><path d="M5.25 5.653c0-.856.917-1.398 1.667-.986l11.54 6.348a1.125 1.125 0 010 1.971l-11.54 6.347a1.125 1.125 0 01-1.667-.985V5.653z"/></svg>
            </div>
          </div>
        </div>
        <div class="port-info">
          <div class="port-platform">
            <svg viewBox="0 0 24 24" fill="none" stroke="currentColor" stroke-width="1.8" aria-hidden="true"><path stroke-linecap="round" stroke-linejoin="round" d="M13.5 4.5L21 12m0 0l-7.5 7.5M21 12H3"/></svg>
            AI Assistant
          </div>
          <div class="port-title">IBA GPT: AI assistant answering questions about IBA university</div>
        </div>
      </div>

      <div class="port-card" style="--i:4">
        <div class="port-thumb pt5">
          <div class="port-badge">24/7 sales</div>
          <div class="port-thumb-overlay">
            <div class="port-play">
              <svg viewBox="0 0 24 24" fill="white" aria-hidden="true"><path d="M5.25 5.653c0-.856.917-1.398 1.667-.986l11.54 6.348a1.125 1.125 0 010 1.971l-11.54 6.347a1.125 1.125 0 01-1.667-.985V5.653z"/></svg>
            </div>
          </div>
        </div>
        <div class="port-info">
          <div class="port-platform">
            <svg viewBox="0 0 24 24" fill="none" stroke="currentColor" stroke-width="1.8" aria-hidden="true"><path stroke-linecap="round" stroke-linejoin="round" d="M13.5 4.5L21 12m0 0l-7.5 7.5M21 12H3"/></svg>
            Health + E-commerce
          </div>
          <div class="port-title">Prestige Sleep: clinic booking and product shop in one site</div>
        </div>
      </div>

      <div class="port-card" style="--i:5">
        <div class="port-thumb pt6">
          <div class="port-badge">4 dashboards</div>
          <div class="port-thumb-overlay">
            <div class="port-play">
              <svg viewBox="0 0 24 24" fill="white" aria-hidden="true"><path d="M5.25 5.653c0-.856.917-1.398 1.667-.986l11.54 6.348a1.125 1.125 0 010 1.971l-11.54 6.347a1.125 1.125 0 01-1.667-.985V5.653z"/></svg>
            </div>
          </div>
        </div>
        <div class="port-info">
          <div class="port-platform">
            <svg viewBox="0 0 24 24" fill="none" stroke="currentColor" stroke-width="1.8" aria-hidden="true"><path stroke-linecap="round" stroke-linejoin="round" d="M13.5 4.5L21 12m0 0l-7.5 7.5M21 12H3"/></svg>
            Analytics
          </div>
          <div class="port-title">IBA CED: analytics platform tracking entrepreneurship outcomes</div>
        </div>
      </div>

    </div><!-- /portfolio-grid -->
  </div>
</section>
\end{lstlisting}

The final three cards complete the set of six, following the same template with
indices 3, 4 and 5 and gradients \code{pt4} through \code{pt6}. Two of the six
projects are for the same institution (IBA GPT and the analytics platform), which
is the only repetition among them.

Line 410 closes the grid with a naming comment, and lines 411 and 412 close the
container and the section.

\paragraph{Lines 413 to 472: the client proof band.}

\begin{lstlisting}[language=HTML,firstnumber=413,caption={\file{index.html}, lines 413--472}]


<!-- ══════════════════════════════════════════════
     CLIENT PROOF
     UPDATE: replace initials and numbers with real client data
══════════════════════════════════════════════ -->
<section class="proof-section" aria-label="Client results">
  <div class="container">
    <div class="reveal" style="margin-bottom:48px;">
      <p class="section-kicker">Client results</p>
      <h2 class="section-title">Numbers that<br>speak for themselves.</h2>
    </div>

    <div class="proof-grid stagger">

      <div class="proof-card" style="--i:0">
        <div class="proof-header">
          <div class="proof-avatar pa1">K</div>
          <div>
            <div class="proof-name">Hi-Tec Bearings</div>
            <div class="proof-company">Industrial E-commerce</div>
          </div>
        </div>
        <div class="proof-metric">
          <span class="proof-num">300%</span>
          <span class="proof-unit">organic traffic growth</span>
        </div>
      </div>

      <div class="proof-card" style="--i:1">
        <div class="proof-header">
          <div class="proof-avatar pa2">D</div>
          <div>
            <div class="proof-name">DigiMarvixx</div>
            <div class="proof-company">Digital Marketing</div>
          </div>
        </div>
        <div class="proof-metric">
          <span class="proof-num">3x</span>
          <span class="proof-unit">monthly WhatsApp leads</span>
        </div>
      </div>

      <div class="proof-card" style="--i:2">
        <div class="proof-header">
          <div class="proof-avatar pa3">S</div>
          <div>
            <div class="proof-name">Synergy Event</div>
            <div class="proof-company">Events & Conferences</div>
          </div>
        </div>
        <div class="proof-metric">
          <span class="proof-num">800</span>
          <span class="proof-unit">participants, 98/100 Lighthouse</span>
        </div>
      </div>

    </div>
  </div>
</section>
\end{lstlisting}

Three cards, each with a coloured circle bearing an initial, a client name, an
industry, and one headline number. The numbers repeat what the portfolio cards
already said about the same three clients, which is deliberate: repetition is how
a landing page makes a claim stick.

Note line 430: the avatar for Hi-Tec Bearings shows the letter \textbf{K}, not H.
That is almost certainly a leftover from the placeholder set in the hero (K, H, M,
A), and it is the clearest surviving trace of the \code{UPDATE} comment on line
417 having been only partly acted on.

Line 461 contains \code{Events \& Conferences}. The ampersand is written as a bare
character rather than as the escape \code{\&amp;}. Browsers forgive this, and it
is technically invalid: \code{\&} begins a character reference in HTML. It works
here only because what follows is a space, so no valid reference can be formed.
It is worth knowing as a habit to avoid rather than a bug to fix.

\begin{jargon}{Character reference, or entity}
A way of writing characters that would otherwise be interpreted as markup:
\code{\&amp;} for an ampersand, \code{\&lt;} for a less-than sign. They exist
because those characters have structural meaning in HTML.
\end{jargon}

\paragraph{Lines 473 to 549: the services grid.}

\begin{lstlisting}[language=HTML,firstnumber=473,caption={\file{index.html}, lines 473--495}]


<!-- ══════════════════════════════════════════════
     SERVICES
══════════════════════════════════════════════ -->
<section class="services-section" aria-label="Services">
  <div class="container">
    <div class="reveal" style="margin-bottom:48px;">
      <p class="section-kicker">What we do</p>
      <h2 class="section-title">Everything your content<br>operation needs.</h2>
    </div>

    <div class="services-grid stagger">

      <div class="service-item" style="--i:0">
        <div class="service-icon">
          <svg viewBox="0 0 24 24" fill="none" stroke="currentColor" stroke-width="1.8" aria-hidden="true">
            <path stroke-linecap="round" stroke-linejoin="round" d="M15.75 10.5l4.72-4.72a.75.75 0 011.28.53v11.38a.75.75 0 01-1.28.53l-4.72-4.72M4.5 18.75h9a2.25 2.25 0 002.25-2.25v-9a2.25 2.25 0 00-2.25-2.25h-9A2.25 2.25 0 002.25 7.5v9a2.25 2.25 0 002.25 2.25z"/>
          </svg>
        </div>
        <div class="service-name">AI Personal Brand</div>
        <span class="service-tag">Your voice, published</span>
      </div>
\end{lstlisting}

The services grid follows the established shape: a revealed heading pair, then a
staggered grid. Each of the six items is an icon box, a name, and a small
uppercase tag line. The first item's icon is, once again, the same video-camera
path from tile 1.

The remaining five items are structurally identical and are quoted in full below,
since the coverage rule admits no ``and so on''.

\begin{lstlisting}[language=HTML,firstnumber=496,caption={\file{index.html}, lines 496--549: services two to six}]

      <div class="service-item" style="--i:1">
        <div class="service-icon">
          <svg viewBox="0 0 24 24" fill="none" stroke="currentColor" stroke-width="1.8" aria-hidden="true">
            <path stroke-linecap="round" stroke-linejoin="round" d="M21 21l-5.197-5.197m0 0A7.5 7.5 0 105.196 5.196a7.5 7.5 0 0010.607 10.607z"/>
          </svg>
        </div>
        <div class="service-name">AI SEO & Content</div>
        <span class="service-tag">Rank for what they ask</span>
      </div>

      <div class="service-item" style="--i:2">
        <div class="service-icon">
          <svg viewBox="0 0 24 24" fill="none" stroke="currentColor" stroke-width="1.8" aria-hidden="true">
            <path stroke-linecap="round" stroke-linejoin="round" d="M8.625 12a.375.375 0 11-.75 0 .375.375 0 01.75 0zm0 0H8.25m4.125 0a.375.375 0 11-.75 0 .375.375 0 01.75 0zm0 0H12m4.125 0a.375.375 0 11-.75 0 .375.375 0 01.75 0zm0 0h-.375M21 12c0 4.556-4.03 8.25-9 8.25a9.764 9.764 0 01-2.555-.337A5.972 5.972 0 015.41 20.97a5.969 5.969 0 01-.474-.065 4.48 4.48 0 00.978-2.025c.09-.457-.133-.901-.467-1.226C3.93 16.178 3 14.189 3 12c0-4.556 4.03-8.25 9-8.25s9 3.694 9 8.25z"/>
          </svg>
        </div>
        <div class="service-name">AI Chat Agent</div>
        <span class="service-tag">Reply in seconds</span>
      </div>

      <div class="service-item" style="--i:3">
        <div class="service-icon">
          <svg viewBox="0 0 24 24" fill="none" stroke="currentColor" stroke-width="1.8" aria-hidden="true">
            <path stroke-linecap="round" stroke-linejoin="round" d="M2.25 6.75c0 8.284 6.716 15 15 15h2.25a2.25 2.25 0 002.25-2.25v-1.372c0-.516-.351-.966-.852-1.091l-4.423-1.106c-.44-.11-.902.055-1.173.417l-.97 1.293c-.282.376-.769.542-1.21.38a12.035 12.035 0 01-7.143-7.143c-.162-.441.004-.928.38-1.21l1.293-.97c.363-.271.527-.734.417-1.173L6.963 3.102a1.125 1.125 0 00-1.091-.852H4.5A2.25 2.25 0 002.25 4.5v2.25z"/>
          </svg>
        </div>
        <div class="service-name">AI Voice Agent</div>
        <span class="service-tag">Answer calls 24/7</span>
      </div>

      <div class="service-item" style="--i:4">
        <div class="service-icon">
          <svg viewBox="0 0 24 24" fill="none" stroke="currentColor" stroke-width="1.8" aria-hidden="true">
            <path stroke-linecap="round" stroke-linejoin="round" d="M21.75 6.75v10.5a2.25 2.25 0 01-2.25 2.25h-15a2.25 2.25 0 01-2.25-2.25V6.75m19.5 0A2.25 2.25 0 0019.5 4.5h-15a2.25 2.25 0 00-2.25 2.25m19.5 0v.243a2.25 2.25 0 01-1.07 1.916l-7.5 4.615a2.25 2.25 0 01-2.36 0L3.32 8.91a2.25 2.25 0 01-1.07-1.916V6.75"/>
          </svg>
        </div>
        <div class="service-name">AI Email Outreach</div>
        <span class="service-tag">Follow up until they respond</span>
      </div>

      <div class="service-item" style="--i:5">
        <div class="service-icon">
          <svg viewBox="0 0 24 24" fill="none" stroke="currentColor" stroke-width="1.8" aria-hidden="true">
            <path stroke-linecap="round" stroke-linejoin="round" d="M9.53 16.122a3 3 0 00-5.78 1.128 2.25 2.25 0 01-2.4 2.245 4.5 4.5 0 008.4-2.245c0-.399-.078-.78-.22-1.128zm0 0a15.998 15.998 0 003.388-1.62m-5.043-.025a15.994 15.994 0 011.622-3.395m3.42 3.42a15.995 15.995 0 004.764-4.648l3.876-5.814a1.151 1.151 0 00-1.597-1.597L14.146 6.32a15.996 15.996 0 00-4.649 4.763m3.42 3.42a6.776 6.776 0 00-3.42-3.42"/>
          </svg>
        </div>
        <div class="service-name">AI Ad Creatives</div>
        <span class="service-tag">On-brand at volume</span>
      </div>

    </div>
  </div>
</section>
\end{lstlisting}

The six services are AI Personal Brand, AI SEO and Content, AI Chat Agent, AI
Voice Agent, AI Email Outreach and AI Ad Creatives. Their icons in order are a
video camera, a magnifying glass, a speech bubble, a telephone handset, an
envelope and a paintbrush, each one an obvious visual match for its service, which
is why they need no labels of their own.

Line 503 contains another unescaped ampersand in ``AI SEO \& Content'', the same
point as line 461.

Note the grid gets its dividing lines by a trick rather than by borders: the
stylesheet gives the grid a 2-pixel gap with the border colour as the grid's own
background, then gives each cell the page background (lines 608 to 618). The gaps
show the layer beneath, producing hairlines with no border rules at all. It is
elegant, and it is the kind of thing that looks broken if you change one of the
two backgrounds without the other.

\paragraph{Lines 550 to 580: the booking panel.}

\begin{lstlisting}[language=HTML,firstnumber=550,caption={\file{index.html}, lines 550--580}]


<!-- ══════════════════════════════════════════════
     BOOK A CALL
     UPDATE: replace YOUR_CAL_LINK with Salman's Cal.com or Calendly URL.
     For inline Cal.com embed, swap the anchor tag for the Cal embed snippet.
══════════════════════════════════════════════ -->
<section class="book-section" id="book" aria-label="Book a call">
  <div class="container">
    <div class="book-wrap reveal">
      <p class="book-eyebrow">Ready to grow?</p>
      <h2 class="book-title">30 minutes.<br>No slides.</h2>
      <p class="book-sub">Pick a time. We talk about your content goals.</p>
      <div class="book-cta-row">
        <!-- UPDATE: replace href with your Cal.com or Calendly link -->
        <a href="https://cal.com/davonex/30min" target="_blank" rel="noopener" class="btn btn-primary" style="padding:16px 36px;font-size:17px;">
          Book a call
          <svg width="18" height="18" viewBox="0 0 24 24" fill="none" stroke="currentColor" stroke-width="2.5" aria-hidden="true">
            <path stroke-linecap="round" stroke-linejoin="round" d="M13.5 4.5L21 12m0 0l-7.5 7.5M21 12H3"/>
          </svg>
        </a>
      </div>
      <p class="book-detail">
        <svg viewBox="0 0 24 24" fill="none" stroke="currentColor" stroke-width="2" aria-hidden="true">
          <path stroke-linecap="round" stroke-linejoin="round" d="M12 6v6h4.5m4.5 0a9 9 0 11-18 0 9 9 0 0118 0z"/>
        </svg>
        Free · <span>30 minutes · Zoom</span>
      </p>
    </div>
  </div>
</section>
\end{lstlisting}

This is the destination of every ``Book a call'' button on the page, via
\code{id="book"} on line 557. It is the point of the whole document.

Line 559 combines two classes on one element, \code{book-wrap} for the panel's
appearance and \code{reveal} so the JavaScript fades it in. That is the only place
in the file where a layout class and a behaviour class share an element.

Line 565 is the actual booking link, and three of its attributes matter.
\code{target="\_blank"} opens it in a new tab, so the visitor does not lose the
page. \code{rel="noopener"} is a security measure that must accompany it: without
it, the newly opened page is handed a reference back to this one and can navigate
it somewhere else, a trick known as tab-nabbing. Its presence here is correct and
is the kind of detail that is usually forgotten.

\begin{honest}{The booking link is a placeholder that was never replaced}
\code{cal.com/davonex/30min} is the calendar of Davonex, a different agency
project, not a Digitalise AI calendar. Commit \code{4dffc12} on 2026-07-28 set
this deliberately, and the same placeholder calendar is recorded as an outstanding
item on the agency's main website too, so it is a known shared loose end rather
than a slip here.

Two things around it are worse than the link itself. The comment on line 554 still
tells the reader to search for \code{YOUR\_CAL\_LINK}, which no longer exists
anywhere in the file, and \file{ASSETS.md} repeats that same stale instruction.
Somebody following the written instructions would search, find nothing, and
conclude the work was done.
\end{honest}

Lines 572 to 577 are the reassurance line under the button: a clock icon, then
``Free'', then the duration and the platform. The middle dots here are the
U+00B7 character discussed in section 1.4, and they render exactly.

\paragraph{Lines 581 to 597: the footer, and the script tag.}

\begin{lstlisting}[language=HTML,firstnumber=581,caption={\file{index.html}, lines 581--597}]


<!-- ══════════════════════════════════════════════
     FOOTER
══════════════════════════════════════════════ -->
<footer>
  <div class="container">
    <p>
      <a href="/">digitalise.agency</a> &nbsp;·&nbsp;
      <a href="mailto:hello@digitalise.agency">hello@digitalise.agency</a>
    </p>
  </div>
</footer>

<script src="/js/main.js"></script>
</body>
</html>
\end{lstlisting}

The footer is one line of text: a link home and an email address, separated by a
middle dot padded with \code{\&nbsp;} on each side.

\begin{jargon}{Non-breaking space}
\code{\&nbsp;} is a space that a browser will not use as a place to wrap a line,
and which is not collapsed when several spaces sit together. Here it does the
second job: ordinary spaces around the dot would collapse to one, and these keep
the separator visually padded.
\end{jargon}

Line 589 links to \code{/}, meaning the site's own home page. The visible text
says \code{digitalise.agency}, which is a different site from
\code{ai.digitalise.agency}. So the footer link appears to promise the main agency
site and in fact reloads this one. On the live site that would be a small
misdirection; here it is another loose end.

Line 590 is a \code{mailto:} link, which opens the visitor's email program with
the address filled in.

Line 595 loads the JavaScript, and \textbf{where it sits is the whole point}. It
is the last thing before \code{</body>}, so by the time the browser reaches it,
every element on the page already exists. \file{main.js} begins by searching for
elements immediately, with no wrapper telling it to wait, and that only works
because of this placement. Move this one line up into the \code{<head>} and the
script would run against an empty document, find nothing, and the page would
quietly lose all three of its behaviours with no error message.

Lines 596 and 597 close the body and the document.

\subsubsection*{What it connects to}

\textbf{This file names:} \file{new-design-reference/css/site.css} (line 16),
\file{new-design-reference/js/main.js} (line 595), Google Fonts (lines 13 to 15),
and cal.com (line 565).

\textbf{This file is named by:} nothing inside the repository. The deploy workflow
copies the whole folder rather than referring to files individually, and
\file{netlify.toml} publishes the folder as a whole. \file{ASSETS.md} discusses it
at length but cannot affect it.

\textbf{Internal links:} the hero and navigation both point at \code{\#book}
(lines 33 and 162), and the hero also points at \code{\#portfolio} (line 168).
Those two ids, on lines 557 and 293, are the only link targets in the file.

\subsubsection*{How it breaks}

\begin{itemize}
\item \textbf{Rename a class here without renaming it in the stylesheet} and that
  element loses its appearance silently. Nothing warns you.
\item \textbf{Move the \code{<script>} tag into the \code{<head>}} and all three
  behaviours stop working, with no error in the console. The metrics band would
  then permanently display four zeros.
\item \textbf{Reorder the eight hero tiles} and they jump to different positions,
  because the stylesheet positions them by their number, not by their class.
\item \textbf{Open the file directly off a disk} and it renders unstyled, because
  of the leading slash on line 16. This looks like catastrophic breakage and is
  not. See section 11.
\item \textbf{Disable JavaScript} and the page is fully readable but static: the
  metrics show zeros, nothing fades in, and the anchor links jump instead of
  gliding. Only the metrics are a real loss.
\item \textbf{Delete the id on line 557} and every ``Book a call'' button on the
  page silently does nothing, because \file{main.js} returns early when it cannot
  find a link's target.
\end{itemize}

\subsubsection*{Honest findings for this file}

Collected here, each already argued above: the invented tile metrics (lines 55,
68, 74, 91, 116, 141); the missing \code{og:image} and favicon; the mismatched
avatar initial on line 429; five stale \code{UPDATE} comments (lines 187, 290,
417, 554, 564); the placeholder booking link on line 565; two unescaped ampersands
(lines 461 and 503); \code{<em>} used as a colour hook on line 158; play buttons
that play nothing; a footer link whose text and destination disagree; the
literal \code{\#34d399} used four times where a token exists; and one inline
padding that permanently overrides a stylesheet rule.

% ===========================================================================
\subsection{\file{new-design-reference/css/site.css}}
% ===========================================================================

\textbf{CSS, 768 lines.} Every visual decision on the page: the colours, the
typefaces, the spacing, the layout, the movement, and how all of it rearranges on
a narrow screen. It is the longest file in the repository and, in the judgement
of this document, the most valuable thing in it.

\begin{provenance}{new-design-reference/css/site.css}
Hand-written, and it says where its ideas come from in its own second line:
\code{Design tokens mirror the Digitalise Agency system (tokens.css in Gantry).}
That is a direct statement by the author that the colour and spacing names were
taken from another project in this workspace, \code{gantry-digitalise}, which is
strong evidence of a deliberate shared design language rather than a coincidence.

The file arrived complete in the first commit, \code{0d77d51}, and has not been
edited since except by being moved. No generator banner, no vendor prefix block
of the sort a preprocessor emits, and no minified section: it is written to be
read. \textbf{Confidence: certain for hand-written, high for the token lineage,
which rests on the author's own comment rather than on a file comparison.}
\end{provenance}

\subsubsection*{Why it exists}

Without this file the page still contains every word it needs, and looks like a
1994 document: black serif text on white, everything in one column, no colour, no
layout. Nothing would be broken and nothing would be persuasive. All eight hero
tiles would appear stacked in a list, since it is CSS and not HTML that scatters
them.

\subsubsection*{Line by line}

\paragraph{Lines 1 to 25: the header, and the token block.}

\begin{lstlisting}[firstnumber=1,caption={\file{css/site.css}, lines 1--25}]
/* ai.digitalise.agency · site.css
   Design tokens mirror the Digitalise Agency system (tokens.css in Gantry). */

/* ─── TOKENS ─── */
:root {
  --s0:    #05070d;
  --s1:    #0d1220;
  --s2:    #151c2e;
  --ink1:  #f4f7fc;
  --ink2:  #c3ccdd;
  --ink3:  #93a0b8;
  --accent:  #1967ff;
  --accentH: #3f82ff;
  --accent2: #ff6719;
  --green:   #34d399;
  --border:  #222b40;
  --borderS: #1a2233;
  --font-d: 'Space Grotesk', ui-sans-serif, system-ui, sans-serif;
  --font-b: 'Inter', ui-sans-serif, system-ui, sans-serif;
  --r:   16px;
  --rsm: 10px;
  --ease: cubic-bezier(.22,1,.36,1);
  --maxw: 1320px;
  --shadow: 0 8px 40px rgba(0,0,0,.45);
}
\end{lstlisting}

Lines 1 and 2 are a comment. In CSS, anything between \code{/*} and \code{*/} is
ignored, and unlike HTML there is only this one comment form. These two lines
carry the file's entire provenance, which is why they matter out of all
proportion to their length.

Line 4 is the first of eleven section banners drawn with box-drawing characters.
As with the HTML, they are the author's navigation aid.

Line 5 opens the token block. \code{:root} is a selector that matches the
document's outermost element, which is \code{<html>}. Custom properties defined
there are inherited by every element on the page, which is precisely why token
blocks are written against \code{:root}: it is the one place a definition reaches
everything.

The twenty tokens fall into six groups, and understanding the naming scheme is
most of understanding the file:

\begin{table}[H]\centering\small
\begin{tabular}{@{}llp{6.4cm}@{}}
\toprule
Tokens & Line & What the group is for \\
\midrule
\code{--s0} to \code{--s2}      & 6--8   & Surfaces, darkest to lightest. \code{s0} is the page, \code{s1} is a card, \code{s2} is a detail inside a card \\
\code{--ink1} to \code{--ink3}  & 9--11  & Text, brightest to dimmest. \code{ink1} is a heading, \code{ink2} is body text, \code{ink3} is a caption \\
\code{--accent}, \code{--accentH}, \code{--accent2} & 12--14 & The brand blue, its hover variant, and an orange \\
\code{--green}                  & 15     & The one success colour \\
\code{--border}, \code{--borderS} & 16--17 & A visible hairline, and a subtler one \\
\code{--font-d}, \code{--font-b} & 18--19 & The display typeface and the body typeface \\
\code{--r}, \code{--rsm}, \code{--ease}, \code{--maxw}, \code{--shadow} & 20--24 & Corner radius, small radius, motion curve, content width, drop shadow \\
\bottomrule
\end{tabular}
\caption{The token vocabulary. The numeric suffix always means the same thing:
0 is furthest back or brightest, and the number rises as you come forward or
dim down.}
\end{table}

\begin{keyidea}{Why a numbered scale beats named colours}
Calling a colour \code{--dark-navy} ties the name to how it looks. Calling it
\code{--s1} ties the name to what it is \emph{for}: the surface one step up from
the page. Redesign the site in a light theme and \code{--s1} still means the same
thing while \code{--dark-navy} becomes a lie. This is the single most transferable
idea in this repository, and it is why the author borrowed the scheme from
another project rather than inventing one.
\end{keyidea}

Line 18 defines the display typeface as Space Grotesk with three fallbacks after
it. A browser reads that list left to right and uses the first one it can find.
\code{ui-sans-serif} means the interface font of the operating system,
\code{system-ui} is the same idea by an older name, and \code{sans-serif} is the
last resort every browser guarantees. This is the chain that keeps the page
readable when the Google Fonts request fails.

Line 22 defines the motion curve. \code{cubic-bezier(.22,1,.36,1)} describes how
an animation's speed varies across its duration, by giving two control points of a
curve. This particular one starts fast and settles gently, the shape usually
described as ``ease out''. Every transition in the file refers to it, so all
movement on the page shares one personality.

\begin{honest}{Two tokens are defined and never used}
\code{--accent2} (line 14, an orange) and \code{--shadow} (line 24) are read zero
times in the entire file. Counted with \code{grep -c 'var(--accent2)'}, which
returns \code{0}, and likewise for \code{--shadow}.

The orange is not entirely absent from the page: it appears as the literal
\code{\#ff6719} in the second avatar gradient (line 320) and \code{\#ff9a5c} in
the second proof avatar (line 580). So the colour is used and the token for it is
not, which is the worst of both worlds. Either delete the token or use it.
\end{honest}

\paragraph{Lines 26 to 43: the reset, and the one utility class.}

\begin{lstlisting}[firstnumber=26,caption={\file{css/site.css}, lines 26--43}]

*, *::before, *::after { box-sizing: border-box; margin: 0; padding: 0; }
html { color-scheme: dark; scroll-behavior: smooth; }
body {
  font-family: var(--font-b);
  background: var(--s0);
  color: var(--ink1);
  -webkit-font-smoothing: antialiased;
  text-rendering: optimizeLegibility;
  overflow-x: hidden;
  line-height: 1.5;
}
img, video, svg { max-width: 100%; display: block; }
a { color: inherit; text-decoration: none; }
button { cursor: pointer; border: none; font-family: inherit; }

/* ─── UTILITY ─── */
.container { max-width: var(--maxw); margin: 0 auto; padding: 0 24px; }
\end{lstlisting}

Line 27 is a reset, and it is doing three separate jobs. The selector \code{*}
means every element; \code{*::before} and \code{*::after} extend that to the two
invisible pseudo-elements CSS lets you attach to anything.

\code{box-sizing: border-box} is the most consequential declaration in the file.
By default, if you give an element a width of 200 pixels and then 20 pixels of
padding, it occupies 240 pixels: the padding is added outside the stated width.
That default surprises everyone. \code{border-box} changes it so the stated width
is the final width and padding eats into it. Every layout calculation in this
file assumes it.

\code{margin: 0; padding: 0} removes the browser's built-in spacing on headings,
paragraphs and lists, so that all spacing in the design is deliberate rather than
inherited from a default that varies between browsers.

Line 28 does two things to the document root. \code{color-scheme: dark} tells the
browser this page is dark, so that scrollbars, text-selection highlights and form
controls are drawn in their dark variants rather than staying stubbornly white.
\code{scroll-behavior: smooth} makes every jump-link glide rather than snap.

\begin{honest}{The smooth scroll is implemented twice}
Line 28 already makes anchor links glide. \file{main.js} lines 70 to 77 then
intercept every anchor link and call \code{scrollIntoView(\{behavior: 'smooth'\})}
to do the same thing again. Either alone is sufficient. The JavaScript version is
the redundant one, and it is also the one that carries a latent bug, described in
section 7.3.
\end{honest}

Lines 29 to 37 set the page's defaults: the body typeface, the darkest surface as
the background, the brightest ink as the text, and a line height of 1.5, meaning
each line of text occupies one and a half times the font's size. That last number
is the most important typographic decision in the file and the one most often set
badly.

\code{-webkit-font-smoothing: antialiased} on line 33 thins text rendering on Mac
displays. It is a Safari and Chrome specific property with no standard equivalent,
and it is the conventional companion to light text on a dark background, which
otherwise looks heavier than intended.

\code{overflow-x: hidden} on line 35 forbids horizontal scrolling. Here it is
load-bearing rather than cosmetic: the hero tiles are positioned by percentage and
some of them extend past the right edge of the window, which would otherwise let
the visitor scroll sideways into empty space.

\begin{honest}{Hiding horizontal overflow treats a symptom}
\code{overflow-x: hidden} on the body is a common fix and a blunt one: it hides
the consequence of something overflowing rather than stopping the overflow. The
cost is that if a genuine layout bug ever pushes content off screen, it will be
invisible instead of obvious. Here the cause is known and intentional (the drifting
tiles), so the trade is defensible, but it is worth recognising it as a patch.
\end{honest}

Line 38 makes images, videos and SVGs never exceed their container's width, and
makes them block-level. That second part fixes a genuinely confusing default:
images are inline by default, which reserves space beneath them for the descenders
of text that is not there, producing a mysterious few pixels of gap under every
image.

Line 39 strips links of their blue colour and underline, making them inherit
whatever colour surrounds them. That is a deliberate design choice and an
accessibility trade: colour alone now distinguishes a link from text, which fails
the usual guidance. In this file every link is either a button or a footer item, so
the practical harm is small.

Line 40 styles buttons, and is dead in the strict sense that this page contains no
\code{<button>} element at all. It costs nothing and would come alive the moment
one was added.

Line 43 is the only utility class in the file: a centred column no wider than 1320
pixels with 24 pixels of breathing room on each side. \code{margin: 0 auto} is the
standard centring idiom, meaning no vertical margin and equal automatic margins
left and right. Every section on the page wraps its content in this.

\paragraph{Lines 45 to 75: the button system.}

\begin{lstlisting}[firstnumber=44,caption={\file{css/site.css}, lines 44--75}]

.btn {
  display: inline-flex;
  align-items: center;
  gap: 8px;
  padding: 13px 28px;
  border-radius: 100px;
  font-family: var(--font-d);
  font-size: 15px;
  font-weight: 600;
  letter-spacing: -0.01em;
  transition: transform 0.18s var(--ease), box-shadow 0.18s var(--ease),
              background 0.18s var(--ease);
}
.btn-primary {
  background: var(--accent);
  color: #fff;
}
.btn-primary:hover {
  background: var(--accentH);
  transform: translateY(-2px);
  box-shadow: 0 8px 24px rgba(25,103,255,.4);
}
.btn-ghost {
  background: transparent;
  color: var(--ink2);
  border: 1px solid var(--border);
}
.btn-ghost:hover {
  background: var(--s1);
  color: var(--ink1);
}
\end{lstlisting}

This is the two-class pattern the HTML uses everywhere: \code{.btn} holds
everything both buttons share, and \code{.btn-primary} or \code{.btn-ghost} adds
only what differs. Add a third button style and you write four lines, not thirty.

\code{display: inline-flex} on line 46 makes the button a flex container that sits
in the flow of text rather than on its own line. That matters because the primary
buttons contain both a word and an icon, and flex is what lines those up.

\begin{jargon}{Flexbox}
A layout system that arranges an element's children in a row or a column and
controls how they align and how leftover space is shared. Set
\code{display: flex} on a parent and its children become flex items.
\code{align-items: center} lines them up on their centres, which is how the icon
in a button sits level with the text rather than on the baseline.
\end{jargon}

\code{gap: 8px} on line 48 puts eight pixels between the word and the icon
without either of them needing a margin. Before \code{gap} existed this required
fiddly rules targeting all-but-the-first child, and its arrival is one of the
genuine improvements in modern CSS.

\code{border-radius: 100px} on line 50 makes a pill. Any radius larger than half
the element's height produces a perfect semicircular end, so 100 is simply a
number comfortably larger than needed.

Lines 55 and 56 declare the transition: three named properties, each over 0.18
seconds, each using the shared curve. Listing properties explicitly rather than
writing \code{transition: all} is the better practice, because \code{all} makes
the browser watch every property for changes and can animate things you never
intended.

Lines 62 to 66 are the hover state. \code{transform: translateY(-2px)} lifts the
button two pixels, and the shadow beneath it grows: together they read as the
button rising towards the cursor. The shadow colour \code{rgba(25,103,255,.4)} is
the brand blue at 40\% opacity, written as red, green and blue numbers plus an
alpha.

\begin{jargon}{rgba and opacity}
\code{rgba(25,103,255,.4)} names a colour by its red, green and blue amounts from
0 to 255, plus an opacity from 0 (invisible) to 1 (solid). It is the same blue as
\code{\#1967ff}, expressed in the form that allows transparency.
\end{jargon}

\begin{honest}{Buttons have hover states and no focus states}
Neither button defines \code{:focus} or \code{:focus-visible}. A visitor
navigating with a keyboard rather than a mouse therefore gets whatever outline the
browser draws by default, which on a dark background is often nearly invisible.
Since these are the page's only interactive elements, this is the most consequential
accessibility gap in the stylesheet, and it is three lines to fix.
\end{honest}

\paragraph{Lines 77 to 92: the two reveal mechanisms.}

\begin{lstlisting}[firstnumber=76,caption={\file{css/site.css}, lines 76--92}]

/* ─── REVEAL (scroll-driven, driven by JS adding .visible) ─── */
.reveal {
  opacity: 0;
  transform: translateY(28px);
  transition: opacity 0.7s var(--ease), transform 0.7s var(--ease);
}
.reveal.visible { opacity: 1; transform: none; }

/* stagger via CSS custom property set inline */
.stagger > * {
  opacity: 0;
  transform: translateY(20px);
  transition: opacity 0.6s var(--ease), transform 0.6s var(--ease);
  transition-delay: calc(var(--i, 0) * 90ms);
}
.stagger.visible > * { opacity: 1; transform: none; }
\end{lstlisting}

These sixteen lines are the entire animation contract between the stylesheet and
the JavaScript, and the comment on line 77 states it: the CSS defines both states
and the transition between them, and the JavaScript's only job is to add the class
\code{visible} at the right moment.

\code{.reveal} starts invisible and 28 pixels low. \code{.reveal.visible}, which
matches an element carrying \emph{both} classes, returns it to fully visible and
its natural position. Because a transition is declared, the browser animates
between the two over 0.7 seconds rather than snapping.

The stagger is the same idea with one addition. \code{.stagger > *} means every
direct child of a staggered container, and line 90 delays each child's transition
by \code{calc(var(--i, 0) * 90ms)}: the number the HTML wrote inline, multiplied by
90 milliseconds. Child 0 starts immediately, child 3 starts 270 milliseconds later.

\begin{jargon}{calc, and a custom property's fallback}
\code{calc()} performs arithmetic in CSS, and can mix units.
\code{var(--i, 0)} reads the custom property \code{--i} and falls back to
\code{0} if it was never set, which is what keeps this rule safe on any element
whose inline \code{--i} was forgotten: it simply gets no delay instead of breaking.
\end{jargon}

\begin{keyidea}{Why the animation lives in CSS and the trigger lives in JavaScript}
This split is the reason the page degrades well. If JavaScript never runs, the
class \code{visible} is never added, and every revealed element stays at opacity
zero, invisible forever. That would be catastrophic, and it is exactly why line
767 exists: the reduced-motion block at the end of the file forces
\code{.reveal} and \code{.stagger > *} back to full opacity. That block is doing
double duty as an accessibility feature and as a safety net.
\end{keyidea}

\paragraph{Lines 94 to 133: the navigation bar.}

\begin{lstlisting}[firstnumber=93,caption={\file{css/site.css}, lines 93--133}]

/* ─── NAV ─── */
.nav {
  position: fixed;
  top: 0; left: 0; right: 0;
  z-index: 100;
  padding: 18px 24px;
  display: flex;
  align-items: center;
  justify-content: space-between;
  transition: background 0.3s, border-color 0.3s, box-shadow 0.3s;
}
.nav.scrolled {
  background: rgba(5,7,13,0.88);
  border-bottom: 1px solid var(--border);
  box-shadow: 0 4px 24px rgba(0,0,0,.3);
  backdrop-filter: blur(16px);
  -webkit-backdrop-filter: blur(16px);
}
.nav-inner {
  max-width: var(--maxw);
  width: 100%;
  margin: 0 auto;
  display: flex;
  align-items: center;
  justify-content: space-between;
}
.nav-logo {
  display: flex;
  align-items: center;
  gap: 10px;
  font-family: var(--font-d);
  font-size: 18px;
  font-weight: 700;
  letter-spacing: -0.02em;
  color: var(--ink1);
}
.nav-logo .logo-ai {
  color: var(--accent);
}
.nav-logo svg { width: 28px; height: 28px; }
\end{lstlisting}

\code{position: fixed} on line 96 pins the bar to the window rather than to the
document, so it stays put while the page scrolls beneath it. Lines 97 pins it to
the top and stretches it across.

\code{z-index: 100} on line 98 controls stacking order: when elements overlap, the
higher number is drawn in front. 100 is arbitrary and generous, chosen so nothing
else on the page can accidentally cover the bar.

Lines 105 to 111 are the scrolled state, the class the JavaScript adds once the
page has moved more than 40 pixels. It gives the bar a background at 88\% opacity,
a bottom hairline, and a shadow.

\code{backdrop-filter: blur(16px)} on line 109 is the interesting one: it blurs
whatever is \emph{behind} the element, producing the frosted-glass effect that only
works because the background is slightly transparent. Line 110 repeats it with a
\code{-webkit-} prefix for Safari, which needed the prefixed form for years.

\begin{honest}{A doubled flex layout}
\code{.nav} sets \code{display: flex} with \code{justify-content: space-between}
(lines 100 to 102), and \code{.nav-inner} sets exactly the same three properties
again (lines 116 to 118). The outer one has only one child, so its
\code{space-between} has nothing to space. It is harmless and it is redundant, and
it makes the layout harder to reason about than it needs to be.
\end{honest}

Line 130 is a descendant selector: \code{.nav-logo .logo-ai}, with a space between
the two names, means ``an element with class \code{logo-ai} anywhere inside an
element with class \code{nav-logo}''. That space is meaningful syntax, and reading
it as one thing rather than two is a common beginner error.

\paragraph{Lines 135 to 177: the hero shell and the floating animation.}

\begin{lstlisting}[firstnumber=134,caption={\file{css/site.css}, lines 134--177}]

/* ─── HERO ─── */
.hero {
  position: relative;
  min-height: 100vh;
  display: flex;
  align-items: center;
  justify-content: center;
  overflow: hidden;
  padding-top: 80px;
}

/* animated tile backdrop */
.hero-bg {
  position: absolute;
  inset: 0;
  pointer-events: none;
}
.tile {
  position: absolute;
  width: 176px;
  background: var(--s1);
  border: 1px solid var(--border);
  border-radius: var(--r);
  padding: 14px;
  opacity: 0;
  animation: tileFloat 12s var(--ease) infinite;
}
.tile:nth-child(1) { top: 12%; left: 3%;   animation-delay:  0.0s; }
.tile:nth-child(2) { top: 18%; right: 4%;  animation-delay:  1.5s; }
.tile:nth-child(3) { top: 55%; left: 5%;   animation-delay:  3.0s; }
.tile:nth-child(4) { bottom: 18%; right: 3%; animation-delay: 4.5s; }
.tile:nth-child(5) { top: 42%; right: 26%; animation-delay:  6.0s; }
.tile:nth-child(6) { top:  8%; left: 32%;  animation-delay:  7.5s; }
.tile:nth-child(7) { bottom: 28%; left: 22%; animation-delay: 9.0s; }
.tile:nth-child(8) { top: 30%; left: 64%;  animation-delay: 10.5s; }

@keyframes tileFloat {
  0%   { opacity: 0; transform: translateY(18px) scale(0.94); }
  12%  { opacity: 0.8; transform: translateY(0) scale(1); }
  72%  { opacity: 0.8; transform: translateY(-10px) scale(1); }
  85%  { opacity: 0; transform: translateY(-18px) scale(0.94); }
  100% { opacity: 0; transform: translateY(-18px) scale(0.94); }
}
\end{lstlisting}

\code{position: relative} on line 137 looks like it does nothing, and it is
essential. It establishes this element as the reference point for any absolutely
positioned descendant. Without it, the tiles would position themselves against the
whole window instead of against the hero.

\begin{jargon}{Positioning: relative and absolute}
\code{position: absolute} takes an element out of the normal flow and places it at
coordinates you give. Those coordinates are measured from the nearest ancestor
that is itself positioned, which in practice means the nearest one with
\code{position: relative}. This pairing, a relative parent and absolute children,
is the standard way to place things inside a box.
\end{jargon}

\code{min-height: 100vh} on line 138 makes the hero at least as tall as the
window. \code{vh} means one percent of the viewport height, so \code{100vh} is the
full height.

\code{pointer-events: none} on line 149 makes the tile container transparent to the
mouse: clicks pass straight through to whatever is beneath. Without it, the tiles
would sit over the hero buttons and swallow clicks aimed at ``Book a call''. It is
a one-line fix for a bug that would otherwise be maddening to diagnose, since the
button would look perfectly normal and simply not respond.

Lines 152 to 161 define a tile: absolutely positioned, 176 pixels wide, using the
card surface, the standard border and the standard 16-pixel radius, and starting
at \code{opacity: 0}. It is invisible until the animation makes it visible.

Line 160 runs the animation: the named sequence \code{tileFloat}, over 12 seconds,
using the shared curve, repeating forever.

Lines 162 to 169 place the eight tiles individually. \code{:nth-child(1)} matches
the first child of its parent, and so on.

\begin{keyidea}{This is where the HTML order becomes load-bearing}
Nothing here mentions \code{tile-blue} or \code{tile-orange}. The tiles are
positioned purely by their \emph{number}, so tile 5 is wherever the fifth child
happens to be. Reorder the blocks in \file{index.html} and the layout of the hero
changes, with no change to this file. That coupling is invisible from either file
alone, and it is the kind of thing this document exists to make visible.
\end{keyidea}

Note also that the eight delays are 1.5 seconds apart, from 0.0 up to 10.5, across
a 12-second loop. That is a deliberate arrangement: eight evenly spaced entrances
inside one cycle, so at any moment several tiles are visible at different stages
and the backdrop never appears empty or synchronised.

Lines 171 to 177 define the animation itself.

\begin{jargon}{Keyframes}
An animation defined by describing what the element looks like at particular
points through its run, given as percentages of the total duration. The browser
fills in every moment in between. \code{0\%} is the start and \code{100\%} is the
end.
\end{jargon}

Read as a story: at the start the tile is invisible, 18 pixels low and slightly
small. By 12\% of the way through (about 1.4 seconds) it has risen into place at
full size and 80\% opacity. It holds there until 72\%, drifting up 10 pixels. By
85\% it has faded out completely while continuing to rise and shrink. The final
15\% is spent invisible, which is the gap before it begins again.

\paragraph{Lines 179 to 218: what is inside a tile.}

\begin{lstlisting}[firstnumber=178,caption={\file{css/site.css}, lines 178--218}]

.tile-thumb {
  border-radius: 10px;
  height: 90px;
  margin-bottom: 10px;
  display: flex;
  align-items: center;
  justify-content: center;
  position: relative;
  overflow: hidden;
}
.tile-thumb svg { width: 28px; height: 28px; color: rgba(255,255,255,.6); position: relative; z-index: 1; }
.tile-meta { display: flex; align-items: center; justify-content: space-between; gap: 6px; }
.tile-label { font-size: 11px; color: var(--ink3); font-weight: 500; }
.tile-metric {
  font-family: var(--font-d);
  font-size: 11px;
  font-weight: 700;
  color: var(--green);
  background: rgba(52,211,153,.1);
  border-radius: 20px;
  padding: 2px 7px;
}

/* tile gradient variants */
.tile-blue .tile-thumb  { background: linear-gradient(135deg,#0a1e50,#0d2c6e); }
.tile-orange .tile-thumb { background: linear-gradient(135deg,#3a1800,#5a2800); }
.tile-teal .tile-thumb  { background: linear-gradient(135deg,#003838,#005252); }
.tile-violet .tile-thumb { background: linear-gradient(135deg,#1a0a40,#2a1060); }

/* fake text bars inside "post" tiles */
.tile-lines { padding: 0; }
.tile-line {
  height: 7px;
  border-radius: 4px;
  background: var(--s2);
  margin-bottom: 6px;
}
.tile-line:nth-child(1) { width: 100%; }
.tile-line:nth-child(2) { width: 85%; }
.tile-line:nth-child(3) { width: 60%; }
\end{lstlisting}

Lines 179 to 188 shape the thumbnail area: a fixed 90-pixel-tall rounded box that
centres its contents both ways using flex.

Line 189 styles the icon inside it. Note \code{color:} rather than \code{fill:} or
\code{stroke:}: this works because the SVG in the HTML was written with
\code{stroke="currentColor"}, so setting a text colour here reaches into the
drawing. That is the pay-off for the attribute noted in section 7.1.

Line 190 lays out the caption row with \code{space-between}, which pushes the label
to the left edge and the metric to the right, the flex idiom for a two-ended row.

Lines 192 to 200 style the metric as a small green pill: green text on a green
background at 10\% opacity, which is the standard recipe for a tinted badge that
suits any surface beneath it.

Lines 203 to 206 are the four colour variants. Each is a descendant selector: the
gradient applies to a \code{.tile-thumb} \emph{inside} a \code{.tile-blue}, not to
the tile itself. \code{linear-gradient(135deg, A, B)} blends from colour A to
colour B along a line at 135 degrees, which runs from top-left to bottom-right.

Line 209 is the dead rule identified in section 7.1: both places that use
\code{.tile-lines} override its padding inline, so this never applies.

Lines 210 to 218 draw the fake text bars: 7 pixels tall, rounded, in the \code{s2}
surface, with the three widths that make them read as a paragraph. Note that tile
8 overrides the third width inline to 50\%, so these are defaults rather than
rules.

\paragraph{Lines 220 to 243: the statistic tiles and the readability overlay.}

\begin{lstlisting}[firstnumber=219,caption={\file{css/site.css}, lines 219--243}]

/* stats tile */
.tile-stat-num {
  font-family: var(--font-d);
  font-size: 26px;
  font-weight: 700;
  color: var(--ink1);
  line-height: 1;
  margin-bottom: 4px;
}
.tile-stat-lbl { font-size: 11px; color: var(--ink3); }

/* overlay - keeps center readable */
.hero-overlay {
  position: absolute;
  inset: 0;
  background:
    radial-gradient(ellipse 55% 65% at 50% 50%,
      rgba(5,7,13,.9) 0%,
      rgba(5,7,13,.55) 70%,
      transparent 100%),
    linear-gradient(to bottom,
      rgba(5,7,13,.2) 0%,
      rgba(5,7,13,.85) 100%);
}
\end{lstlisting}

\code{line-height: 1} on line 226 is deliberate. The page-wide default of 1.5 is
right for paragraphs and wrong for a large number, where it leaves a visible gap
above and below. Setting it to 1 makes the line exactly as tall as the digits.

Lines 231 to 243 are the overlay, and they are the most sophisticated piece of
styling in the file. Two gradients are stacked in one declaration, separated by a
comma, and in CSS the \emph{first} one listed is drawn on top.

The first is a radial gradient: an ellipse 55\% wide and 65\% tall, centred, going
from 90\% opaque near-black at its centre, through 55\% at seven-tenths of the way
out, to fully transparent at the edge. That is a dark cloud over the middle of the
hero, exactly where the headline sits.

The second is a linear gradient running down the whole area, from 20\% opaque at
the top to 85\% at the bottom, which darkens the lower half so the hero blends into
the section beneath it.

Together they solve the real problem with a busy animated backdrop: text over
moving content is hard to read, and the contrast changes as the content moves.
Rather than dimming the tiles themselves, the design puts a shaped veil between
them and the text.

\begin{jargon}{inset}
\code{inset: 0} is shorthand for \code{top: 0; right: 0; bottom: 0; left: 0},
which on an absolutely positioned element means ``stretch to fill the parent
completely''.
\end{jargon}

\paragraph{Lines 245 to 324: the hero's content.}

\begin{lstlisting}[firstnumber=244,caption={\file{css/site.css}, lines 244--275}]

.hero-content {
  position: relative;
  z-index: 2;
  text-align: center;
  max-width: 700px;
  padding: 0 24px;
}
.hero-badge {
  display: inline-flex;
  align-items: center;
  gap: 8px;
  background: rgba(25,103,255,.12);
  border: 1px solid rgba(25,103,255,.3);
  border-radius: 100px;
  padding: 6px 14px;
  font-size: 13px;
  font-weight: 600;
  color: var(--accentH);
  letter-spacing: 0.01em;
  margin-bottom: 28px;
}
.hero-badge-dot {
  width: 7px; height: 7px;
  border-radius: 50%;
  background: var(--accent);
  animation: pulse 2s ease-in-out infinite;
}
@keyframes pulse {
  0%, 100% { opacity: 1; transform: scale(1); }
  50%       { opacity: 0.5; transform: scale(0.7); }
}
\end{lstlisting}

\code{z-index: 2} on line 247 lifts the content above the tiles and the overlay,
which is what makes the headline readable rather than buried. Note that
\code{z-index} only works on positioned elements, which is why line 246 sets
\code{position: relative} even though nothing is being moved.

The badge (lines 252 to 265) uses the tinted-pill recipe again: brand blue at 12\%
for the fill and 30\% for the border, with the brighter hover blue for the text so
it stays legible against its own tint.

Lines 266 to 275 define the pulsing dot: a 7-pixel square made circular with a 50\%
radius, running a two-second pulse forever. The keyframes group \code{0\%, 100\%}
on one line, which is the idiom for an animation that returns to where it started:
full size at both ends, shrunk to 70\% and half-faded in the middle.

\begin{lstlisting}[firstnumber=276,caption={\file{css/site.css}, lines 276--324}]
.hero-h1 {
  font-family: var(--font-d);
  font-size: clamp(42px, 6.5vw, 80px);
  font-weight: 700;
  letter-spacing: -0.03em;
  line-height: 1.06;
  color: var(--ink1);
  margin-bottom: 36px;
}
.hero-h1 em {
  font-style: normal;
  color: var(--accent);
}
.hero-actions {
  display: flex;
  align-items: center;
  justify-content: center;
  gap: 14px;
  flex-wrap: wrap;
  margin-bottom: 48px;
}
.hero-social-proof {
  display: flex;
  align-items: center;
  gap: 12px;
  justify-content: center;
}
.avatars { display: flex; }
.avatar {
  width: 34px; height: 34px;
  border-radius: 50%;
  border: 2px solid var(--s0);
  margin-left: -10px;
  display: flex;
  align-items: center;
  justify-content: center;
  font-family: var(--font-d);
  font-size: 12px;
  font-weight: 700;
  color: #fff;
  flex-shrink: 0;
}
.avatars .avatar:first-child { margin-left: 0; }
.av1 { background: linear-gradient(135deg,#1967ff,#3f82ff); }
.av2 { background: linear-gradient(135deg,#ff6719,#ff8c4a); }
.av3 { background: linear-gradient(135deg,#34d399,#10b981); }
.av4 { background: linear-gradient(135deg,#a78bfa,#7c3aed); }
.hero-proof-text { font-size: 13px; color: var(--ink3); }
.hero-proof-text strong { color: var(--ink2); }
\end{lstlisting}

Line 278 is the most useful single line in the file to learn from.

\begin{jargon}{clamp}
\code{clamp(minimum, preferred, maximum)} picks the preferred value, but never lets
it fall below the minimum or rise above the maximum. Here the preferred value is
\code{6.5vw}, meaning 6.5\% of the window width, so the headline scales
continuously with the window between a floor of 42 pixels and a ceiling of 80. One
line replaces what used to need several media queries.
\end{jargon}

Line 280 sets \code{letter-spacing: -0.03em}, pulling the letters slightly closer
together. Large display type looks loose at its default spacing, and negative
tracking at large sizes is standard practice. The unit \code{em} means ``relative
to this element's font size'', so the tightening scales with the text.

Lines 285 to 288 are the rule flagged in section 7.1: it takes the \code{<em>} tag,
removes the italics it normally implies, and colours it blue instead.

Line 294 sets \code{flex-wrap: wrap} on the button row, so on a narrow screen the
second button drops onto its own line instead of overflowing. That single property
is why the hero needs no special mobile rule for its buttons.

Lines 303 to 318 build the overlapping avatar row, and the technique is worth
naming: each circle has a negative left margin of 10 pixels, pulling it over its
neighbour, and a 2-pixel border in the page background colour, which reads as a gap
punched out of the circle behind it. Line 318 then cancels the negative margin on
the first one so the row does not start shifted left.

\code{flex-shrink: 0} on line 316 stops the circles being squashed when space is
tight. Flex items shrink by default, and circles that shrink become ovals.

\begin{honest}{Four literal colours where tokens exist}
Lines 319 to 322 write \code{\#1967ff}, \code{\#3f82ff}, \code{\#ff6719} and
\code{\#34d399} directly. Three of those four have tokens (\code{--accent},
\code{--accentH}, \code{--green}) and the fourth has one that is otherwise unused
(\code{--accent2}). Gradients accept custom properties perfectly well, so there is
no technical reason for the literals. This is the same inconsistency as the inline
styles in the HTML, and together they mean the token system governs most of the
page rather than all of it.
\end{honest}

\paragraph{Lines 326 to 359: the metrics band.}

\begin{lstlisting}[firstnumber=325,caption={\file{css/site.css}, lines 325--359}]

/* ─── METRICS BAND ─── */
.metrics {
  padding: 80px 0;
  border-top: 1px solid var(--borderS);
  border-bottom: 1px solid var(--borderS);
}
.metrics-grid {
  display: grid;
  grid-template-columns: repeat(4, 1fr);
  gap: 2px;
}
.metric-item {
  padding: 36px 24px;
  text-align: center;
  border-right: 1px solid var(--borderS);
}
.metric-item:last-child { border-right: none; }
.metric-num {
  font-family: var(--font-d);
  font-size: clamp(36px, 4.5vw, 60px);
  font-weight: 700;
  letter-spacing: -0.03em;
  line-height: 1;
  color: var(--ink1);
  margin-bottom: 8px;
}
.metric-num .accent { color: var(--accent); }
.metric-label {
  font-size: 13px;
  color: var(--ink3);
  font-weight: 500;
  text-transform: uppercase;
  letter-spacing: 0.06em;
}
\end{lstlisting}

Line 333 introduces the second layout system in the file.

\begin{jargon}{CSS Grid}
A layout system for two dimensions, as opposed to flexbox's one.
\code{grid-template-columns: repeat(4, 1fr)} means four columns, each taking one
equal share of the space. \code{1fr} is a fraction of what is left over after any
fixed sizes are accounted for.
\end{jargon}

The separators here are drawn as a right border on every item (line 340), with the
last one removed (line 342). That is a different technique from the one used later
by the services grid, which draws its lines using gaps. Both work; using both in
one file is an inconsistency, and the border approach is the one that needed a
special case.

Line 352, \code{.metric-num .accent}, styles an element with class \code{accent}
inside a metric number. \textbf{No such element exists} anywhere in
\file{index.html}. This rule is dead.

Lines 353 to 359 style the label beneath each number: small, dim, uppercase, with
positive letter-spacing. Note the tracking is positive here and negative on the
big numbers above. That is not an inconsistency but a rule of typography: large
text needs tightening, small uppercase text needs opening up or it becomes a
smudge.

\paragraph{Lines 361 to 379: shared section headings.}

\begin{lstlisting}[firstnumber=360,caption={\file{css/site.css}, lines 360--379}]

/* ─── PIPELINE ─── */
.pipeline-section { padding: 120px 0; }
.section-kicker {
  font-size: 12px;
  font-weight: 700;
  text-transform: uppercase;
  letter-spacing: 0.1em;
  color: var(--accent);
  margin-bottom: 14px;
}
.section-title {
  font-family: var(--font-d);
  font-size: clamp(30px, 4vw, 52px);
  font-weight: 700;
  letter-spacing: -0.025em;
  line-height: 1.1;
  color: var(--ink1);
  margin-bottom: 72px;
}
\end{lstlisting}

Although these sit under the pipeline banner, \code{.section-kicker} and
\code{.section-title} are used by all four remaining sections. Their placement here
is a small organisational untidiness: a reader looking for the heading styles would
reasonably expect them near the top with the other shared pieces, not inside the
pipeline block.

The kicker is the small blue uppercase line (``How it works'', ``The work'') and the
title is the large heading beneath it. Both follow the same tracking rule as the
metrics band.

\begin{honest}{The kicker's contrast is worth measuring}
\code{--accent} is \code{\#1967ff}, a saturated blue, and the kicker sets it at 12
pixels in uppercase on the near-black page background. Small text needs more
contrast than large text, not less. This document does not have a measured
contrast ratio for that exact pairing and so does not assert a failure, but it is
the first thing to check in an accessibility audit, along with the 11-pixel service
tags which use the dimmest ink on the same background.
\end{honest}

\paragraph{Lines 381 to 450: the pipeline.}

\begin{lstlisting}[firstnumber=380,caption={\file{css/site.css}, lines 380--426}]

.pipeline {
  display: flex;
  align-items: center;
  gap: 0;
  position: relative;
  overflow: hidden;
}
.pipeline-node {
  flex: 1;
  display: flex;
  flex-direction: column;
  align-items: center;
  gap: 16px;
  position: relative;
  z-index: 1;
  opacity: 0;
  transform: scale(0.85);
  transition: opacity 0.5s var(--ease), transform 0.5s var(--ease);
}
.pipeline.visible .pipeline-node:nth-child(1) { opacity: 1; transform: none; transition-delay: 0.0s; }
.pipeline.visible .pipeline-node:nth-child(3) { opacity: 1; transform: none; transition-delay: 0.2s; }
.pipeline.visible .pipeline-node:nth-child(5) { opacity: 1; transform: none; transition-delay: 0.4s; }
.pipeline.visible .pipeline-node:nth-child(7) { opacity: 1; transform: none; transition-delay: 0.6s; }
.pipeline.visible .pipeline-node:nth-child(9) { opacity: 1; transform: none; transition-delay: 0.8s; }

.pipeline-connector {
  flex: 0 0 auto;
  width: 60px;
  height: 2px;
  background: var(--border);
  position: relative;
  overflow: hidden;
}
.pipeline-connector::after {
  content: '';
  position: absolute;
  top: 0; left: 0;
  height: 100%;
  width: 0;
  background: linear-gradient(to right, var(--accent), var(--accentH));
  transition: width 0.6s var(--ease);
}
.pipeline.visible .pipeline-connector:nth-child(2)::after { width: 100%; transition-delay: 0.15s; }
.pipeline.visible .pipeline-connector:nth-child(4)::after { width: 100%; transition-delay: 0.35s; }
.pipeline.visible .pipeline-connector:nth-child(6)::after { width: 100%; transition-delay: 0.55s; }
.pipeline.visible .pipeline-connector:nth-child(8)::after { width: 100%; transition-delay: 0.75s; }
\end{lstlisting}

This is the most intricate animation in the file and it is worth reading closely,
because it is doing something the stagger mechanism could not.

The pipeline is a flex row of nine children alternating node, connector, node,
connector and so on: five nodes at odd positions and four connectors at even ones.
That alternation is why the selectors count 1, 3, 5, 7, 9 for nodes and 2, 4, 6, 8
for connectors.

\code{flex: 1} on the nodes (line 389) tells each to take an equal share of the
leftover space. \code{flex: 0 0 auto} on the connectors (line 407) tells them the
opposite: do not grow, do not shrink, stay exactly 60 pixels. So the nodes spread
and the connectors stay fixed.

Each node begins invisible and at 85\% size, and the \code{.visible} rules bring it
to full size with an increasing delay, 0.2 seconds apart. The connectors do
something different: each has an \code{::after} pseudo-element, an invisible extra
box CSS creates on demand, absolutely positioned, starting at zero width and
filled with a blue gradient. When the section becomes visible its width transitions
to 100\%, so the blue appears to flow along the line.

\begin{jargon}{Pseudo-element}
An extra box that CSS creates inside an element without any corresponding tag in
the HTML, written as \code{::before} or \code{::after}. It requires the
\code{content} property to exist at all, which is why line 415 sets
\code{content: ''}: an empty string is enough to bring it into being. It is the
standard way to draw a decorative shape without polluting the markup.
\end{jargon}

The delays interleave deliberately: node 1 at 0.0, connector at 0.15, node 3 at
0.2, connector at 0.35, and so on. The line starts drawing before the next node
lands, which is what makes the sequence read as flowing rather than as a series of
separate appearances.

\begin{honest}{This animation cannot survive an edit to the HTML}
Every selector here depends on the exact child positions in \file{index.html}. Add
a sixth step, or wrap the nodes in a container, and every one of these nine rules
silently stops matching: the pipeline would simply never appear, staying at opacity
zero forever, with no error anywhere. Selecting by class
(\code{.pipeline-node:nth-of-type(n)}) would have been sturdier for the same
effect.
\end{honest}

\begin{lstlisting}[firstnumber=427,caption={\file{css/site.css}, lines 427--450}]

.pipeline-icon {
  width: 72px; height: 72px;
  border-radius: 20px;
  background: var(--s1);
  border: 1px solid var(--border);
  display: flex;
  align-items: center;
  justify-content: center;
  transition: border-color 0.3s, box-shadow 0.3s;
}
.pipeline-icon svg { width: 30px; height: 30px; color: var(--accent); }
.pipeline-node:hover .pipeline-icon {
  border-color: var(--accent);
  box-shadow: 0 0 24px rgba(25,103,255,.25);
}
.pipeline-lbl {
  font-family: var(--font-d);
  font-size: 13px;
  font-weight: 600;
  color: var(--ink2);
  text-align: center;
  white-space: nowrap;
}
\end{lstlisting}

Lines 439 to 442 give each icon box a blue border and a soft blue glow when the
mouse is over its node. A shadow with no offset and a large blur, as here, reads as
a glow rather than a shadow.

\code{white-space: nowrap} on line 449 stops labels such as ``Auto-Post'' wrapping
onto two lines and disturbing the row's alignment. It is overridden at line 739 for
narrow screens, where the pipeline stacks vertically and there is no alignment left
to protect.

\paragraph{Lines 452 to 539: the portfolio cards.}

\begin{lstlisting}[firstnumber=451,caption={\file{css/site.css}, lines 451--509}]

/* ─── PORTFOLIO ─── */
.portfolio-section { padding: 0 0 120px; }
.portfolio-grid {
  display: grid;
  grid-template-columns: repeat(3, 1fr);
  gap: 16px;
}
.port-card {
  background: var(--s1);
  border: 1px solid var(--border);
  border-radius: var(--r);
  overflow: hidden;
  transition: transform 0.22s var(--ease), border-color 0.22s, box-shadow 0.22s;
  cursor: pointer;
}
.port-card:hover {
  transform: translateY(-4px);
  border-color: rgba(25,103,255,.4);
  box-shadow: 0 16px 48px rgba(0,0,0,.4);
}
.port-thumb {
  aspect-ratio: 16/9;
  position: relative;
  display: flex;
  align-items: center;
  justify-content: center;
  overflow: hidden;
}
.port-thumb-overlay {
  position: absolute;
  inset: 0;
  background: rgba(5,7,13,.25);
  display: flex;
  align-items: center;
  justify-content: center;
  opacity: 0;
  transition: opacity 0.22s;
}
.port-card:hover .port-thumb-overlay { opacity: 1; }
.port-play {
  width: 52px; height: 52px;
  border-radius: 50%;
  background: rgba(255,255,255,.15);
  backdrop-filter: blur(8px);
  display: flex;
  align-items: center;
  justify-content: center;
  border: 1.5px solid rgba(255,255,255,.3);
}
.port-play svg { width: 22px; height: 22px; color: #fff; margin-left: 2px; }

/* gradient thumbnails - replace with real video posters */
.pt1 { background: linear-gradient(135deg,#0a1840 0%,#0d2c6e 50%,#0a1840 100%); }
.pt2 { background: linear-gradient(135deg,#2a0e00 0%,#5a2800 50%,#2a0e00 100%); }
.pt3 { background: linear-gradient(135deg,#001a28 0%,#003855 50%,#001a28 100%); }
.pt4 { background: linear-gradient(135deg,#150050 0%,#2d0080 50%,#150050 100%); }
.pt5 { background: linear-gradient(135deg,#0a2800 0%,#1a4800 50%,#0a2800 100%); }
.pt6 { background: linear-gradient(135deg,#280014 0%,#500028 50%,#280014 100%); }
\end{lstlisting}

\code{overflow: hidden} on line 463 is what makes the rounded corners work: without
it, the gradient thumbnail inside would square off the top two corners it is
supposed to follow.

\code{cursor: pointer} on line 465 changes the mouse to a pointing hand, which
conventionally means clickable. As section 7.1 noted, the card is not clickable,
so this is a promise the page does not keep.

\code{aspect-ratio: 16/9} on line 473 reserves a widescreen-shaped area regardless
of width. This is worth appreciating: it means the layout does not shift when a
real image eventually loads, because the space is already the right shape. It is
the modern replacement for a well-known hack involving percentage padding.

The overlay (lines 480 to 489) sits over the whole thumbnail at zero opacity and
fades to one on hover, revealing the play button beneath a 25\% black wash. The play
button itself (lines 491 to 500) is a frosted circle: white at 15\% opacity with an
8-pixel backdrop blur and a brighter rim.

Line 501 gives the play triangle \code{margin-left: 2px}. That is not an accident.
A triangle's visual centre of mass sits left of its geometric centre, so a
perfectly centred play icon looks slightly left. Nudging it two pixels right is a
standard correction, and noticing it is a good sign that whoever built this cared
about detail.

Lines 504 to 509 are the six placeholder gradients, each a three-stop blend that
returns to its starting colour, so it reads as a soft sheen rather than a slope.
The comment on line 503 says outright that they are standing in for real video
posters that never arrived.

\begin{lstlisting}[firstnumber=510,caption={\file{css/site.css}, lines 510--539}]

.port-badge {
  position: absolute;
  top: 10px; right: 10px;
  background: rgba(5,7,13,.75);
  backdrop-filter: blur(8px);
  border: 1px solid var(--border);
  border-radius: 20px;
  padding: 4px 10px;
  font-size: 11px;
  font-weight: 600;
  color: var(--green);
}
.port-info { padding: 16px; }
.port-platform {
  display: flex;
  align-items: center;
  gap: 6px;
  font-size: 11px;
  color: var(--ink3);
  margin-bottom: 6px;
}
.port-platform svg { width: 13px; height: 13px; }
.port-title {
  font-family: var(--font-d);
  font-size: 15px;
  font-weight: 600;
  color: var(--ink1);
  letter-spacing: -0.01em;
}
\end{lstlisting}

The badge is pinned to the thumbnail's top-right corner, which works because
\code{.port-thumb} set \code{position: relative} on line 474.

The card title on lines 533 to 539 is the only place in the file where the display
typeface is used at a small size (15 pixels). Everywhere else Space Grotesk
appears it is either large (a heading, a metric) or tiny and uppercase (a label),
so this is the one instance of it doing ordinary work.

\paragraph{Lines 541 to 604: the client proof cards.}

\begin{lstlisting}[firstnumber=540,caption={\file{css/site.css}, lines 540--604}]

/* ─── CLIENT PROOF ─── */
.proof-section { padding: 0 0 120px; }
.proof-grid {
  display: grid;
  grid-template-columns: repeat(3, 1fr);
  gap: 16px;
}
.proof-card {
  background: var(--s1);
  border: 1px solid var(--border);
  border-radius: var(--r);
  padding: 28px;
  display: flex;
  flex-direction: column;
  gap: 20px;
  transition: border-color 0.22s, transform 0.22s var(--ease);
}
.proof-card:hover {
  border-color: rgba(25,103,255,.3);
  transform: translateY(-3px);
}
.proof-header {
  display: flex;
  align-items: center;
  gap: 14px;
}
.proof-avatar {
  width: 52px; height: 52px;
  border-radius: 50%;
  display: flex;
  align-items: center;
  justify-content: center;
  font-family: var(--font-d);
  font-size: 20px;
  font-weight: 700;
  color: #fff;
  flex-shrink: 0;
}
.pa1 { background: linear-gradient(135deg,#1967ff,#3f82ff); }
.pa2 { background: linear-gradient(135deg,#ff6719,#ff9a5c); }
.pa3 { background: linear-gradient(135deg,#7c3aed,#a78bfa); }

.proof-name {
  font-family: var(--font-d);
  font-size: 16px;
  font-weight: 700;
  color: var(--ink1);
  letter-spacing: -0.01em;
}
.proof-company { font-size: 13px; color: var(--ink3); margin-top: 2px; }
.proof-metric {
  display: flex;
  align-items: baseline;
  gap: 6px;
}
.proof-num {
  font-family: var(--font-d);
  font-size: 36px;
  font-weight: 700;
  color: var(--accent);
  letter-spacing: -0.03em;
  line-height: 1;
}
.proof-unit { font-size: 14px; color: var(--ink3); }
\end{lstlisting}

The proof card is a vertical flex column with a 20-pixel gap, so its two halves
(the header and the metric) space themselves without margins.

\code{align-items: baseline} on line 594 is the detail worth learning. The metric
row holds a 36-pixel number beside 14-pixel text. Centring them would look wrong,
because the eye reads the two as sitting on a shared line of writing.
\code{baseline} aligns them on that invisible line the letters sit on, which is
what a typesetter would do.

The three avatar gradients (lines 579 to 581) repeat the same literal-colour issue
as the hero avatars, and \code{pa3} introduces two purples that exist nowhere in
the token block at all.

\paragraph{Lines 606 to 654: the services grid.}

\begin{lstlisting}[firstnumber=605,caption={\file{css/site.css}, lines 605--654}]

/* ─── SERVICES ─── */
.services-section { padding: 0 0 120px; }
.services-grid {
  display: grid;
  grid-template-columns: repeat(3, 1fr);
  gap: 2px;
  background: var(--border);
  border-radius: var(--r);
  overflow: hidden;
}
.service-item {
  background: var(--s0);
  padding: 40px 32px;
  display: flex;
  flex-direction: column;
  gap: 16px;
  transition: background 0.22s;
}
.service-item:hover { background: var(--s1); }
.service-icon {
  width: 52px; height: 52px;
  border-radius: var(--rsm);
  background: rgba(25,103,255,.1);
  border: 1px solid rgba(25,103,255,.2);
  display: flex;
  align-items: center;
  justify-content: center;
  transition: background 0.22s, border-color 0.22s;
}
.service-icon svg { width: 22px; height: 22px; color: var(--accentH); }
.service-item:hover .service-icon {
  background: rgba(25,103,255,.18);
  border-color: rgba(25,103,255,.45);
}
.service-name {
  font-family: var(--font-d);
  font-size: 17px;
  font-weight: 700;
  color: var(--ink1);
  letter-spacing: -0.01em;
}
.service-tag {
  display: inline-block;
  font-size: 11px;
  font-weight: 600;
  color: var(--ink3);
  text-transform: uppercase;
  letter-spacing: 0.07em;
}
\end{lstlisting}

Lines 608 to 615 use the technique mentioned in section 7.1 and it deserves its own
explanation, because it is genuinely clever and genuinely fragile.

\begin{keyidea}{Hairlines made of gaps}
The grid is given a 2-pixel gap and a background of the border colour. Each cell is
then given the page background. Because the cells cover the grid's background
everywhere except in the gaps, the gaps show through as 2-pixel lines. The result
is a perfect internal grid of hairlines with no border rules and no special case
for the last item in a row, which is exactly what the metrics band above needed
(\code{:last-child \{ border-right: none \}}).

The fragility is that it depends on the cell background matching the page. Put this
grid on a different surface without changing \code{.service-item}, and the cells
become visible rectangles of the wrong colour.
\end{keyidea}

\code{border-radius} plus \code{overflow: hidden} on lines 613 and 614 round the
whole grid and clip the corner cells to match.

Line 647 sets \code{display: inline-block} on the tag. Since the tag is inside a
flex column, it would already behave as a block, so this has no visible effect: it
is a harmless leftover.

\paragraph{Lines 656 to 725: the booking panel and the footer.}

\begin{lstlisting}[firstnumber=655,caption={\file{css/site.css}, lines 655--725}]

/* ─── BOOK A CALL ─── */
.book-section { padding: 0 0 140px; }
.book-wrap {
  background: var(--s1);
  border: 1px solid var(--border);
  border-radius: 24px;
  padding: 72px 48px;
  text-align: center;
  position: relative;
  overflow: hidden;
}
.book-wrap::before {
  content: '';
  position: absolute;
  top: -100px; left: 50%;
  transform: translateX(-50%);
  width: 600px; height: 300px;
  background: radial-gradient(ellipse, rgba(25,103,255,.12) 0%, transparent 70%);
  pointer-events: none;
}
.book-eyebrow {
  font-size: 12px;
  font-weight: 700;
  text-transform: uppercase;
  letter-spacing: 0.1em;
  color: var(--accent);
  margin-bottom: 16px;
}
.book-title {
  font-family: var(--font-d);
  font-size: clamp(28px, 4vw, 48px);
  font-weight: 700;
  letter-spacing: -0.025em;
  color: var(--ink1);
  margin-bottom: 12px;
}
.book-sub {
  font-size: 16px;
  color: var(--ink3);
  margin-bottom: 36px;
}
.book-cta-row {
  display: flex;
  align-items: center;
  justify-content: center;
  gap: 14px;
  flex-wrap: wrap;
}
.book-detail {
  display: flex;
  align-items: center;
  justify-content: center;
  gap: 6px;
  font-size: 13px;
  color: var(--ink3);
  margin-top: 20px;
}
.book-detail svg { width: 15px; height: 15px; }
.book-detail span { color: var(--ink2); }

/* ─── FOOTER ─── */
footer {
  border-top: 1px solid var(--borderS);
  padding: 32px 24px;
  text-align: center;
  font-size: 13px;
  color: var(--ink3);
}
footer a { color: var(--ink3); transition: color 0.15s; }
footer a:hover { color: var(--ink1); }
\end{lstlisting}

Lines 667 to 675 add a blue glow above the booking panel using an \code{::before}
pseudo-element: a 600 by 300 ellipse of brand blue at 12\% opacity, positioned 100
pixels above the panel's top edge and centred with the standard
\code{left: 50\%} plus \code{translateX(-50\%)} idiom. Because the panel has
\code{overflow: hidden}, only the bottom part of the glow is visible, which is what
makes it read as light spilling in from above rather than as a floating oval.

\code{pointer-events: none} on line 674 appears for the second time in the file, for
the same reason as before: this decorative layer sits over the panel and must not
intercept the click aimed at the booking button.

The footer is deliberately unremarkable: one hairline, centred dim text, and links
that brighten on hover over 0.15 seconds.

\paragraph{Lines 727 to 768: the responsive rules.}

\begin{lstlisting}[firstnumber=726,caption={\file{css/site.css}, lines 726--768}]

/* ─── RESPONSIVE ─── */
@media (max-width: 900px) {
  .metrics-grid { grid-template-columns: repeat(2, 1fr); }
  .metric-item { border-right: none; border-bottom: 1px solid var(--borderS); }
  .metric-item:nth-child(odd) { border-right: 1px solid var(--borderS); }
  .metric-item:nth-child(3),
  .metric-item:nth-child(4) { border-bottom: none; }

  .pipeline { flex-direction: column; }
  .pipeline-connector { width: 2px; height: 40px; }
  .pipeline-connector::after { width: 100% !important; height: 0; }
  .pipeline.visible .pipeline-connector::after { height: 100%; width: 100% !important; }
  .pipeline-lbl { white-space: normal; }

  .portfolio-grid { grid-template-columns: repeat(2, 1fr); }
  .proof-grid     { grid-template-columns: 1fr; }
  .services-grid  { grid-template-columns: repeat(2, 1fr); }

  .tile { display: none; }
  .tile:nth-child(1),
  .tile:nth-child(2),
  .tile:nth-child(3) { display: block; }
  .tile:nth-child(1) { top: 10%; left: 2%; width: 140px; }
  .tile:nth-child(2) { top: 15%; right: 2%; width: 140px; }
  .tile:nth-child(3) { display: none; }
}

@media (max-width: 600px) {
  .portfolio-grid { grid-template-columns: 1fr; }
  .services-grid  { grid-template-columns: 1fr; }
  .book-wrap { padding: 48px 24px; }
  .hero-h1 { font-size: 38px; }
  .tile { display: none; }
}

@media (prefers-reduced-motion: reduce) {
  *, *::before, *::after {
    animation-duration: .001ms !important;
    transition-duration: .001ms !important;
  }
  .reveal, .stagger > * { opacity: 1; transform: none; }
}
\end{lstlisting}

\begin{jargon}{Media query}
A block of CSS that only applies when a condition about the display is true.
\code{@media (max-width: 900px)} means ``apply everything in here only when the
window is 900 pixels wide or narrower''. Because it comes later in the file, its
rules override the equivalents above.
\end{jargon}

The 900-pixel block does five things. The metrics band folds from four columns to
two, and its dividing lines are rebuilt for the new shape: the right border is
removed from everything, a bottom border added to everything, the right border put
back on odd-numbered items so there is still a line down the middle, and the bottom
border removed from items 3 and 4 so the block does not end in a stray line. That
is four rules to reposition two lines, and it is the price of the border technique
that the services grid avoided.

The pipeline turns vertical (line 735) and its connectors become 2 pixels wide and
40 tall (line 736). Lines 737 and 738 then rebuild the fill animation to run
downwards instead of across, growing height rather than width, and use
\code{!important} to force the width value past the more specific rules further up
the file.

\begin{jargon}{!important}
A marker that forces a declaration to win against the normal precedence rules. It
is generally considered a last resort, because it breaks the predictability that
makes CSS manageable. It appears exactly twice in this file, both on line 737 and
738, both to override the horizontal fill animation.
\end{jargon}

The three content grids each drop a column, and the hero tiles are reduced from
eight to two.

\begin{honest}{Three lines of dead responsive code}
Line 745 hides every tile. Lines 746 to 748 then show tiles 1, 2 and 3. Line 751
hides tile 3 again. Because line 751 comes later and has identical specificity, it
wins, so tile 3 is hidden and the \code{display: block} it was given three lines
earlier does nothing at all.

The author evidently decided three tiles were too many for a tablet and removed the
third by adding a rule instead of editing the existing one. It is harmless, and it
is exactly the sort of thing that makes a reader doubt their understanding of the
cascade. It should be a two-line rule listing tiles 1 and 2.
\end{honest}

The 600-pixel block hides the tiles entirely, collapses the last two grids to one
column, tightens the booking panel's padding, and fixes the headline at 38 pixels,
overriding the \code{clamp} whose floor was 42.

The final block, lines 762 to 768, is the accessibility one, and it is the best
thing in the file.

\begin{keyidea}{What \code{prefers-reduced-motion} is and why this block matters}
Operating systems let people ask for reduced motion, usually because animation
causes them nausea, dizziness or migraine. The browser passes that preference to
the page, and this block responds to it.

It does two things. It crushes every animation and transition on the page to a
thousandth of a millisecond, which is effectively instant rather than disabled, a
deliberate choice because setting them to zero can stop some browsers firing the
events that mark an animation as finished. Then, crucially, line 767 forces every
revealed and staggered element to full opacity, because those elements start
invisible and would otherwise never appear once their transitions were neutered.

That second line is what turns a well-meaning block into a correct one, and it is
the detail most implementations of this pattern get wrong.
\end{keyidea}

\begin{honest}{The reduced-motion block cannot reach the counting numbers}
This block governs CSS only. The four metric numbers are animated by JavaScript,
counting up over 2.2 seconds, and \file{main.js} never checks this preference. A
visitor who has asked their system for reduced motion still gets four numbers
whirring upward.

The fix is four lines in \file{main.js}, given in section 7.3. This is the single
clearest defect in the project: not because it is severe, but because the intent to
handle it correctly is demonstrably there in the CSS and simply was not carried
across to the other file.
\end{honest}

\subsubsection*{What it connects to}

\textbf{Named by:} \file{new-design-reference/index.html}, line 16, and nothing
else.

\textbf{Names:} no other file. It refers to no images, no fonts by file (it names
font \emph{families}, which the HTML is responsible for downloading), and imports
nothing.

\textbf{Coupled to \file{index.html}} in three ways that no tool checks: by class
name throughout, by child position for the hero tiles (lines 162 to 169) and the
pipeline (lines 400 to 426), and by the shared contract that JavaScript adds the
class \code{visible}.

\textbf{Coupled to \file{js/main.js}} through exactly two class names,
\code{scrolled} and \code{visible}, both of which the JavaScript adds and this file
defines.

\subsubsection*{Concepts introduced}

This file is where most of the page's technical vocabulary first appears: custom
properties, the cascade, flexbox, grid, pseudo-elements, pseudo-classes, keyframe
animations, transitions, media queries, \code{clamp}, \code{calc},
\code{aspect-ratio}, \code{backdrop-filter}, stacking with \code{z-index}, and
\code{prefers-reduced-motion}. Each is defined at its first use above and collected
in the glossary.

\subsubsection*{How it breaks}

\begin{itemize}
\item \textbf{Delete this file and the page still works}, in the sense that every
  word remains readable. It simply looks like an unstyled document, and the hero
  tiles appear as a plain stacked list.
\item \textbf{Rename a class here without renaming it in the HTML} and that
  element silently loses its styling. Nothing reports it.
\item \textbf{Add or remove a hero tile or a pipeline step} and the position rules
  stop lining up, silently.
\item \textbf{Remove \code{pointer-events: none} on line 149} and the hero tiles
  begin swallowing clicks meant for the buttons beneath them.
\item \textbf{Remove \code{position: relative} on line 137} and all eight tiles
  jump to positions measured against the whole window instead of the hero.
\item \textbf{A missing semicolon} makes the browser discard the rest of that
  declaration block. Unlike a program, CSS never reports an error: it drops what it
  cannot parse and carries on, which is why a stylesheet bug shows up as something
  looking wrong rather than as a message.
\end{itemize}

\subsubsection*{Honest findings for this file}

Two unused tokens (lines 14 and 24); one dead rule with no matching element (line
352); one rule permanently overridden by inline styles (line 209); three lines of
dead responsive code (lines 746 to 748 versus 751); a doubled flex layout in the
navigation (lines 100 to 102 and 116 to 118); smooth scrolling implemented twice,
here and in JavaScript; literal colours used where tokens exist (lines 319 to 322,
579 to 581); two different techniques for drawing dividing lines in one file; no
focus states on either button; and a reduced-motion block that is exemplary within
CSS and cannot reach the one animation that lives outside it.

% ===========================================================================
\subsection{\file{new-design-reference/js/main.js}}
% ===========================================================================

\textbf{JavaScript, 79 lines.} Everything on the page that reacts. It does three
things and nothing else: it makes the navigation bar grow a frosted background
once you scroll, it counts the four metric numbers up from zero when they come
into view, and it fades most blocks upward as they arrive.

\begin{provenance}{new-design-reference/js/main.js}
Hand-written, and its own header comment says as much on line 3: \code{Zero
external dependencies.} That claim is true and checkable: the file imports
nothing, requires nothing, and the repository has no \file{package.json} that
could supply a library.

It arrived complete in the first commit, \code{0d77d51}, and has never been
edited. The patterns in it (an immediately invoked function wrapper, an
\code{IntersectionObserver} factory, a cubic ease-out counter) are the standard
vocabulary of vanilla JavaScript rather than the fingerprint of any particular
source, so no upstream is claimed. \textbf{Confidence: certain for hand-written
and dependency-free; no provenance evidence found for any specific inspiration,
and none is invented here.}
\end{provenance}

\subsubsection*{Why it exists}

Without this file the page is complete, readable and static. Two of its three
jobs are pure decoration and would not be missed. The third is not: the metric
numbers are written into the HTML as the character \code{0}, with their real
values hidden in data attributes, so with this file gone the page permanently
advertises zero projects delivered and zero views.

\begin{keyidea}{The single most important line in this file is in another file}
\file{index.html} loads this script at line 595, immediately before
\code{</body>}. Everything below runs the instant it loads, with no wrapper
telling it to wait for the document. That only works because, by the time the
browser reaches line 595, every element already exists. Move the
\code{<script>} tag into the \code{<head>} and this file would run against an
empty page, find nothing, attach nothing, and report no error at all.
\end{keyidea}

\subsubsection*{Line by line}

\paragraph{Lines 1 to 6: the header, and the wrapper.}

\begin{lstlisting}[language=Java,firstnumber=1,caption={\file{js/main.js}, lines 1--6}]
/* ai.digitalise.agency · main.js
   Scroll-driven animations: nav gloss, counter triggers, section reveals.
   Zero external dependencies. */

(function () {
  'use strict';
\end{lstlisting}

Lines 1 to 3 are a comment, using the same \code{/* */} form as CSS. It is an
accurate summary of the file, which is worth noting because a comment that
describes what the code actually does is less common than it should be.

Line 5 opens a wrapper worth understanding, because it looks like noise and is
not.

\begin{jargon}{Immediately invoked function expression}
Writing \code{(function () \{ ... \})();} defines a function and runs it once,
straight away, without ever giving it a name. The point is not the running: it is
the boundary. Anything declared inside is invisible outside.

Without it, the names \code{nav}, \code{animateCounter} and \code{observe}
declared below would become global, meaning any other script on the page could
see them, and any other script that happened to use the same name would collide
with them. One stray global called \code{observe} is exactly the kind of bug that
takes an afternoon to find.
\end{jargon}

Line 6 is \code{'use strict';}, which switches the language into a stricter mode
for everything inside this wrapper.

\begin{jargon}{Strict mode}
A setting that makes JavaScript refuse things it would otherwise silently allow.
The most valuable one: without it, assigning to a name you never declared quietly
creates a global variable, so a typo such as \code{navv = 3} succeeds and does
nothing useful. In strict mode it throws an error and you find the typo
immediately.
\end{jargon}

\paragraph{Lines 8 to 16: the navigation bar.}

\begin{lstlisting}[language=Java,firstnumber=7,caption={\file{js/main.js}, lines 7--16}]

  /* ── Nav: add .scrolled glass on scroll ── */
  const nav = document.querySelector('.nav');
  if (nav) {
    const onScroll = () => {
      nav.classList.toggle('scrolled', window.scrollY > 40);
    };
    window.addEventListener('scroll', onScroll, { passive: true });
    onScroll();
  }
\end{lstlisting}

Line 9 finds the navigation bar. \code{document.querySelector} takes a CSS
selector, exactly the same syntax the stylesheet uses, and returns the first
matching element, or the special value \code{null} if there is none.

\begin{jargon}{const, and why not var}
\code{const} declares a name whose binding cannot later be pointed at something
else. It is the default choice in modern JavaScript, with \code{let} used when a
name genuinely must change. The older \code{var} has confusing scoping rules and
appears nowhere in this file, which is a small sign of care.
\end{jargon}

Line 10 is a guard: everything only happens \code{if (nav)}, that is, only if
something was actually found. This is the pattern that makes the whole file
tolerant of a changing page. Delete the navigation bar from the HTML and this
block is skipped rather than crashing.

\begin{jargon}{Truthy and falsy}
JavaScript treats certain values as false when tested in a condition:
\code{false}, \code{0}, an empty string, \code{null} and \code{undefined}.
Everything else is treated as true. So \code{if (nav)} reads as ``if an element
was found'', because \code{querySelector} returns \code{null} when it finds
nothing.
\end{jargon}

Lines 11 to 13 define what to do when the page scrolls, using an arrow function.

\begin{jargon}{Arrow function}
\code{() => \{ ... \}} is a compact way to write a function. Here the compactness
is the only benefit; the older \code{function () \{ ... \}} would behave
identically in this file.
\end{jargon}

Line 12 is the actual work, and it is one line doing the job of five.
\code{classList.toggle('scrolled', condition)} adds the class when the condition
is true and removes it when false. The condition is \code{window.scrollY > 40},
meaning the page has been scrolled more than 40 pixels from the top. The
stylesheet at lines 105 to 111 defines what \code{scrolled} looks like, and the
transition on line 103 makes the change fade rather than snap.

Line 14 registers the function to run on every scroll event. \code{\{ passive:
true \}} is a performance promise to the browser: this handler will never cancel
the scroll. That matters because without the promise, the browser must wait for
the handler to finish before it knows whether the scroll may proceed, which can
make scrolling feel sticky. It is the correct flag here and it is often forgotten.

Line 15 calls the function once immediately. This handles the case where the page
is loaded already scrolled, for example when a visitor refreshes halfway down or
arrives on a link ending in \code{\#book}. Without it the bar would be transparent
over content until the first scroll.

\paragraph{Lines 18 to 36: the counter.}

\begin{lstlisting}[language=Java,firstnumber=17,caption={\file{js/main.js}, lines 17--36}]

  /* ── Counter animation (cubic ease-out) ── */
  function animateCounter(el) {
    const raw    = el.dataset.target || '0';
    const suffix = el.dataset.suffix || '';
    const target = parseFloat(raw.replace(/[^0-9.]/g, ''));
    const isFloat = raw.includes('.');
    const dec     = isFloat ? (raw.split('.')[1] || '').length : 0;
    const dur     = 2200;
    const start   = performance.now();

    const step = (now) => {
      const p      = Math.min((now - start) / dur, 1);
      const eased  = 1 - Math.pow(1 - p, 3);
      const val    = eased * target;
      el.textContent = (isFloat ? val.toFixed(dec) : Math.round(val).toLocaleString()) + suffix;
      if (p < 1) requestAnimationFrame(step);
    };
    requestAnimationFrame(step);
  }
\end{lstlisting}

This is the most substantial piece of logic in the repository, and it is worth
taking slowly.

Line 20 reads the target. \code{el.dataset.target} is how JavaScript reads the
attribute \code{data-target} written in the HTML: strip the \code{data-} prefix
and what remains is a property of \code{dataset}. The \code{|| '0'} at the end
supplies a default: if the attribute is missing, \code{dataset.target} is
\code{undefined}, which is falsy, so the whole expression evaluates to
\code{'0'}.

Line 21 does the same for the suffix, defaulting to an empty string.

Line 22 turns the text into a number. \code{raw.replace(/[\^{}0-9.]/g, '')}
deletes every character that is not a digit or a dot, and \code{parseFloat} then
reads what is left as a number.

\begin{jargon}{Regular expression}
A pattern for matching text, written between slashes. \code{[0-9.]} means ``a
digit or a dot''; putting a caret first negates it, so it means ``anything that
is not a digit or a dot''. The \code{g} after the closing slash means global:
replace every occurrence, not just the first.
\end{jargon}

Line 23 asks whether the original text contained a dot, which is how the code
distinguishes \code{1.6} from \code{17}. Line 24 counts the digits after the dot,
so that \code{1.6} is displayed to one decimal place throughout the animation
rather than as \code{1.5999999}. \code{(raw.split('.')[1] || '')} takes everything
after the dot, falling back to an empty string if there is nothing there, and
\code{.length} counts it.

\begin{jargon}{Ternary operator}
\code{condition ? a : b} evaluates to \code{a} when the condition is true and
\code{b} otherwise. It is an \code{if} statement compressed into an expression,
and it appears twice in this file, on lines 24 and 32.
\end{jargon}

Line 25 fixes the duration at 2200 milliseconds, 2.2 seconds. Line 26 records the
start time using \code{performance.now()}, a clock designed for measuring
durations precisely.

Lines 28 to 34 define one frame of the animation, and this is the heart of it.

\code{const p = Math.min((now - start) / dur, 1)} computes progress: how far
through the animation we are, from 0 at the start to 1 at the end. The
\code{Math.min(..., 1)} caps it, so that if the browser is busy and a frame
arrives late, progress never exceeds 1 and the number never overshoots its target.
That cap is the difference between a counter that lands on 17 and one that
occasionally flashes 19 before settling.

\code{const eased = 1 - Math.pow(1 - p, 3)} is the easing. Linear progress feels
mechanical; this curve makes the count start fast and slow dramatically as it
approaches the target, which is what makes it feel like it is settling.

\begin{keyidea}{What the easing formula actually does}
Take progress \code{p}, subtract it from 1, cube the result, and subtract that
from 1. At \code{p = 0.5}, halfway through in time, the formula gives
$1 - 0.5^3 = 0.875$: the counter is already 87.5\% of the way to its target. The
whole second half of the animation covers the last eighth of the distance. That
asymmetry is the entire effect.
\end{keyidea}

Line 32 renders the number. If the target had a decimal point, \code{toFixed(dec)}
formats it to that many places. Otherwise \code{Math.round} makes it a whole
number and \code{toLocaleString()} inserts the thousands separators appropriate to
the visitor's region, so a large number reads as \code{1,200} for a British
visitor and \code{1.200} for a German one. The suffix is then glued on the end.

Line 33 asks for the next frame if the animation is not finished, and line 35
starts the first one.

\begin{jargon}{requestAnimationFrame}
A request to run a function just before the browser's next screen repaint, about
60 times a second. Calling it repeatedly from inside itself, as here, is the
standard way to run an animation: it stays in step with the display, and the
browser pauses it entirely when the tab is in the background, which a timer would
not.
\end{jargon}

\begin{honest}{The counter ignores the reduced-motion preference}
This function animates for 2.2 seconds regardless of whether the visitor has asked
their operating system to reduce motion. The stylesheet honours that preference
carefully for everything it controls (lines 762 to 768), so the intent is plainly
there; it simply was not carried into this file.

The fix is small and this document states it precisely so nobody has to design it:

\begin{lstlisting}[language=Java,caption={A four-line fix. This code is a
suggestion in this document and is NOT present in the repository.}]
const reduce = window.matchMedia('(prefers-reduced-motion: reduce)').matches;
if (reduce) {
  el.textContent = (isFloat ? target.toFixed(dec)
                            : target.toLocaleString()) + suffix;
  return;
}
\end{lstlisting}

Placed after line 26, that shows the final number immediately and skips the
animation entirely, which is the correct behaviour: the information is preserved
and only the movement is removed.
\end{honest}

\begin{honest}{The metric values are hidden from anything that does not run
JavaScript}
The real numbers live in \code{data-} attributes, and the visible text is
\code{0}. Search engines, link previews, and any visitor with JavaScript disabled
therefore see a page claiming zero of everything. Writing the true value in the
HTML and having the counter animate towards it from zero would remove that risk
entirely at no cost to the effect.
\end{honest}

\paragraph{Lines 38 to 53: the observer factory.}

\begin{lstlisting}[language=Java,firstnumber=37,caption={\file{js/main.js}, lines 37--53}]

  /* ── Intersection observer factory ── */
  function observe(selector, className, threshold, callback) {
    const els = document.querySelectorAll(selector);
    if (!els.length) return;

    const io = new IntersectionObserver((entries) => {
      entries.forEach((e) => {
        if (!e.isIntersecting) return;
        e.target.classList.add(className || 'visible');
        if (callback) callback(e.target);
        io.unobserve(e.target);
      });
    }, { threshold: threshold || 0.18 });

    els.forEach((el) => io.observe(el));
  }
\end{lstlisting}

This is the best-designed thing in the repository: a small general-purpose helper
that the four lines following it use to do four different jobs.

\begin{jargon}{IntersectionObserver}
A browser facility that watches elements and tells you when they enter or leave
the visible area. Before it existed, the only way to do this was to listen for
every scroll event and measure each element's position by hand, which ran hundreds
of times a second and made pages stutter. The observer does the work outside the
main thread and reports only when something actually changes.
\end{jargon}

Line 40 finds every matching element. \code{querySelectorAll} returns a list, which
may be empty, rather than \code{null}.

Line 41 guards against an empty list and returns early. \code{!els.length} reads as
``if the count is zero'', because \code{0} is falsy and \code{!} inverts it. This
is the same defensive habit as line 10.

Lines 43 to 50 create the observer. Its callback receives a list of entries, one
for each element whose visibility changed, and for each one it:

\begin{enumerate}
\item skips anything that is leaving rather than entering
  (\code{if (!e.isIntersecting) return});
\item adds the class, defaulting to \code{visible} when the caller did not name
  one;
\item runs the optional extra callback, which only the counter uses;
\item stops watching that element with \code{io.unobserve}.
\end{enumerate}

That fourth step is the one that matters most.

\begin{keyidea}{Why it stops watching}
Every reveal in this design is meant to happen once. Without \code{unobserve}, an
element that scrolled out of view and back would fire again, and for the metrics
that would mean the numbers restarting their count every time the visitor scrolled
past. Unobserving also releases the browser's bookkeeping for that element, so
what starts as thirty watched elements steadily becomes zero.
\end{keyidea}

Line 50 sets the threshold, defaulting to 0.18: the callback fires once 18\% of
the element is visible, rather than at the first pixel. Waiting for a fraction to
appear is what makes the reveal feel connected to the scroll instead of triggering
just off screen.

\begin{honest}{The default threshold is written in a way that cannot express zero}
\code{threshold || 0.18} substitutes the default whenever the argument is falsy,
and \code{0} is falsy. So a caller asking for a threshold of \code{0}, meaning
``fire the moment a single pixel appears'', would silently get \code{0.18}
instead. The same objection applies to \code{className || 'visible'}, though an
empty class name is not a sensible request.

No caller in this file passes \code{0}, so nothing is broken today. It is a latent
trap of a kind worth recognising, and the modern remedy is the \code{??} operator,
which substitutes only for \code{null} and \code{undefined} and leaves \code{0}
alone.
\end{honest}

Line 52 starts watching every element found.

\paragraph{Lines 55 to 67: the four things being watched.}

\begin{lstlisting}[language=Java,firstnumber=54,caption={\file{js/main.js}, lines 54--67}]

  /* reveal (fade + slide) */
  observe('.reveal', 'visible', 0.15);

  /* staggered children */
  observe('.stagger', 'visible', 0.15);

  /* pipeline animation */
  observe('.pipeline', 'visible', 0.2);

  /* counters */
  observe('.metric-num[data-target]', 'counted', 0.5, (el) => {
    animateCounter(el);
  });
\end{lstlisting}

Four calls, and the whole page's reveal behaviour. The first three add the class
\code{visible} to anything carrying \code{reveal}, \code{stagger} or
\code{pipeline}, at slightly different thresholds. The stylesheet does everything
after that.

The fourth is the different one. Its selector,
\code{.metric-num[data-target]}, matches elements that carry the class \emph{and}
have the attribute, which is a small piece of defensiveness: a metric element
without a target would be skipped rather than animated to \code{NaN}. Its
threshold is 0.5, so the numbers only start counting when the band is half visible,
and its callback runs the counter.

\begin{honest}{The class \code{counted} is added and styled nowhere}
The fourth call passes \code{'counted'} rather than \code{'visible'}, so each
metric number gains a class that no CSS rule and no other JavaScript ever reads. It
is harmless, and it serves as a marker a developer could inspect, but it is worth
knowing that it does nothing. Note also that the metric numbers are children of
\code{.metrics-grid}, which is a \code{.stagger}, so they are already being faded
in by the second call: the class here is genuinely spare.
\end{honest}

\paragraph{Lines 69 to 79: smooth anchor scrolling, and the close.}

\begin{lstlisting}[language=Java,firstnumber=68,caption={\file{js/main.js}, lines 68--79}]

  /* ── Smooth anchor scroll (book-a-call button) ── */
  document.querySelectorAll('a[href^="#"]').forEach((a) => {
    a.addEventListener('click', (e) => {
      const target = document.querySelector(a.getAttribute('href'));
      if (!target) return;
      e.preventDefault();
      target.scrollIntoView({ behavior: 'smooth', block: 'start' });
    });
  });

})();
\end{lstlisting}

Line 70 finds every link whose destination starts with a hash.
\code{a[href\^{}="\#"]} is an attribute selector, and \code{\^{}=} means ``starts
with''. On this page that matches three links: the navigation's ``Book a call'',
the hero's ``Book a call'', and the hero's ``See the work''.

Line 72 looks up the destination by using the link's own \code{href} as a
selector, which works because \code{\#book} is valid CSS meaning ``the element
with id book''.

Line 73 returns early if the destination does not exist, leaving the browser to
handle the click normally. Line 74 then cancels the browser's default jump, but
only after the destination has been confirmed, which is the correct order: cancel
first and a broken link would do nothing at all instead of at least trying.

Line 75 performs the scroll. \code{block: 'start'} aligns the destination with the
top of the window.

\begin{honest}{This whole block is redundant, and it carries a latent crash}
\file{css/site.css} line 28 already sets \code{scroll-behavior: smooth} on the
document, which makes every anchor link glide without any JavaScript at all. These
eleven lines reproduce that.

They also introduce a failure the CSS version does not have. A link written as
\code{href="\#"}, which is a common way to write a placeholder link, would reach
line 72 as \code{document.querySelector('\#')}. That is not a valid selector, so
the browser throws a \code{SyntaxError}, and because the error happens inside a
click handler it would break that click and leave an error in the console.

No such link exists in the page today, so this is latent rather than live. The
honest summary is that the safest change to this file would be to delete lines 69
to 77 entirely: it would remove a redundancy and a hazard in one edit, and nothing
visible would change.
\end{honest}

Line 79 closes and immediately runs the wrapper opened on line 5.

\subsubsection*{What it connects to}

\textbf{Named by:} \file{new-design-reference/index.html}, line 595, and nothing
else.

\textbf{Names:} no file at all. It imports nothing and fetches nothing.

\textbf{Its entire interface with the rest of the project is seven strings}, and
they are worth listing because they are the contract nothing checks:

\begin{table}[H]\centering\small
\begin{tabular}{@{}lll@{}}
\toprule
String & Used at & Where it must also exist \\
\midrule
\code{.nav}        & line 9  & \file{index.html} line 23, \file{site.css} line 95 \\
\code{scrolled}    & line 12 & \file{site.css} line 105 \\
\code{.reveal}     & line 56 & \file{index.html} (five places), \file{site.css} line 78 \\
\code{.stagger}    & line 59 & \file{index.html} (four places), \file{site.css} line 86 \\
\code{.pipeline}   & line 62 & \file{index.html} line 223, \file{site.css} line 381 \\
\code{visible}     & lines 56--62 & \file{site.css} lines 83, 92, 400 \\
\code{.metric-num} & line 65 & \file{index.html} (four places), \file{site.css} line 343 \\
\bottomrule
\end{tabular}
\caption{Every point of contact between this file and the other two. Break any
one of these strings in one place and the page keeps working, minus one
behaviour, with no error reported anywhere.}
\end{table}

\subsubsection*{How it breaks}

\begin{itemize}
\item \textbf{Load it in the \code{<head>} instead of at the end of the body} and
  all three behaviours vanish silently, since every lookup runs immediately and
  finds nothing.
\item \textbf{Rename a class in the HTML} and the corresponding behaviour stops,
  again silently, because of the early-return guards.
\item \textbf{Delete this file} and the page is static, with four zeros in the
  metrics band. That last part is the only real damage.
\item \textbf{Add a link with \code{href="\#"}} and clicking it throws a
  \code{SyntaxError} from line 72.
\item \textbf{A visitor who has asked for reduced motion} still gets the counting
  animation, because this file never checks.
\item \textbf{An ancient browser without \code{IntersectionObserver}} would throw
  on line 43 and lose all four reveal behaviours, which would leave every
  \code{.reveal} element permanently invisible were it not for the CSS
  reduced-motion block. Every browser in current use supports it, so this is a
  historical note, not a live concern.
\end{itemize}

\subsubsection*{Honest findings for this file}

The counter ignores \code{prefers-reduced-motion}; the real metric values are
invisible without JavaScript; the smooth-scroll block duplicates a CSS feature and
carries a latent \code{SyntaxError}; the \code{counted} class is inert; and the
\code{||} defaults on line 47 and 50 cannot express a falsy argument. Against
those: the file is genuinely well written, with consistent guards, a sensible
shared helper, correct use of \code{unobserve} and \code{passive}, and no
dependencies at all.

% ===========================================================================
\subsection{\file{new-design-reference/netlify.toml}}
% ===========================================================================

\textbf{TOML, 11 lines.} Instructions for a hosting service called Netlify: how
to publish the folder, and four security headers to send with every response.

\begin{provenance}{new-design-reference/netlify.toml}
Hand-written, added in the first commit \code{0d77d51} and never edited. The
structure is the standard shape from Netlify's own documented configuration
format, which is what any developer writing one of these would produce; that is
convention rather than copying, and no specific source is claimed.
\textbf{Confidence: certain for hand-written, no provenance evidence found for
any particular template.}
\end{provenance}

\subsubsection*{Why it exists}

Today, for no reason at all: nothing uses it.

The honest reading of the history is that the site was first expected to be
hosted on Netlify, which is a common choice for a static page like this, and the
plan changed three commits later. Commit \code{2c66a7a} on 2026-07-28 is titled
in part ``switch deploy to VPS rsync'', which is the moment Netlify stopped being
the destination. This file was not removed at that point and has sat unused since.

\begin{jargon}{TOML}
A file format for configuration, designed to be obvious to read. A name in square
brackets starts a section; \code{name = "value"} sets a setting. Double square
brackets, as on line 5, mean ``another entry in a list of these'', so a file can
carry several header rules.
\end{jargon}

\subsubsection*{Line by line}

\begin{lstlisting}[firstnumber=1,caption={\file{netlify.toml}, lines 1--11, the complete file}]
[build]
  publish = "."
  command = ""

[[headers]]
  for = "/*"
  [headers.values]
    X-Frame-Options = "DENY"
    X-Content-Type-Options = "nosniff"
    Referrer-Policy = "strict-origin-when-cross-origin"
    Permissions-Policy = "camera=(), microphone=(), geolocation=()"
\end{lstlisting}

Lines 1 to 3 describe how to build the site. \code{publish = "."} means publish
this directory itself, and \code{command = ""} means there is no build step to run
first. Those two lines are the whole of the ``no build step'' decision, stated to
a hosting provider.

Line 5 opens a headers rule and line 6 applies it to \code{/*}, meaning every path
on the site.

\begin{jargon}{HTTP header}
An extra line of information a web server sends alongside a file, which the visitor
never sees. Some describe the file; the four below are instructions to the browser
about what it is allowed to do with the page.
\end{jargon}

Line 8, \code{X-Frame-Options = "DENY"}, forbids any other site from embedding this
page inside a frame. That defends against clickjacking, where an attacker loads your
real page invisibly over their own and tricks a visitor into clicking your buttons
while believing they are clicking something else.

Line 9, \code{X-Content-Type-Options = "nosniff"}, tells the browser to believe the
declared type of every file rather than guessing from its contents. Guessing is a
historical browser behaviour that can be manipulated into treating an uploaded image
as a script.

Line 10, \code{Referrer-Policy = "strict-origin-when-cross-origin"}, limits what is
disclosed when a visitor follows a link away from the page. Within this site the full
address is passed; to another site only the domain is sent, and to an insecure site
nothing at all. This is the modern default and a sensible privacy choice.

Line 11, \code{Permissions-Policy}, switches off three device capabilities for this
page entirely: camera, microphone and location. The empty parentheses mean ``nobody,
not even this page''. Since the page is static text and images, it needs none of
them, and turning them off means a future compromise cannot quietly request them.

\begin{keyidea}{These four headers are good practice and they are not in effect}
Every one of these is a sensible, defensive choice, and the file they live in is
read by a service this project does not use. The page is served by Caddy on the
agency VPS, which never sees this file.

Moving these four lines into the Caddy configuration on the server would take a few
minutes and would make them real. As things stand, the repository looks like it has
thought about security headers and the live site has none of them.
\end{keyidea}

\subsubsection*{What it connects to}

Nothing in this repository refers to it. Netlify appears nowhere else in the tree
and no workflow reads it. The deploy workflow does not even exclude it, so if that
workflow were ever restored, this file would be copied to the VPS where Caddy would
ignore it.

\subsubsection*{How it breaks}

It cannot break, because nothing reads it. The realistic failure is the opposite: a
future maintainer reads this file, concludes the site sends these headers, and
builds on an assumption that is false.

\begin{honest}{An unused configuration file is worse than no configuration file}
This is a small file and a real trap. It states four security guarantees that the
live site does not provide, and there is nothing in the repository to indicate it is
inert. Either delete it, or move its contents to the Caddy configuration and delete
it. Leaving it is the only option with no upside.
\end{honest}


% ===========================================================================
\subsection{\file{new-design-reference/assets/ASSETS.md}}
% ===========================================================================

\textbf{Markdown, 33 lines.} A checklist of the twelve files somebody would need
to produce before this design could be considered finished, and three notes about
placeholder content to replace.

\begin{provenance}{new-design-reference/assets/ASSETS.md}
Hand-written, added in the first commit \code{0d77d51} alongside the design it
describes, and never edited since. Written by the same person in the same sitting
as the HTML, which is evident from the fact that it references the exact class
names (\code{pt1} through \code{pt6}, \code{pa1}, \code{pa2}) used in that file.
\textbf{Confidence: high, from the commit history and the internal references.}
\end{provenance}

\subsubsection*{Why it exists}

It is a handover note. The design was built with placeholder visuals, and this
file records precisely what would replace each one and what HTML change each
replacement needs. As an artefact it is genuinely good practice: it converts
``this is not finished'' from a feeling into a list.

\begin{jargon}{Markdown}
A way of writing formatted text in a plain file, using punctuation instead of
menus: \code{\#} starts a heading, backticks mark code, and pipe characters draw a
table. It is designed to be readable as plain text and is the standard format for
notes kept alongside code.
\end{jargon}

\subsubsection*{Line by line}

\begin{lstlisting}[firstnumber=1,caption={\file{ASSETS.md}, lines 1--21}]
# Assets checklist: ai.digitalise.agency

Drop real assets here, then update the HTML as noted.

## Required before launch

| File | Size | Where used | HTML change needed |
|---|---|---|---|
| `videos/hero-1.mp4` | under 8MB | Hero background tiles | Swap .tile-blue .tile-thumb div for `<video autoplay muted loop playsinline>` |
| `logos/client-1.svg` | any | Proof section | Replace `pa1`/`pa2` avatar initials with `<img>` |
| `logos/client-2.svg` | any | Proof section | Same |
| `logos/client-3.svg` | any | Proof section | Same |
| `thumbs/port-1.jpg` | 800x450px | Portfolio card 1 | Add `style="background-image:url('/assets/thumbs/port-1.jpg')"` to `.pt1` |
| `thumbs/port-2.jpg` | 800x450px | Portfolio card 2 | Same pattern for `.pt2` through `.pt6` |
| `thumbs/port-3.jpg` | 800x450px | Portfolio card 3 | |
| `thumbs/port-4.jpg` | 800x450px | Portfolio card 4 | |
| `thumbs/port-5.jpg` | 800x450px | Portfolio card 5 | |
| `thumbs/port-6.jpg` | 800x450px | Portfolio card 6 | |
| `favicon.svg` | 32x32 | Browser tab | Add `<link rel="icon" href="/assets/favicon.svg">` to `<head>` |
| `og-image.png` | 1200x630px | Social share preview | Update `og:image` meta tag |
\end{lstlisting}

Line 1 is the title, and line 3 the instruction: put the real files in this folder,
then change the HTML as the table says.

Lines 7 and 8 open a table, with line 8 being the separator row that Markdown
requires between a table's headings and its body.

The twelve rows are the substance. Reading them tells you exactly what the design
is pretending to have:

\begin{itemize}
\item \textbf{One video}, under 8 megabytes, to play inside the hero tiles instead
  of the static icons. The suggested tag carries \code{autoplay muted loop
  playsinline}, and the \code{muted} is not optional: browsers refuse to autoplay
  video with sound, so an unmuted autoplaying video simply does not start.
  \code{playsinline} stops iPhones taking the video full screen.
\item \textbf{Three client logos}, to replace the coloured circles with initials in
  the proof section. Note the row says \code{pa1}/\code{pa2} and there are three
  logos and three avatars (\code{pa1} to \code{pa3}), so the instruction is one
  short.
\item \textbf{Six portfolio thumbnails} at 800 by 450, which is exactly the 16 by 9
  shape the stylesheet reserves with \code{aspect-ratio} on line 473.
\item \textbf{A favicon}, with the exact tag needed to install it.
\item \textbf{A social preview image} at 1200 by 630, the standard Open Graph size,
  to fix the missing \code{og:image} noted in section 7.1.
\end{itemize}

\begin{honest}{None of the twelve exist}
The \file{assets/} folder contains exactly one file: this checklist. Every asset it
calls ``required before launch'' is absent, which means the design was never
finished to the standard its own author set for it. The portfolio thumbnails are
still gradients, the client logos are still lettered circles, and the page has
neither a favicon nor a social preview image.
\end{honest}

\begin{lstlisting}[firstnumber=21,caption={\file{ASSETS.md}, lines 21--33}]

## Numbers to update (index.html metrics band)

Search for `data-target` in `index.html` and replace all four values with real numbers verified from Gantry / n8n data.

## Booking link

Search for `YOUR_CAL_LINK` in `index.html` and replace with the actual Cal.com or Calendly URL.
For a full inline calendar embed, replace the `<a>` tag with the Cal.com embed snippet from cal.com/docs/embeds.

## Client proof section

Three `<!-- UPDATE -->` comments in the proof section mark where real client names and companies go.
\end{lstlisting}

Line 24 is the most interesting line in the file, because of one phrase: the four
metric numbers are to be replaced with ``real numbers verified from Gantry / n8n
data''. Gantry is another project in this workspace, the agency's content
publishing system, and n8n is the automation tool behind several of these projects.
So the author's intention was that the numbers on this marketing page would be
traceable to the agency's own operational systems rather than estimated. That is a
higher standard than most marketing pages hold themselves to, and it makes the
invented hero tile numbers in \file{index.html} more surprising rather than less.

Line 28 tells the reader to search for \code{YOUR\_CAL\_LINK}.

\begin{honest}{Two of this file's three instructions are stale}
Searching \file{index.html} for \code{YOUR\_CAL\_LINK} finds nothing. The literal
placeholder was replaced on 2026-07-28 by commit \code{4dffc12}, which set the link
to \code{cal.com/davonex/30min}. So somebody following this instruction searches,
finds nothing, and has no way to tell whether the job is done or the instruction is
wrong. The answer is that the job is half done: the placeholder text is gone and the
link still points at another project's calendar.

The metrics instruction on line 24 is stale in the same way: commit \code{2c66a7a}
claims to have replaced all placeholder content with real agency data, and the four
values do look like considered numbers rather than round placeholders. Yet the
\code{UPDATE} comment sits in \file{index.html} at line 187 and this instruction sits
here, both saying the work is outstanding.

Only the third instruction, on line 33 about the client proof section, is verifiably
still true: the avatar initial on line 430 of \file{index.html} is a leftover.
\end{honest}

\subsubsection*{What it connects to}

It describes \file{new-design-reference/index.html} and
\file{new-design-reference/css/site.css} by name and by class, and it is the only
file in the repository that documents the project's intent. Nothing reads it: it is
for people.

The deploy workflow explicitly excludes it from publication (line 82 of the
workflow), which is the right instinct, and that exclusion no longer matches its
path. See section 7.6.

\subsubsection*{How it breaks}

As a document it cannot break. As a source of truth it already has: two of its
three instructions describe a state of the code that no longer exists, which is the
characteristic failure of documentation kept beside code and not updated with it.


% ===========================================================================
\subsection{\file{.github/workflows-disabled/deploy.yml.disabled}}
% ===========================================================================

\textbf{YAML, 103 lines.} The automation that used to publish this site: it
watched for pushes, checked for a marker in the commit message, copied the files
to the agency's server over an encrypted connection, and then checked that the
site still answered. It is currently switched off.

\begin{provenance}{.github/workflows-disabled/deploy.yml.disabled}
Hand-written, added in commit \code{d49e034} on 2026-07-28 at the path
\file{.github/workflows/deploy.yml}, rewritten in \code{2c66a7a} the same day when
the deploy target changed from Netlify to the VPS, fixed three times on 2026-07-29,
and finally renamed to its present path in \code{da508a7}.

Its shape is shared with the sibling project \code{salmanadnan-com}, whose
workflow was written in the same week and carries the same retry loops and the
same \code{[deploy]} marker check. Neither file says which came first, and this
document does not guess. The \code{[deploy]} convention itself is
workspace-wide and predates both. \textbf{Confidence: certain for hand-written and
for the rename history, which is recorded in git; the relationship to the sibling
repository is stated as a resemblance and not as a lineage.}
\end{provenance}

\subsubsection*{Why it exists}

To remove a manual step. Without it, publishing means somebody connecting to the
server and copying files by hand, which is slow and easy to get wrong. With it,
publishing is a consequence of writing a commit message in a particular way.

\begin{jargon}{YAML}
A file format for configuration, where structure is expressed by indentation
rather than by brackets. Two spaces of indent mean ``this belongs to the line
above''. It is unforgiving about indentation, which is its main drawback: a single
misplaced space changes the meaning.
\end{jargon}

\begin{jargon}{GitHub Actions, and a workflow}
A service that runs jobs on GitHub's computers when something happens in a
repository. The jobs are described in a workflow file in
\file{.github/workflows/}. Each job runs on a fresh, empty machine, which is why
the file below has to install its own SSH key every time.
\end{jargon}

\subsubsection*{Line by line}

\paragraph{Lines 1 to 19: the header comment and the trigger.}

\begin{lstlisting}[firstnumber=1,caption={\file{deploy.yml.disabled}, lines 1--19}]
# ai.digitalise.agency: deploy to VPS when commit subject starts with [deploy].
#
# A push to main always checks the marker. It DEPLOYS only when the commit
# subject line STARTS with [deploy], matching every other agency repository.
# Mentioning the marker anywhere else in the message does nothing.
#
# The site is served by Caddy from /root/ai.digitalise.agency on the VPS.
# Caddy picks up the new files immediately, no restart needed.
#
# Required repository secrets:
#   VPS_HOST     46.202.194.85
#   VPS_USER     deploy
#   VPS_SSH_KEY  private key for that user
name: Deploy ai.digitalise.agency

on:
  push:
    branches: [main]
\end{lstlisting}

Lines 1 to 13 are a comment block, and it is the best documentation in the
repository. In YAML a \code{\#} starts a comment. It states the marker rule twice
and in capitals, which tells you it had caught somebody out; it names the serving
path and the fact that Caddy needs no restart; and it lists the three secrets the
workflow needs.

\begin{honest}{The server's address is written in the comment that explains it is
a secret}
Line 11 records \code{VPS\_HOST 46.202.194.85} in plain text, immediately under the
heading ``Required repository secrets''. Whatever the value of keeping the address
in GitHub's secret store, writing it here defeats it.

Proportion matters: an IP address is not a credential, this repository is private,
and the server is reachable on the public internet anyway. The genuine cost is that
it is now in the git history permanently, so if the server is ever replaced this
line becomes stale and misleading rather than merely redundant. The private key,
which is the thing that actually matters, is correctly absent.
\end{honest}

Line 14 names the workflow, which is the label shown in GitHub's interface.

Lines 16 to 18 set the trigger: run on every push to the branch \code{main}, and
nothing else. Not on pull requests, not on a schedule, not manually. That last
absence is worth noting: there is no \code{workflow\_dispatch}, so there is no way
to publish without making a commit.

\paragraph{Lines 20 to 38: the marker check.}

\begin{lstlisting}[firstnumber=19,caption={\file{deploy.yml.disabled}, lines 19--38}]

jobs:
  check-deploy:
    name: Check Deploy Marker
    runs-on: ubuntu-latest
    outputs:
      should_deploy: ${{ steps.check.outputs.should_deploy }}
    steps:
      - id: check
        env:
          COMMIT_MSG: ${{ github.event.head_commit.message }}
        run: |
          SUBJECT="${COMMIT_MSG%%$'\n'*}"
          if [[ "$SUBJECT" == "[deploy]"* ]]; then
            echo "should_deploy=true" >> "$GITHUB_OUTPUT"
            echo "Deploy marker found on the subject line."
          else
            echo "should_deploy=false" >> "$GITHUB_OUTPUT"
            echo "No deploy marker: skipping deploy."
          fi
\end{lstlisting}

The first job exists only to answer a yes-or-no question, and it runs on a fresh
Ubuntu machine to do it. Lines 24 and 25 declare that this job produces an output
called \code{should\_deploy}, which the second job will read.

Line 28 passes the commit message into the script as an environment variable rather
than pasting it directly into the script text. That distinction is a genuine
security measure.

\begin{keyidea}{Why the commit message goes through an environment variable}
If line 31 had been written as \code{SUBJECT="\$\{\{ github.event.head\_commit.message
\}\}"}, GitHub would paste the commit message directly into the script before the
shell ever sees it. Anyone who could write a commit message could then write shell
commands into it and have them run on the machine holding the deploy key.

Passing it as an environment variable and reading it as \code{\$COMMIT\_MSG} means
the message is data, never code. This is a well-known class of vulnerability in
GitHub Actions, and this file gets it right.
\end{keyidea}

Line 31 extracts the subject line: \code{\$\{COMMIT\_MSG\%\%\$'\textbackslash n'*\}}
is shell syntax meaning ``delete the longest match of a newline followed by anything,
from the end'', which leaves everything before the first newline.

Line 32 is the test itself. \code{"\$SUBJECT" == "[deploy]"*} compares with a
trailing wildcard, which is why the marker must \emph{begin} the subject.

Lines 33 and 36 write the answer to \code{\$GITHUB\_OUTPUT}, a file GitHub reads to
collect a step's outputs. Lines 34 and 37 print a human-readable explanation into
the log, which is a small kindness to whoever is wondering why nothing deployed.

\paragraph{Lines 40 to 63: starting the deploy, and the SSH setup.}

\begin{lstlisting}[firstnumber=39,caption={\file{deploy.yml.disabled}, lines 39--63}]

  deploy:
    name: Deploy to VPS
    runs-on: ubuntu-latest
    needs: check-deploy
    if: needs.check-deploy.outputs.should_deploy == 'true'
    steps:
      - uses: actions/checkout@v4

      - name: Authorize SSH
        env:
          KEY: ${{ secrets.VPS_SSH_KEY }}
          HOST: ${{ secrets.VPS_HOST }}
        run: |
          mkdir -p ~/.ssh && chmod 700 ~/.ssh
          printf '%s\n' "$KEY" > ~/.ssh/deploy_key
          chmod 600 ~/.ssh/deploy_key
          for i in 1 2 3 4 5; do
            if ssh-keyscan -T 10 -H "$HOST" >> ~/.ssh/known_hosts 2>/tmp/keyscan.err; then
              [ -s ~/.ssh/known_hosts ] && break
            fi
            echo "ssh-keyscan attempt $i failed: $(cat /tmp/keyscan.err 2>/dev/null)"
            sleep 5
          done
          [ -s ~/.ssh/known_hosts ] || { echo "::error::could not fetch the host key"; exit 1; }
\end{lstlisting}

Line 43 makes this job wait for the first, and line 44 runs it only if the answer
was yes. Everything below is therefore skipped entirely on an ordinary commit.

Line 46 checks out the repository onto the machine, using GitHub's own action.
Without this the machine would be empty.

Lines 53 to 55 install the private key: create the \file{.ssh} folder,
write the key into it, and set permissions so only the owner can read it.
\code{chmod 600} is not optional; SSH refuses to use a key file that other users
can read, and its error message when it does is famously unhelpful.

\begin{jargon}{SSH key pair, and the host key}
Connecting securely involves two separate keys and confusing them is common. The
\emph{deploy key} here is the private half of a pair proving \emph{we} are allowed
in; its public half sits on the server. The \emph{host key} is the server's own
identity, which the client checks to be sure it is talking to the right machine and
not an impostor. Lines 53 to 55 install the first. Lines 56 to 63 fetch the second.
\end{jargon}

Lines 56 to 62 are the retry loop, and they exist because this step failed
repeatedly. \code{ssh-keyscan} asks the server for its host key, with a 10-second
timeout, appending the result to \file{known\_hosts} and any complaint to a
temporary file. \code{[ -s ~/.ssh/known\_hosts ]} then tests that the file is
non-empty, because \code{ssh-keyscan} can succeed as a command while producing
nothing at all. That distinction is the subtle part: checking the exit code alone
would not have been enough.

Line 63 fails the whole job with a message if five attempts produced nothing.
\code{::error::} is GitHub's syntax for marking a line in the log as an error so it
surfaces in the interface.

\begin{keyidea}{Why fetching the host key each time is a real trade-off}
Fetching a server's identity at the moment you connect, and trusting whatever comes
back, means you are not actually verifying anything: an attacker positioned between
the runner and the server could supply their own. The secure alternative is to store
the expected host key as a secret and compare.

That is a real weakness and it is a widespread convention, and the sibling project
\code{passbolt-digitalise} in this workspace does it the stricter way, keeping a
\file{.github/known\_hosts} file in the repository. Two projects in one workspace
solving the same problem two ways is worth knowing about.
\end{keyidea}

\paragraph{Lines 65 to 83: copying the files.}

\begin{lstlisting}[firstnumber=64,caption={\file{deploy.yml.disabled}, lines 64--83}]

      - name: Sync files
        env:
          HOST: ${{ secrets.VPS_HOST }}
          USER: ${{ secrets.VPS_USER }}
        run: |
          retry() {
            for attempt in 1 2 3; do
              "$@" && return 0
              echo "attempt $attempt failed: $*"
              [ "$attempt" -lt 3 ] && sleep $(( attempt * 15 ))
            done
            return 1
          }
          retry rsync -az --no-owner --no-group --omit-dir-times --delete \
            -e "ssh -i ~/.ssh/deploy_key -o BatchMode=yes -o ConnectTimeout=25" \
            --exclude '.git' \
            --exclude '.github' \
            --exclude 'assets/ASSETS.md' \
            ./ "$USER@$HOST:/root/ai.digitalise.agency/"
\end{lstlisting}

Lines 70 to 77 define a small shell function that runs whatever it is given, up to
three times, waiting longer between each attempt: 15 seconds, then 30. \code{"\$@"}
means ``all the arguments, run as a command''. This is the second retry loop in the
file, and like the first it is evidence of a step that failed in practice.

Line 78 is the actual copy. Each flag earns its place:

\begin{table}[H]\centering\small
\begin{tabular}{@{}ll@{}}
\toprule
Flag & What it does \\
\midrule
\code{-a} & archive mode: recurse, and preserve timestamps, permissions and links \\
\code{-z} & compress during transfer \\
\code{--no-owner} & do not try to set file ownership on the far end \\
\code{--no-group} & do not try to set group ownership either \\
\code{--omit-dir-times} & do not try to set directory timestamps \\
\code{--delete} & remove files on the server that no longer exist here \\
\bottomrule
\end{tabular}
\caption{The \code{rsync} flags, and why each is present.}
\end{table}

The three \code{--no}/\code{--omit} flags are all working around the same thing: the
\code{deploy} user is not the owner of the target directory and lacks the privilege
to change ownership or directory timestamps. Without them \code{rsync} reports
errors and fails. These flags are the surviving trace of the problem that commit
\code{f698421} was fixing.

\code{--delete} deserves respect: it makes the server an exact mirror, so removing a
file here removes it there. That is usually what you want and it is also the flag
that turns a mistaken source path into a wiped site.

Line 79 tells \code{rsync} how to connect: use this key, never prompt for a password
(\code{BatchMode=yes}, essential because there is no human to answer), and give up
after 25 seconds.

Lines 80 to 82 exclude three things from the copy: the git history, the workflow
folder, and the internal checklist.

Line 83 is the source and destination. \code{./} means the current directory's
\emph{contents}, and the trailing slash is significant: without it, \code{rsync}
would copy the directory itself into the target, producing a nested folder.

\begin{honest}{This line would now break the live site}
\code{./} was the right source when \file{index.html} sat at the repository root.
Since commit \code{da508a7} moved everything down one level, \code{./} now contains
a single folder called \file{new-design-reference}. Restoring this workflow
unchanged would therefore copy that folder to the server and, because
\code{--delete} is set, remove the real site's files in the same operation.

The correct source after the move would be \code{new-design-reference/}. The same
commit that moved the files disabled this workflow, so the two are consistent; the
trap only exists for somebody who re-enables one without reading the other.
\end{honest}

\begin{honest}{The exclusion for the checklist no longer matches}
\code{--exclude 'assets/ASSETS.md'} was written when that path was correct. The file
now lives at \file{new-design-reference/assets/ASSETS.md}, which this pattern does
not match, so a restored workflow would publish the internal checklist, including its
line about replacing placeholder numbers, to the public internet.
\end{honest}

\paragraph{Lines 85 to 103: verifying, and cleaning up.}

\begin{lstlisting}[firstnumber=84,caption={\file{deploy.yml.disabled}, lines 84--103}]

      - name: Verify
        env:
          HOST: ${{ secrets.VPS_HOST }}
          USER: ${{ secrets.VPS_USER }}
        run: |
          for i in $(seq 1 12); do
            code=$(ssh -i ~/.ssh/deploy_key -o BatchMode=yes "$USER@$HOST" \
              'curl -sf -o /dev/null -w "%{http_code}" --max-time 10 https://ai.digitalise.agency/' || echo 000)
            case "$code" in
              2*|3*) echo "Site answering with $code after ${i} check(s)."; exit 0 ;;
            esac
            echo "waiting, got $code ($i/12)"
            sleep 5
          done
          echo "::warning::site did not respond 2xx/3xx; check Caddy on the VPS"

      - name: Clean up key
        if: always()
        run: rm -f ~/.ssh/deploy_key
\end{lstlisting}

The verification step asks, up to twelve times over about a minute, whether the site
is answering. It does this by connecting to the server and running \code{curl}
\emph{there} rather than from the runner, which tests the server's own view of
itself.

The \code{curl} flags: \code{-s} quiet, \code{-f} treat an error status as a failure,
\code{-o /dev/null} throw the page away, \code{-w "\%\{http\_code\}"} print just the
status number, and a 10-second limit. The \code{|| echo 000} supplies a fallback so
that a failed connection yields the harmless string \code{000} instead of an empty
one.

Line 93 accepts any status beginning with 2 or 3, meaning success or redirect, and
exits successfully.

\begin{honest}{A failed verification is only a warning}
Line 99 uses \code{::warning::} rather than \code{::error::}, and the step does not
exit with a failure code. So if the site never came back, the workflow would still
finish green with a warning buried in the log. Given that the preceding step ran
\code{rsync --delete} against the live site, a failure here is the one moment you
would want a loud red failure and an email. This is the clearest correctness
weakness in the workflow, and it is one word to fix.
\end{honest}

Lines 101 to 103 delete the private key from the runner. \code{if: always()} makes
this run even when an earlier step failed, which is the point: a key must not be left
behind on a machine because something else went wrong. In practice GitHub destroys
the machine anyway, so this is belt and braces, and it is the right instinct.

\subsubsection*{What it connects to}

\textbf{Reads:} three GitHub secrets (\code{VPS\_SSH\_KEY}, \code{VPS\_HOST},
\code{VPS\_USER}) and the commit message of the push that triggered it.

\textbf{Writes:} the contents of \file{/root/ai.digitalise.agency/} on the VPS at
\code{46.202.194.85}, which Caddy serves at \code{https://ai.digitalise.agency/}.

\textbf{Excludes:} \file{.git}, \file{.github} and \file{assets/ASSETS.md}.

\textbf{Runs when:} never, currently. It sits in
\file{.github/workflows-disabled/} with the extension \file{.disabled}, and GitHub
only reads \file{.github/workflows/*.yml}.

\subsubsection*{How it breaks}

\begin{itemize}
\item \textbf{Restored unchanged}, it would copy the wrong directory and, with
  \code{--delete}, take down whatever is currently served.
\item \textbf{A missing or wrong \code{VPS\_SSH\_KEY}} fails at the \code{rsync}
  step after three attempts. This is what commit \code{ed5c154} fixed.
\item \textbf{A \code{deploy} user without permission to traverse into
  \code{/root}} fails the same way. This is what \code{f698421} fixed, and it is the
  least obvious of the three: the user can hold the key, connect successfully, and
  still be unable to reach the target directory.
\item \textbf{GitHub's runner being refused by the server} fails the host-key fetch,
  which is what both retry loops exist for.
\item \textbf{A commit whose subject does not start with \code{[deploy]}} does
  nothing at all, by design, and the log says so.
\item \textbf{The site failing to come back} produces a warning and a green tick,
  which is the failure mode worth fixing first.
\end{itemize}

\subsubsection*{Honest findings for this file}

The server address is recorded in a comment describing it as a secret; the
\code{rsync} source path and the \file{ASSETS.md} exclusion are both stale after
the move; the host key is trusted on first sight rather than pinned, unlike the
sibling \code{passbolt-digitalise} repository; and a failed post-deploy check
produces a warning rather than a failure. Against those: the commit message is
handled safely as data rather than as code, both retry loops are well built, the
key is cleaned up unconditionally, and the comment block at the top is the most
useful documentation in the repository.


% ===========================================================================
\section{Generated and third-party files}
% ===========================================================================

\textbf{There are none.} Every one of the six tracked files was written by hand
by a person. There is no lock file, no minified bundle, no vendored dependency
tree, no binary asset, and no file carrying a generator's banner. This was
checked mechanically rather than assumed:

\begin{lstlisting}[language=bash,caption={The check that was run, and its result}]
$ git grep -lIE '(DO NOT EDIT|auto-?generated|@generated|Generated by)'
(no output)
\end{lstlisting}

So the tier 2 count in Appendix A is zero, and every file is walked line by line.

% ===========================================================================
\section{Configuration and secrets}
% ===========================================================================

Nothing in this repository is configurable at run time. There is no settings
file, no environment file, and no \file{.env.example}. The page is entirely
fixed: to change what it says, you edit the HTML.

Three values are effectively configuration that has been written directly into
the files, which is worth knowing because it is where you would look to change
them:

\begin{table}[H]\centering\small
\begin{tabular}{@{}lll@{}}
\toprule
Value & Where it is & What it controls \\
\midrule
\code{https://cal.com/davonex/30min} & \file{index.html}, line 565 & the booking calendar that opens \\
\code{hello@digitalise.agency} & \file{index.html}, line 590 & the contact address in the footer \\
The colour tokens & \file{css/site.css}, lines 5--25 & the whole colour scheme \\
\bottomrule
\end{tabular}
\caption{Values that behave like settings but are written into the source.}
\end{table}

\subsection{Secrets}

The repository contains \textbf{no passwords, keys or tokens}, and it should not:
a page that runs entirely in a visitor's browser can keep no secret, because the
visitor can read every byte of it.

Three secrets are \emph{referred to} by the deploy workflow, and stored in GitHub
rather than here: \code{VPS\_HOST}, \code{VPS\_USER} and \code{VPS\_SSH\_KEY}.
Section 12 covers how they are used.

\begin{honest}{The server's IP address is written in a comment}
Line 11 of the deploy workflow records \code{VPS\_HOST 46.202.194.85} in a
comment, next to the note that it is supposed to be a repository secret. Storing
the address in a comment defeats the point of storing the value in a secret. This
is a small exposure rather than a large one, because an IP address is not a
credential and this repository is private, but it is inconsistent with its own
stated design, and the same address is reachable from anywhere on the internet.
\end{honest}

% ===========================================================================
\section{Dependencies}
% ===========================================================================

The project has \textbf{no build-time dependencies at all}: no
\file{package.json}, no \file{requirements.txt}, nothing to install. That is
unusual enough to be worth stating plainly rather than leaving as an absence.

It has exactly two run-time dependencies, both of them services on other
people's computers, both reached by the visitor's browser rather than by the
project:

\begin{table}[H]\centering\small
\begin{tabular}{@{}lll@{}}
\toprule
Dependency & Reached from & What it is for \\
\midrule
Google Fonts & \file{index.html}, lines 13--15 & downloads the Space Grotesk and Inter typefaces \\
cal.com & \file{index.html}, line 565 & shows the calendar when ``Book a call'' is pressed \\
\bottomrule
\end{tabular}
\caption{The two external services the page depends on.}
\end{table}

\begin{honest}{Fonts are loaded from Google, which the agency's main site
deliberately does not do}
The agency's main website (the \code{digitalise-website} project in this same
workspace) self-hosts its fonts, meaning it serves the font files from its own
server. This page instead fetches them from \code{fonts.googleapis.com} on every
visit. That has two consequences: every visitor's browser makes a request to
Google, which is a privacy and consent question in some jurisdictions, and the
page's appearance depends on a third party staying up.

The page does degrade gracefully if Google is unreachable, because both font
definitions list fallbacks (\code{ui-sans-serif}, \code{system-ui},
\code{sans-serif}), so text still appears in whatever the visitor's device
provides. The layout would shift, but nothing would disappear.
\end{honest}

% ===========================================================================
\section{How to run it}
% ===========================================================================

There is nothing to install and nothing to start. Two ways work:

\begin{lstlisting}[language=bash,caption={The simplest way: open the file directly}]
# From the repository root
xdg-open new-design-reference/index.html      # Linux
open      new-design-reference/index.html     # macOS
\end{lstlisting}

This works, with one caveat that will confuse a beginner. The page asks for its
stylesheet at \code{/css/site.css}, with a leading slash. A leading slash means
``from the very top of this site''. When the file is opened directly off a disk,
the very top is the root of your hard drive, so the browser looks for
\code{/css/site.css} there, does not find it, and renders the page with no
styling at all: black text on white, everything stacked in a column.

To see it as intended, serve the folder over HTTP instead:

\begin{lstlisting}[language=bash,caption={The way that actually looks right}]
cd new-design-reference
python3 -m http.server 8000
# then open http://localhost:8000/ in a browser
\end{lstlisting}

\begin{keyidea}{Why the leading slash matters so much}
\code{/css/site.css} means ``start at the top of the site''.
\code{css/site.css}, without the slash, would mean ``start from wherever the
current page is''. The second form would have worked both on a server and off a
disk. The first form was chosen, which is correct for the live site and
inconvenient for local inspection. This one character is the single most likely
thing to make a newcomer think the project is broken when it is not.
\end{keyidea}

\subsection{What you should see}

A full-screen dark navy hero with the headline ``AI systems built to grow while
you sleep'', eight small cards drifting in and out behind it, and below that, in
order: four numbers that count up as they scroll into view, a five-step
horizontal pipeline, six project cards, three client cards, six service tiles, a
booking panel, and a one-line footer.

% ===========================================================================
\section{The agent-facing surface}
% ===========================================================================

Several projects in this workspace are operated through Claude Code and ship
their own commands, skills and hooks. \textbf{This project ships none.} That was
checked rather than assumed:

\begin{lstlisting}[language=bash,caption={The check, and its result}]
$ find . -maxdepth 2 \( -name 'CLAUDE.md' -o -name 'skills.md' \)
(no output)
$ find .claude -type f 2>/dev/null
(no output)
\end{lstlisting}

There is no \file{.claude/} folder, no \file{CLAUDE.md}, no skill and no hook.

\begin{honest}{No CLAUDE.md and no README means the repository explains nothing
about itself}
There is no README either. A person cloning this repository is given six files
and no statement of what it is, whether it is live, or why everything sits inside
a folder called \file{new-design-reference}. The only explanation anywhere is the
message on commit \code{da508a7}, which is invisible unless you know to run
\code{git log}. For a repository whose central fact is ``this is parked, the live
site is elsewhere'', that fact belongs in a README where it cannot be missed. This
is the most consequential gap in the project and it would take five minutes to
close.
\end{honest}

The document you are reading, and its companion Brief, are produced by a
workspace-level command, \code{/explain-projects}, which lives outside this
repository and writes into \file{docs/explainer/}. That folder is the only part
of this repository that command owns.

% ===========================================================================
\section{Deployment}
% ===========================================================================

\begin{figure}[H]\centering
\begin{tikzpicture}[node distance=9mm and 14mm]
  \node[box]                   (push) {A commit is pushed\\to \code{main}};
  \node[dec, right=of push]    (chk)  {subject starts\\with \code{[deploy]}?};
  \node[box, right=of chk]     (ssh)  {Fetch the host key,\\write the SSH key};
  \node[box, below=of ssh]     (sync) {\code{rsync} the files\\to the VPS};
  \node[box, left=of sync]     (ver)  {Ask the VPS to fetch\\the live URL};
  \node[box, left=of ver]      (stop) {Stop, do nothing};
  \draw[flow] (push) -- (chk);
  \draw[flow] (chk)  -- node[lbl]{yes} (ssh);
  \draw[flow] (chk)  |- node[lbl,pos=0.3]{no} (stop);
  \draw[flow] (ssh)  -- (sync);
  \draw[flow] (sync) -- (ver);
\end{tikzpicture}
\caption{The publishing pipeline as written in the workflow file. It is drawn in
the past tense on purpose: this pipeline is switched off and has been since
2026-07-29.}
\end{figure}

The mechanism is explained line by line in section 7.6. The short version: a
push to the \code{main} branch checks whether the commit's first line begins with
the literal text \code{[deploy]}. If it does, the files are copied to the VPS with
\code{rsync} and the workflow then asks the VPS to fetch the live address to
confirm the site still answers. If it does not, nothing happens.

\begin{keyidea}{The \code{[deploy]} marker must start the subject line}
Every repository in this workspace shares this convention, and it catches people
out. Writing ``fix the header and \code{[deploy]}'' does nothing at all: the check
is \code{"\$SUBJECT" == "[deploy]"*}, which tests the \emph{beginning} of the
string. The marker must be the first thing in the commit message.
\end{keyidea}

\subsection{Why it is disabled, and what would happen if it were not}

The workflow is disabled twice over, as section 4 explained. If somebody restored
it by moving the file back and renaming it, it would not work, and it is worth
knowing why before anybody tries.

The \code{rsync} command copies \code{./}, the repository root, to
\code{/root/ai.digitalise.agency/} on the server. When the workflow was written,
\file{index.html} was at that root. Since commit \code{da508a7} it is one level
down. So a restored workflow would copy a folder called
\file{new-design-reference} to the server, the server would look for a home page
at the top and find none, and the live site would break. The same commit that
moved the files is the one that disabled the workflow, so the two changes are
consistent with each other; the trap is only there for somebody who re-enables
one without the other.

\begin{honest}{One exclusion is now stale as well}
The \code{rsync} command excludes \code{assets/ASSETS.md} so the internal
checklist is not published. After the move, that file is at
\code{new-design-reference/assets/ASSETS.md}, which the exclusion pattern no
longer matches. Re-enable the workflow unchanged and the checklist, including the
line about replacing placeholder numbers, would be published to the public
internet.
\end{honest}

% ===========================================================================
\section{Honest findings}
% ===========================================================================

Everything here is also stated at the point in the walkthrough where it arises.
It is collected once so it can be read as a list.

\subsection{Findings about correctness}

\begin{enumerate}
\item \textbf{The live site's canonical link points at a different domain.} The
  page served at ai.digitalise.agency declares salmanadnan.com as its canonical
  URL, which asks search engines to index the other site instead. Verified live on
  2026-08-11. Not fixable from this repository, but whoever fixes it should know.
\item \textbf{A restored deploy workflow would break the live site}, because the
  files moved and the \code{rsync} source path did not. Section 13.1.
\item \textbf{The \code{rsync} exclusion for \file{ASSETS.md} no longer matches}
  its path, so re-enabling the workflow would publish the internal checklist.
\end{enumerate}

\subsection{Findings about dead or contradictory code}

\begin{enumerate}\setcounter{enumi}{3}
\item \textbf{Two colour tokens are defined and never used.} \code{-{}-accent2}
  (an orange) and \code{-{}-shadow} are declared in \file{site.css} and read zero
  times. Counted mechanically: \code{grep -c 'var(-{}-accent2)'} returns 0.
\item \textbf{The rule \code{.metric-num .accent} styles nothing.} No element in
  \file{index.html} carries the class \code{accent}.
\item \textbf{One responsive rule contradicts itself.} In the 900-pixel
  breakpoint, tile 3 is shown by one rule (line 748) and hidden by another five
  lines later (line 751). The hide wins, so the show is dead.
\item \textbf{\code{.tile-lines \{ padding: 0 \}} is overridden every time it is
  used}, because both places that use the class set \code{padding:10px 0} inline
  in the HTML.
\item \textbf{The class \code{counted} is added by the JavaScript and styled
  nowhere.} It is harmless, and it does serve as a marker, but nothing reads it.
\end{enumerate}

\subsection{Findings about unfinished work}

\begin{enumerate}\setcounter{enumi}{8}
\item \textbf{Five ``UPDATE'' comments remain in \file{index.html}} telling a
  future editor to replace placeholder content, even though the commit that
  populated the page claimed to replace all placeholder content with real agency
  data. Some were acted on and the comments were not removed.
\item \textbf{\file{ASSETS.md} is stale in two specific ways.} It tells the
  reader to search \file{index.html} for \code{YOUR\_CAL\_LINK}, which no longer
  exists there, and to replace the four metric numbers, which appear to have
  been replaced already.
\item \textbf{Not one of the twelve assets \file{ASSETS.md} calls
  ``required before launch'' exists.} The \file{assets/} folder contains only the
  checklist itself. In particular there is no favicon and no social preview
  image, and \file{index.html} has neither a \code{<link rel="icon">} nor an
  \code{og:image} tag.
\item \textbf{The decorative hero tiles display invented engagement numbers}
  (847K views, 234K impressions, 1.2M, 512K, 1.8M, 19K reactions). They are marked
  \code{aria-hidden} and are plainly ornamental, but a visitor reads them as
  claims. This is worth flagging given that the agency's main site was built
  specifically to avoid fabricated numbers.
\item \textbf{\file{netlify.toml} configures a host that is not used.} The
  deploy path is \code{rsync} to a VPS running Caddy. Netlify appears nowhere
  else in the repository.
\end{enumerate}

\subsection{Findings about accessibility}

\begin{enumerate}\setcounter{enumi}{13}
\item \textbf{The counting numbers ignore the reduced-motion preference.} The
  stylesheet honours \code{prefers-reduced-motion} for every CSS animation, but
  the counter is animated in JavaScript over 2.2 seconds and never checks that
  setting. A visitor who has asked their system to reduce motion still gets
  moving numbers. This is the one accessibility defect with a clear fix, given in
  section 7.3.
\item \textbf{\code{\#1967ff} on the dark background is a marginal contrast
  choice} for the small uppercase ``kicker'' labels and the 11-pixel service tags.
  This document does not have a measured contrast figure for those exact pairings,
  so it is raised as worth checking rather than asserted as a failure.
\end{enumerate}

% ===========================================================================
\section{The project's story}
% ===========================================================================

Parts 1 to 4 and part 8 are fact, taken from the repository and its history. Parts
5, 6 and 7 are advice and judgement, and are labelled as such where they begin.

\subsection{Deliverables}

What this project actually produced:

\begin{itemize}
\item A complete, self-contained one-page marketing website: 597 lines of HTML,
  768 of CSS, 79 of JavaScript, with no dependencies to install.
\item A hosting configuration for Netlify (\file{netlify.toml}) that sets four
  security response headers.
\item A GitHub Actions workflow that publishes the site to the agency VPS over
  SSH, gated on a \code{[deploy]} marker, with retries on both the host-key fetch
  and the file copy, and a post-deploy check that the site answers.
\item An asset checklist (\file{ASSETS.md}) naming the twelve files that would be
  needed to finish the design.
\end{itemize}

What it did not produce: a live site. The address ai.digitalise.agency serves
something else, as established in section 1.2.

\subsection{Timeline and effort}

The entire history is nine commits across two calendar days.

\begin{table}[H]\centering\small
\begin{tabular}{@{}llp{7.2cm}@{}}
\toprule
Date & Commit & Subject \\
\midrule
2026-07-28 & \code{0d77d51} & Launch ai.digitalise.agency, a visual-first marketing site \\
2026-07-28 & \code{4dffc12} & Update booking link to cal.com/davonex/30min \\
2026-07-28 & \code{d49e034} & Add deploy workflow \\
2026-07-28 & \code{2c66a7a} & \code{[deploy]} Replace all placeholder content with real agency data and switch deploy to VPS rsync \\
2026-07-29 & \code{f6d810f} & \code{[deploy]} Trigger deploy after VPS secrets added \\
2026-07-29 & \code{ed5c154} & \code{[deploy]} Fix VPS deploy key \\
2026-07-29 & \code{f698421} & \code{[deploy]} Fix /root traverse permission for deploy user \\
2026-07-29 & \code{0c3cad1} & \code{[deploy]} Remove em-dashes from titles, portfolio copy, and comments \\
2026-07-29 & \code{da508a7} & Move new design into new-design-reference and disable auto-deploy \\
\bottomrule
\end{tabular}
\caption{The complete commit history. Subjects are reproduced as written except
that commit \code{0d77d51} uses a long dash where this table uses a comma; that
character is not printable in this document's typeface.}
\end{table}

Two days, start to finish, and the second day is entirely deployment: four of the
five commits on 2026-07-29 are fixing the publishing pipeline, and the fifth
switches it off. Nothing in the history suggests the design itself was difficult.
Everything in the history suggests getting it onto a server was.

\textbf{This is calendar time, not effort.} Git records when commits were made and
nothing about hours spent. No document in this repository records effort, so none
is claimed here.

\subsection{Challenges}

Each of these is claimed because the repository carries a scar for it, not
because it seems plausible.

\textbf{1. Getting SSH access to the VPS working (three commits).} The scars are
\code{f6d810f} (``Trigger deploy after VPS secrets added''), \code{ed5c154}
(``Fix VPS deploy key'') and \code{f698421} (``Fix /root traverse permission for
deploy user''). Read in order they tell a complete story: the workflow was written
before the secrets existed, then the key was wrong, then the key was right but the
\code{deploy} user could not traverse into \code{/root} to reach the target
folder. The last of those is a genuinely unobvious Unix permission problem, and
it is the reason the deploy target is an unusual path.

\textbf{2. Fragile host-key fetching on GitHub's runners.} The workflow wraps
\code{ssh-keyscan} in a five-attempt retry loop with a five-second wait, and
wraps \code{rsync} in a three-attempt loop with a growing wait. Nobody writes two
independent retry loops into a fifty-line script unless both steps have failed on
them. The sibling repository \code{salmanadnan-com}, built the same week from the
same skeleton, carries a commit literally titled ``Retry, runner IP lottery'',
which confirms the diagnosis: GitHub's runners come from a wide pool of addresses
and some of them were being refused.

\textbf{3. Em-dashes.} Commit \code{0c3cad1} exists solely to remove them from
titles, portfolio copy and comments. This is the workspace-wide house rule showing
up as real work; the rule is enforced in this document too.

\textbf{4. The design was overtaken.} The final commit is not a fix, it is a
retreat: the page was finished and working, and then something else became the
live site, so this one was parked. That is the largest challenge recorded here and
it was not technical.

\subsection{Outcomes}

\begin{itemize}
\item \textbf{Built and complete.} The page is finished to the standard of a
  design reference: it renders, it is responsive down to a phone, and it honours
  the reduced-motion preference in CSS.
\item \textbf{Not live, and deliberately so.} The design is parked under
  \file{new-design-reference/} and the publishing workflow is disabled twice over.
\item \textbf{Superseded.} ai.digitalise.agency answers with a different page
  built elsewhere, as verified in section 1.2.
\item \textbf{Never finished as a product.} None of the twelve assets in the
  launch checklist were ever added, so the portfolio thumbnails are still colour
  gradients and the client logos are still coloured circles with a letter in them.
\item \textbf{Dormant.} No commit since 2026-07-29.
\end{itemize}

\subsection{What to do (judgement)}

\emph{From here on this section is advice, not a record of project decisions.}

\begin{itemize}
\item \textbf{Read \code{git log} before you touch anything}, because the only
  statement of this repository's status lives in a commit message.
\item \textbf{Serve the folder over HTTP to look at it}, never open the file
  directly, or the absolute stylesheet path will make it look broken.
\item \textbf{Keep the \code{[deploy]} convention} if the workflow is ever
  restored. It is shared across every repository in this workspace and its
  consistency is the point.
\item \textbf{Keep the retry loops.} They were paid for.
\item \textbf{Treat this as a design source.} Its real remaining value is the
  stylesheet: a coherent, self-contained dark theme in 768 readable lines with no
  tooling required.
\end{itemize}

\subsection{What not to do (judgement)}

\begin{itemize}
\item \textbf{Do not re-enable the deploy workflow without changing the
  \code{rsync} source path.} It would copy a folder rather than a page and take
  the live site down. Section 13.1.
\item \textbf{Do not copy the hero tiles into anything real without replacing
  their numbers.} They are invented, and this agency's main site was built to
  avoid exactly that.
\item \textbf{Do not trust \file{ASSETS.md}.} Two of its instructions refer to
  things that are no longer in the code.
\item \textbf{Do not assume this repository is what visitors see.} It is not, and
  that is the single easiest mistake to make here.
\item \textbf{Do not delete the repository because it is not live.} The
  stylesheet is reusable and the deploy workflow encodes three separate lessons
  about SSH on GitHub runners.
\end{itemize}

\subsection{If you were starting over (judgement)}

\emph{This subsection is entirely opinion, argued from what the code shows.}

\textbf{Keep exactly as it is:} the no-build-step decision, the token block at the
top of the stylesheet, the three-behaviour JavaScript, the reduced-motion block,
and the \code{[deploy]} marker convention. For a page of this size, every one of
those choices is right, and a modern framework would have added a hundred times
the machinery for no benefit a visitor could perceive.

\textbf{Do differently:} write a README on day one saying what the site is and
where it is served from, because the absence of that is what makes this repository
confusing. Use a relative stylesheet path so the page can be opened off a disk.
Self-host the two fonts, matching the agency's own main site and removing the
dependency on Google. Add the favicon and social preview image before the first
deploy rather than listing them on a checklist. And delete a placeholder comment
in the same commit that removes the placeholder.

\textbf{Drop:} \file{netlify.toml}, which configures a host this project does not
use, and the two unused colour tokens.

\textbf{Built today, it would look almost identical}, which is the strongest thing
that can be said for it. The only structural change worth making would be to move
the four security headers out of \file{netlify.toml} and into the Caddy
configuration on the VPS, so that they actually take effect.

\subsection{Open threads}

\begin{itemize}
\item \textbf{Twelve assets never delivered}, listed in \file{ASSETS.md}: one
  hero video, three client logos, six portfolio thumbnails, a favicon and a social
  preview image.
\item \textbf{The booking link is a placeholder.} \code{cal.com/davonex/30min} is
  the Davonex calendar, not a Digitalise AI one. The same placeholder is recorded
  as a known outstanding item on the agency's main website, so this is a shared
  loose end rather than a local mistake.
\item \textbf{Five ``UPDATE'' comments} remain in \file{index.html}, each marking
  work that may or may not have been done.
\item \textbf{The relationship to \code{salmanadnan-com} is undocumented.} The two
  repositories share an identical file list and an identical structure and were
  built in the same week. Neither says so.
\item \textbf{The live site's canonical link} points at salmanadnan.com. Somebody
  needs to fix that in whichever repository produces the live page.
\end{itemize}

% ===========================================================================
\section{Glossary}
% ===========================================================================

\begin{description}[leftmargin=0pt,style=nextline]
\item[Attribute] Extra information inside an opening tag, written as a name and a
  value, such as \code{href="\#book"}.
\item[Breakpoint] A screen width at which the layout changes. This project has
  two: 900 pixels and 600 pixels.
\item[Class] A label attached to elements so that CSS and JavaScript can find
  them again. It has no meaning to the browser by itself.
\item[Commit] One recorded change to a repository, with a message explaining it.
\item[CSS] Cascading Style Sheets: the file that controls appearance.
\item[Custom property] A named value in CSS, defined as \code{-{}-name: value} and
  read back as \code{var(-{}-name)}.
\item[Element] One thing on the page, produced by a tag.
\item[Hex colour] A colour written as six characters after a hash, giving red,
  green and blue amounts.
\item[HTML] HyperText Markup Language: the file that holds the content and its
  structure.
\item[IntersectionObserver] A browser facility that watches elements and reports
  when they scroll into view. This project uses it four times.
\item[JavaScript] The programming language that runs in the browser after the
  page loads.
\item[Media query] A block of CSS that only applies when a condition about the
  screen is true, such as being narrower than 900 pixels.
\item[Pseudo-class] A selector suffix describing a temporary state, such as
  \code{:hover} for ``while the mouse is over it''.
\item[rsync] A command-line program that copies files between computers,
  transferring only what differs.
\item[Selector] The part of a CSS rule that says which elements it applies to.
\item[SSH] Secure Shell: the standard way to connect to and run commands on
  another computer securely.
\item[Tag] The marker that creates an element, written between angle brackets.
\item[Viewport] The visible area of the page in the browser window.
\item[VPS] Virtual Private Server: a rented computer the agency controls
  remotely.
\item[Web server] A program that answers requests for files over the network.
\item[Workflow] A file describing a job GitHub runs automatically when something
  happens in the repository.
\end{description}

% ===========================================================================
\appendix
\section{Appendix A: the file ledger}
% ===========================================================================

Every path git tracks in this repository, with its tier, its line count, and the
section that explains it. This table is what makes the coverage claim in this
document checkable rather than merely asserted: run \code{git ls-files} in the
repository and compare.

\subsection*{\file{.github/}}
\begin{table}[H]\centering\small
\begin{tabular}{@{}llrl@{}}
\toprule
Path & Coverage & Lines & Section \\
\midrule
\file{.github/workflows-disabled/deploy.yml.disabled} & explained in full & 103 & 7.6 \\
\bottomrule
\end{tabular}
\caption{Coverage ledger: \file{.github/}. 1 file, 103 lines.}
\end{table}

\subsection*{\file{new-design-reference/}}
\begin{table}[H]\centering\small
\begin{tabular}{@{}llrl@{}}
\toprule
Path & Coverage & Lines & Section \\
\midrule
\file{new-design-reference/index.html}        & explained in full & 597 & 7.1 \\
\file{new-design-reference/css/site.css}      & explained in full & 768 & 7.2 \\
\file{new-design-reference/js/main.js}        & explained in full &  79 & 7.3 \\
\file{new-design-reference/netlify.toml}      & explained in full &  11 & 7.4 \\
\file{new-design-reference/assets/ASSETS.md}  & explained in full &  33 & 7.5 \\
\bottomrule
\end{tabular}
\caption{Coverage ledger: \file{new-design-reference/}. 5 files, 1{,}488 lines.}
\end{table}

\subsection*{Untracked directories (tier 3)}

\textbf{There are none.} Running \code{git status -{}-porcelain -{}-ignored} in this
repository produces no output at all: there are no ignored files, no build
output, no dependency folder, and not even a \file{.gitignore}. Every file
present on disk is tracked and every tracked file is explained above.

\subsection*{The arithmetic}

\code{git ls-files} lists \textbf{6} paths. All \textbf{6} are tier 1, explained
in full, line by line, in section 7. \textbf{0} files are tier 2 (declared and
summarized): there is no generated, vendored, minified or binary content in this
repository, checked with a search for generator banners which returned nothing.
\textbf{0} directories are tier 3 (excluded), because nothing on disk is
untracked. Total lines accounted for: \textbf{1{,}591}, which is the sum reported
by \code{git ls-files -z | xargs -0 wc -l}.

$6 + 0 = 6$, and $103 + 1{,}488 = 1{,}591$. The ledger reconciles.

\end{document}
